\chapter{时空对称与经典规范场}\label{chap:classical-fields}

普通力学中，空间平移对应动量守恒，空间转动对应角动量守恒；相对论延续的是同一思路，只是需要保持不变的几何量变成 Minkowski 间隔。由此先问两个具体问题：哪些线性坐标变换保持这个间隔？连续变换在恒等元附近由哪些无穷小生成元产生？

第一个问题定义 Lorentz 群，第二个问题给出 Lorentz Lie 代数；再加入时空平移便得到 Poincaré 群与 Poincaré 代数。电子的自旋、Dirac 场的变换律以及粒子的质量--自旋分类，都从这些生成元及其表示中产生。时空对称建立后，同样的“整体对称性 $\rightarrow$ 局域比较”思路用于内部相位：整体 $U(1)$ 给出电荷守恒，局域 $U(1)$ 则要求引入电磁规范势。Maxwell Green 张量随后把局域场方程改写成“给定源、材料与边界，场如何响应”的求解问题。

\section{Lorentz 与 Poincar\'e 对称}

在线性代数中，保持一个二次型不变的矩阵形成一个群。狭义相对论的二次型就是 Minkowski 间隔，因此保持它不变的齐次线性变换构成 Lorentz 群。惯性系的时空原点还可以任意平移，所以完整的自由时空对称性还要把四维平移加入其中；二者的组合称为 Poincar\'e 变换。

从\eqnrefs{eq:lorentz-metric} 所定义的有限变换出发，恒等变换附近的无穷小参数给出六个 Lorentz 生成元；再加入四个平移生成元，就得到 Poincar'e 代数。代数建立后，其 Casimir 算符给出质量与自旋分类；半整数自旋局域场所需的有限维表示则由同一 Lorentz 代数继续导出。

\subsection{Lorentz 生成元}

有限 Lorentz 变换描述惯性系之间的有限映射，但连续群的动力学信息集中在恒等元附近。把有限变换线性化以后得到生成元；它们的对易关系决定有限旋转与推动怎样组合，这组生成元正构成 Poincar\'e 代数的齐次部分。

第一章已经把群、Lie 群、Lie 代数、表示与不可约表示作为一般数学对象建立起来。把这些对象用于时空对称时，需要始终区分两个物理层次：一方面，Lorentz/Poincar\'e 群本身规定“允许做哪些时空变换以及这些变换怎样复合”；另一方面，一个具体的表示规定同一时空变换怎样作用在坐标四矢量、局域场分量或量子态上。后面出现的矩阵维数不同，并不意味着存在不同的 Lorentz 变换，而只是同一抽象对称作用在不同对象上的实现不同。

对粒子态还需要再强调一次“不可约”的物理意义。若 Hilbert 空间中某个子空间在全部 Poincar\'e 变换下都保持封闭，那么对称性不会把它与其它子空间混合；继续分解直到不能再分，就得到不可约表示。质量、自旋或螺旋度之所以能够成为单粒子种类的标签，正因为它们在这样的最小不变子空间内取固定值。Poincar\'e 代数的 Casimir 不变量在不可约表示上取定值，从而直接给出质量以及自旋或螺旋度标签。


连续群的关键不只是“可以取无穷小参数”，还在于两个无穷小变换的复合必须仍落在同一个群里。取 Hilbert 空间中的厄米生成元 $T_a,T_b$，分别定义
\begin{equation*}
 U_a(\epsilon)=\ee^{-\ii\epsilon T_a/\hbar},
 \qquad
 U_b(\delta)=\ee^{-\ii\delta T_b/\hbar}.
\end{equation*}
若依次沿 $a$ 方向前进、沿 $b$ 方向前进，再沿两条方向逆向返回，得到群交换子
\begin{equation}
\boxed{
 U_a(\epsilon)U_b(\delta)U_a(-\epsilon)U_b(-\delta)
 =\id-\frac{\epsilon\delta}{\hbar^2}[T_a,T_b]
 +\order(\epsilon^2\delta,\epsilon\delta^2).}
\label{eq:group-commutator-bch-dp8}
\end{equation}
这个乘积在 $\epsilon,\delta\to0$ 时仍趋近单位元，因此它的最低阶偏离必须仍属于单位元处的切空间，也就是生成元所张成的 Lie 代数。于是存在与参数无关的结构常数 $f_{ab}{}^c$，使
\begin{equation}
\boxed{
 [T_a,T_b]=\ii\hbar f_{ab}{}^cT_c.}
\label{eq:lie-algebra-closure-dp68}
\end{equation}
这条闭合关系把有限群乘法在恒等元附近的局部结构压缩成生成元的对易代数；Lorentz、Poincar\'e、$SU(2)$ 与 $SU(3)$ 的生成元关系都属于这一共同结构。BCH 展开会把这种闭合性逐阶落实到群乘法系数。

\begin{note}[两种表示中的“酉”]
物理量子态的内积必须在时空对称下保持，因此 Poincar\'e 群在物理 Hilbert 空间上由酉算符表示。局域场分量承担的任务不同：它们必须使场方程在 Lorentz 变换后保持同样的张量形式，所以只要求分量按有限维 Lorentz 表示线性混合。Lorentz 群是非紧群，其非平凡有限维表示一般不能同时保持普通欧几里得分量内积；因此有限维多分量局域场的 Lorentz 变换矩阵通常不是普通意义下的酉矩阵；二分量自旋 $1/2$ 表示中的场分量称为旋量。这与物理态空间上的时间演化和 Poincar\'e 对称必须酉实现并不矛盾。
\end{note}



四维坐标写为
\begin{equation*}
 x^\mu=(ct,\mathbf r),\qquad x_\mu=\eta_{\mu\nu}x^\nu,
\end{equation*}
并定义 Minkowski 间隔
\begin{equation*}
 s^2=x_\mu x^\mu=c^2t^2-\mathbf r^2.
\end{equation*}
齐次 Lorentz 变换是保持此二次型不变的线性变换
\begin{equation}
 x'^\mu=\Lambda^\mu{}_{\nu}x^\nu,
 \qquad
 \Lambda^{\mathsf T}\eta\Lambda=\eta.
 \label{eq:lorentz-metric}
\end{equation}
\begin{exbox}[沿 $x$ 轴的有限推动]
一般 Lorentz 变换已经由\eqnrefs{eq:lorentz-metric} 定义。现在取一个最简单但仍然完全有限的特例：惯性系 $S'$ 相对 $S$ 沿 $+x$ 方向以速度 $v$ 匀速运动。定义
\begin{equation*}
 \beta_{\rm rel}\equiv\frac{v}{c},
 \qquad
 \gamma_{\rm L}(\beta_{\rm rel})\equiv\frac1{\sqrt{1-\beta_{\rm rel}^2}},
 \qquad |\beta_{\rm rel}|<1.
\end{equation*}
相应矢量表示矩阵为
\begin{equation}
 \Lambda_x(\beta_{\rm rel})=
 \begin{pmatrix}
 \gamma_{\rm L} &-\gamma_{\rm L}\beta_{\rm rel}&0&0\\
 -\gamma_{\rm L}\beta_{\rm rel}&\gamma_{\rm L}&0&0\\
 0&0&1&0\\
 0&0&0&1
 \end{pmatrix},
 \label{eq:x-boost-matrix-dp61}
\end{equation}
即
\begin{equation*}
 ct'=\gamma_{\rm L}(ct-\beta_{\rm rel}x),
 \qquad
 x'=\gamma_{\rm L}(x-\beta_{\rm rel}ct),
 \qquad y'=y,\ z'=z.
\end{equation*}
直接代入间隔：
\begin{align*}
 c^2t'^2-x'^2
 &=\gamma_{\rm L}^2\!
 \left[(ct-\beta_{\rm rel}x)^2
 -(x-\beta_{\rm rel}ct)^2\right],\\
 &=\gamma_{\rm L}^2(1-\beta_{\rm rel}^2)(c^2t^2-x^2),\\
 &=c^2t^2-x^2.
\end{align*}
再加上不变的 $y,z$ 分量，就有 $x'_\mu x'^\mu=x_\mu x^\mu$。因此\eqnrefs{eq:x-boost-matrix-dp61} 确实是一般保度规条件的一个具体元素，而不是另一套独立公设。

有限推动由无穷小生成元连续累积。用快度 $\zeta$ 参数化这种累积：
\begin{equation}
 \beta_{\rm rel}=\tanh\zeta,
 \qquad
 \gamma_{\rm L}=\cosh\zeta,
 \qquad
 \gamma_{\rm L}\beta_{\rm rel}=\sinh\zeta.
 \label{eq:rapidity-definition-dp61}
\end{equation}
于是 $(ct,x)$ 子块变成
\begin{equation*}
 \begin{pmatrix}\cosh\zeta&-\sinh\zeta\\-\sinh\zeta&\cosh\zeta\end{pmatrix}.
\end{equation*}
连续两次同方向推动只需相乘这两个 $2\times2$ 矩阵；双曲函数加法公式给
\begin{equation*}
 \Lambda_x(\zeta_2)\Lambda_x(\zeta_1)
 =\Lambda_x(\zeta_1+\zeta_2).
\end{equation*}
因此快度线性相加，而速度满足
\begin{equation}
 \boxed{
 \beta_{12}
 =\tanh(\zeta_1+\zeta_2)
 =\frac{\beta_1+\beta_2}{1+\beta_1\beta_2}.}
 \label{eq:velocity-addition-dp61}
\end{equation}
这说明“生成元”最直观的含义就是有限连续变换在恒等元处的切向方向；快度之所以比速度更适合做 Lorentz 群参数，正因为同轴群乘法在 $\zeta$ 上变成普通加法。
\end{exbox}


考虑无穷小变换
\begin{equation*}
 \Lambda^\mu{}_{\nu}=\delta^\mu{}_{\nu}+\omega^\mu{}_{\nu}+\order(\omega^2).
\end{equation*}
把它代入 \eqnrefs{eq:lorentz-metric}，保留 $\omega$ 的一次项：
\begin{align*}
 (\id()+\omega)^{\mathsf T}\eta(\id()+\omega)
 &=\eta+\omega^{\mathsf T}\eta+\eta\omega+\order(\omega^2),\\
 &=\eta.
\end{align*}
因此
\begin{equation*}
 \omega^{\mathsf T}\eta+\eta\omega=0.
\end{equation*}
把一个指标降下，$\omega_{\mu\nu}\equiv\eta_{\mu\rho}\omega^\rho{}_{\nu}$，得到
\begin{equation}
 \omega_{\mu\nu}=-\omega_{\nu\mu}.
 \label{eq:omega-antisym}
\end{equation}
四维反对称张量只有 $4\times3/2=6$ 个独立分量：三个 $\omega_{ij}$ 生成空间转动，三个 $\omega_{0i}$ 生成推动。这六个自由度不是经验上“选出的六个变换”，而是度规不变条件 \eqnrefs{eq:omega-antisym} 的直接结果。

作为矩阵 Lie 群，$SO^+(1,3)$ 在单位元处的切空间本身就是 Lorentz 李代数
\begin{equation}
\boxed{
\mathfrak{so}(1,3)
=T_{\id()}SO^+(1,3)
=\{X\in M_4(\mathbb R)\mid X^{\mathsf T}\eta+\eta X=0\}.}
\label{eq:lorentz-lie-algebra-tangent}
\end{equation}
选取六个实基矩阵 $\ell^{\mu\nu}=-\ell^{\nu\mu}$，定义
\begin{equation}
\boxed{
(\ell^{\mu\nu})^\rho{}_{\sigma}
=\eta^{\mu\rho}\delta^\nu{}_{\sigma}
-\eta^{\nu\rho}\delta^\mu{}_{\sigma}.}
\label{eq:lorentz-matrix-basis}
\end{equation}
则任意切向量可写成 $X=\omega_{\mu\nu}\ell^{\mu\nu}/2$。由前文矩阵 Lie 群的一般结论~\eqnrefs{eq:matrix-lie-bracket-math-dp76}，$\mathfrak{so}(1,3)$ 的抽象 Lie 括号在这个定义表示中就是矩阵对易子。因而 $[\ell^{\mu\nu},\ell^{\rho\sigma}]$ 的系数由保持 Minkowski 度规的矩阵群结构唯一确定。

量子态还带来一个有限四矢量矩阵看不见的可能性：半整数自旋态做 $2\pi$ 空间转动后可以多出负号，而所有四矢量已经回到原值。因此，作用在这类量子态上的连续群必须允许“两个不同群元素对应同一个 Lorentz 四矢量变换”。这种保持群乘法的二对一映射称为 Lorentz 群的\emph{双覆盖}；这一双覆盖可由 $2\times2$ 复矩阵群 $SL(2,\mathbb C)$ 显式实现。二对一只改变全局群结构，单位元附近的 Lie 代数仍是同一个 $\mathfrak{so}(1,3)$。记 Lorentz 的厄米生成元为 $M^{\mu\nu}=-M^{\nu\mu}$，则 Hilbert 空间中的无穷小变换写为
\begin{equation}
 U(\Lambda)
 =\id()-\frac{\ii}{2\hbar}\omega_{\mu\nu}M^{\mu\nu}+\order(\omega^2).
 \label{eq:lorentz-unitary}
\end{equation}
\subsection{Poincar\'e 代数}

Lorentz 变换只改变惯性系之间的线性坐标关系，并保持时空原点不动；自由相对论理论还必须允许任意时空平移。因此完整的时空对称变换由 Lorentz 变换 $\Lambda$ 与四维平移 $a^\mu$ 共同给出：
\begin{equation}
 x'^\mu=\Lambda^\mu{}_{\nu}x^\nu+a^\mu,
 \qquad [a^\mu]=\mathrm m.
 \label{eq:poincare-affine-action-r15}
\end{equation}
连续施加两次这类变换时，第一次平移产生的四矢量还会被第二次 Lorentz 变换旋转或推动，所以 Lorentz 部分与平移部分不能彼此独立相乘。复合律为
\begin{equation}
 \boxed{
 (\Lambda_2,a_2)(\Lambda_1,a_1)
 =(\Lambda_2\Lambda_1,\,a_2+\Lambda_2a_1).}
 \label{eq:poincare-group-law-r15}
\end{equation}
这个复合律有一个值得单独命名的结构：平移本身形成一个交换子群，但 Lorentz 变换会把一个平移四矢量变成另一个平移四矢量，因此两个子群不是彼此完全独立的直积。群论中把“一个子群保持封闭，另一个子群通过共轭作用在它上面”的这种组合称为\emph{半直积}。Poincar\'e 群因此写成
\begin{equation*}
 ISO^+(1,3)=SO^+(1,3)\ltimes\mathbb R^{1,3}.
\end{equation*}
这里符号 $\ltimes$ 正是用来提醒：Lorentz 群会作用在平移四矢量空间上，而平移本身仍构成交换群。

半直积的 Lie 括号可以从 Poincar\'e 群的仿射矩阵表示直接读出
\begin{equation}
 \boxed{
 g(\Lambda,a)=
 \begin{pmatrix}
 \Lambda & a\\
 0&1
 \end{pmatrix}.}
 \label{eq:poincare-affine-matrix-dp75}
\end{equation}
单位元处的 Lie 代数元素因而是一对 $(X,u)\in\mathfrak{so}(1,3)\oplus\mathbb R^{1,3}$，
\begin{equation*}
 \mathsf x(X,u)=
 \begin{pmatrix}
 X&u\\
 0&0
 \end{pmatrix}.
\end{equation*}
使用矩阵 Lie 括号直接相减两个乘积，得到半直积的母公式
\begin{equation}
 \boxed{
 [(X,u),(Y,v)]
 =\bigl([X,Y],\,Xv-Yu\bigr).}
 \label{eq:poincare-semidir-lie-bracket-dp75}
\end{equation}
因此 Poincar\'e 代数的三类括号并不是三个彼此独立的经验事实：$[P,P]$ 来自平移理想本身的零括号，$[M,P]$ 来自 $Xv$ 这一 Lorentz 对四矢量的作用，而 $[M,M]$ 就是 $\mathfrak{so}(1,3)$ 的 Lie 括号。

连续对称在量子 Hilbert 空间中由幺正算符表示。四维平移的厄米生成元记为 $P^\mu$，其 SI 量纲为动量，
\begin{equation}
 U(a)=\exp\!\left(-\frac{\ii}{\hbar}a_\mu P^\mu\right),
 \qquad P^0=\frac{H}{c}.
 \label{eq:translation-unitary-r15}
\end{equation}
在单粒子四动量本征态上
\begin{equation}
 P^\mu|p,r\rangle=p^\mu|p,r\rangle,
 \qquad
 p^\mu=\left(\frac{E}{c},\mathbf p\right),
 \label{eq:P-eigenvalue-r15}
\end{equation}
其中 $r$ 表示除四动量外的内部量子数。Lorentz 变换的厄米生成元记为 $M^{\mu\nu}=-M^{\nu\mu}$，量纲为角动量。令 $m^{\mu\nu}=(\ell^{\mu\nu},0)$、$\tau^\rho=(0,e^\rho)$ 表示抽象 Poincar\'e Lie 代数中的 Lorentz 与平移基。\eqnrefs{eq:poincare-semidir-lie-bracket-dp75} 给出它们的结构常数；再由表示微分同态
\begin{equation*}
 dU(m^{\mu\nu})=-\frac{\ii}{\hbar}M^{\mu\nu},
 \qquad
 dU(\tau^\rho)=-\frac{\ii}{\hbar}P^\rho
\end{equation*}
送到量子 Hilbert 空间，得到三组对易关系：
\begin{align}
 [P^\mu,P^\nu]&=0,
 \label{eq:poincare-pp}\\
 [M^{\mu\nu},P^\rho]
 &=\ii\hbar\left(
 \eta^{\nu\rho}P^\mu-\eta^{\mu\rho}P^\nu
 \right),
 \label{eq:poincare-mp}\\
 [M^{\mu\nu},M^{\rho\sigma}]
 &=\ii\hbar\left(
 \eta^{\nu\rho}M^{\mu\sigma}
 -\eta^{\mu\rho}M^{\nu\sigma}
 -\eta^{\nu\sigma}M^{\mu\rho}
 +\eta^{\mu\sigma}M^{\nu\rho}
 \right).
 \label{eq:poincare-mm}
\end{align}
第一式表达平移群的交换性；第二式说明 $P^\mu$ 在 Lorentz 变换下确实按四矢量混合；第三式则是保持 Minkowski 度规的六个 Lorentz 生成元之间的闭合关系。对带内部自由度的场，总 Lorentz 生成元可分成
\begin{equation}
 M^{\mu\nu}=L^{\mu\nu}+S^{\mu\nu},
 \qquad
 L^{\mu\nu}=x^\mu P^\nu-x^\nu P^\mu,
 \label{eq:orbital-plus-spin}
\end{equation}
其中 $L^{\mu\nu}$ 作用在时空坐标依赖上，$S^{\mu\nu}$ 只混合场的内部分量；二者必须实现相同的 Lorentz 李代数结构常数。

为了与普通三维角动量语言连接，定义三个转动生成元和三个推动生成元
\begin{equation}
 J_i\equiv\frac12\epsilon_{ijk}M^{jk},
 \qquad
 K_i\equiv M^{0i},
 \qquad i=1,2,3.
 \label{eq:rotation-boost-generators-dp34}
\end{equation}
协变 Lorentz 代数随即等价于
\begin{align}
 [J_i,J_j]&=\ii\hbar\epsilon_{ijk}J_k,\\
 [J_i,K_j]&=\ii\hbar\epsilon_{ijk}K_k,\\
 [K_i,K_j]&=-\ii\hbar\epsilon_{ijk}J_k.
 \label{eq:JK-algebra}
\end{align}
前两式分别表示转动生成元自身闭合、推动在转动下按三矢量变换；第三式中的负号来自 Minkowski 度规，并说明两个不同方向的有限推动一般不再是纯推动，而会同时产生转动。这一非交换性是 Thomas--Wigner 转动的群论根源。

\begin{derivation}[从\eqnrefs{eq:poincare-affine-matrix-dp75}到\eqnrefs{eq:poincare-mm}]

正文已经把 Poincar\'e 群写成仿射矩阵。这里仅完成从矩阵乘法到三组结构常数的代数，并用表示微分同态把抽象基送到厄米量子生成元。

取
\begin{equation*}
 \mathsf x(X,u)=
 \begin{pmatrix}X&u\\0&0\end{pmatrix},
 \qquad
 \mathsf x(Y,v)=
 \begin{pmatrix}Y&v\\0&0\end{pmatrix}.
\end{equation*}
直接相乘得到
\begin{align*}
 \mathsf x(X,u)\mathsf x(Y,v)
 &=\begin{pmatrix}XY&Xv\\0&0\end{pmatrix},\\
 \mathsf x(Y,v)\mathsf x(X,u)
 &=\begin{pmatrix}YX&Yu\\0&0\end{pmatrix},
\end{align*}
故
\begin{equation*}
 [\mathsf x(X,u),\mathsf x(Y,v)]
 =\mathsf x([X,Y],Xv-Yu),
\end{equation*}
即正文\eqnrefs{eq:poincare-semidir-lie-bracket-dp75}。

现在取 Lorentz 基 $m^{\mu\nu}=(\ell^{\mu\nu},0)$ 与平移基 $\tau^\rho=(0,e^\rho)$。两个纯平移元素的矩阵乘积在 Lie 括号中没有非零线性块，所以
\begin{equation*}
 [\tau^\mu,\tau^\nu]=0.
\end{equation*}
Lorentz 基对平移四矢量的作用由正文\eqnrefs{eq:lorentz-matrix-basis} 给出，
\begin{equation*}
 [m^{\mu\nu},\tau^\rho]
 =\eta^{\nu\rho}\tau^\mu-\eta^{\mu\rho}\tau^\nu.
\end{equation*}
而两个 Lorentz 基的矩阵对易子为
\begin{equation*}
[m^{\mu\nu},m^{\rho\sigma}]
=\eta^{\nu\rho}m^{\mu\sigma}
-\eta^{\mu\rho}m^{\nu\sigma}
-\eta^{\nu\sigma}m^{\mu\rho}
+\eta^{\mu\sigma}m^{\nu\rho}.
\end{equation*}
最后使用
\begin{equation*}
 dU(m^{\mu\nu})=-\frac{\ii}{\hbar}M^{\mu\nu},
 \qquad
 dU(\tau^\rho)=-\frac{\ii}{\hbar}P^\rho,
 \qquad
 dU([A,B])=[dU(A),dU(B)],
\end{equation*}
三条抽象 Lie 括号分别变成
\begin{align*}
 [P^\mu,P^\nu]&=0,\\
 [M^{\mu\nu},P^\rho]
 &=\ii\hbar(\eta^{\nu\rho}P^\mu-\eta^{\mu\rho}P^\nu),\\
 [M^{\mu\nu},M^{\rho\sigma}]
 &=\ii\hbar\left(
 \eta^{\nu\rho}M^{\mu\sigma}
 -\eta^{\mu\rho}M^{\nu\sigma}
 -\eta^{\nu\sigma}M^{\mu\rho}
 +\eta^{\mu\sigma}M^{\nu\rho}\right).
\end{align*}
这就是正文\eqnrange{eq:poincare-pp}{eq:poincare-mm}。轨道表达式 $L^{\mu\nu}=x^\mu P^\nu-x^\nu P^\mu$ 只是随后的一种具体表示，不参与一般 Poincar\'e 代数的建立。
\end{derivation}

\begin{derivation}[从\eqnrefs{eq:poincare-mm}到\eqnrefs{eq:JK-algebra}]

正文\eqnrefs{eq:poincare-mm} 已给出 $M^{\mu\nu}$ 的协变 Lorentz 代数。把六个独立分量按 $J_i=\epsilon_{ijk}M^{jk}/2$ 与 $K_i=M^{0i}$ 重新组织，并逐项代入协变指标关系，就得到正文\eqnrefs{eq:JK-algebra} 的三组 $J$--$K$ 对易子。


由定义
\begin{equation*}
 J_i=\frac12\epsilon_{ijk}M^{jk},
 \qquad K_i=M^{0i},
\end{equation*}
且反解为 $M^{ij}=\epsilon_{ijk}J_k$。

先取 $(\mu,\nu,\rho,\sigma)=(0,i,0,j)$：
\begin{align*}
 [K_i,K_j]
 &=[M^{0i},M^{0j}],\\
 &=\ii\hbar\left(
 \eta^{i0}M^{0j}-\eta^{00}M^{ij}
 -\eta^{ij}M^{00}+\eta^{0j}M^{i0}\right),\\
 &=-\ii\hbar M^{ij}
 =-\ii\hbar\epsilon_{ijk}J_k.
\end{align*}
这里使用 $\eta^{00}=1$、$\eta^{i0}=0$ 与 $M^{00}=0$。

转动与推动的对易子为
\begin{align*}
 [J_i,K_j]
 &=\frac12\epsilon_{iab}[M^{ab},M^{0j}],\\
 &=\frac{\ii\hbar}{2}\epsilon_{iab}
 (\delta_{bj}M^{a0}-\delta_{aj}M^{b0}),\\
 &=\ii\hbar\epsilon_{ijk}K_k.
\end{align*}

最后对纯空间指标，
\begin{equation*}
 [M^{ab},M^{cd}]
 =\ii\hbar(-\delta_{bc}M^{ad}+\delta_{ac}M^{bd}
 +\delta_{bd}M^{ac}-\delta_{ad}M^{bc}).
\end{equation*}
代入两个 $J$ 并使用
\begin{equation*}
 \epsilon_{iab}\epsilon_{jcb}
 =\delta_{ij}\delta_{ac}-\delta_{ic}\delta_{aj}
\end{equation*}
及其指标置换，可得
\begin{align*}
 [J_i,J_j]
 &=\frac14\epsilon_{iab}\epsilon_{jcd}[M^{ab},M^{cd}],\\
 &=\ii\hbar M^{ij}
 =\ii\hbar\epsilon_{ijk}J_k.
\end{align*}
三组分量关系因此完全由同一个协变 Lorentz 代数得到。
\end{derivation}
\subsection{Pauli--Lubanski 矢量}

非相对论自旋可看成三维角动量的内部部分，但相对论推动会把空间转动与时空推动相互混合，因此任何固定的三维分量都会依赖参考系。已经建立的四动量与反对称时空生成元允许构造一个与四动量正交的协变量，它把真正的内部角动量从参考系依赖的轨道分量中分离出来。

这个协变量称为 Pauli--Lubanski 赝四矢量，定义为
\begin{equation}
 \boxed{
 W^\mu\equiv-\frac12\epsilon^{\mu\nu\rho\sigma}
 P_\nu M_{\rho\sigma}.}
 \label{eq:pauli-lubanski}
\end{equation}
这里采用 $\epsilon^{0123}=+1$。反对称性立即限制它只能描述与四动量方向无关的内部角动量自由度：
\begin{equation}
 W^\mu P_\mu=0.
 \label{eq:W-orthogonality-dp34}
\end{equation}

Pauli--Lubanski 矢量的意义取决于它怎样受 Poincar\'e 变换作用。将其定义代入已经建立的 Poincar\'e 生成元代数可得
\begin{align}
 [W^\mu,P^\lambda]&=0,
 \label{eq:w-p}\\
 [M^{\alpha\beta},W^\mu]
 &=\ii\hbar\left(
 \eta^{\beta\mu}W^\alpha-
 \eta^{\alpha\mu}W^\beta\right),
 \label{eq:m-w}\\
 [W^\mu,W^\nu]
 &=\ii\hbar\epsilon^{\mu\nu\rho\sigma}W_\rho P_\sigma.
 \label{eq:w-w-dp34}
\end{align}
第一式说明 $W^\mu$ 不改变四动量本征值；第二式说明它按 Lorentz 四矢量变换；第三式则是由前两式与 Poincar\'e 代数直接推出的内部闭合关系，并不依赖先选择有质量粒子的静止系。

对质量 $m>0$ 的不可约表示，可以进入静止系
$P^\mu=(mc,\mathbf0)$。此时正交条件给 $W^0=0$，空间分量退化为静止系内部转动生成元：
\begin{equation}
 \boxed{
 W^\mu\big|_{\rm rest}
 =\left(0,\,mc\,\mathbf J_{\rm int}\right).}
 \label{eq:W-rest-r13}
\end{equation}
因此 Poincar\'e 群的两个独立二次 Casimir 算符可取 $P^2$ 与 $W^2$。在质量为 $m$、自旋为 $s$ 的不可约表示上，
\begin{equation}
 \boxed{
 P^2=m^2c^2,
 \qquad
 W^2=-m^2c^2\hbar^2s(s+1).}
 \label{eq:poincare-casimirs}
\end{equation}
这里 $s=0,\frac12,1,\ldots$ 是相对论意义下的粒子自旋。电子对应 $s=1/2$。因此寻找电子的局域波动方程，可以表述为：寻找一个局域场表示，使其一粒子激发落在 $P^2=m_{\mathrm e}^2c^2$ 且 $W^2=-3m_{\mathrm e}^2c^2\hbar^2/4$ 的 Poincar\'e 不可约表示中。

时空对称性先把质量与自旋固定为不可约表示的标签，具体 Hamiltonian 再决定这些自由度怎样传播和相互作用。Pauli--Lubanski 矢量 $W^\mu$ 由 Poincar\'e 生成元协变构造；对有质量表示，在静止参考动量 $p^\mu=(mc,\mathbf0)$ 上，它的空间分量化为 $mc$ 乘内部转动生成元。熟悉的 $SU(2)$ 自旋代数因此是这个协变结构在标准动量上的稳定子表示。


前面的 Pauli--Lubanski 算符已经告诉我们，质量和自旋可以作为单粒子态的 Lorentz 不变量标签。但还缺少一个更具体的问题：把粒子从一个惯性系变到另一个惯性系时，四动量会改变，内部自旋指标究竟怎样跟着变？

固定一条质量壳后，任意四动量 $p$ 都可由某个标准 Lorentz 变换从参考四动量 $\bar p$ 得到。若再施加一般 Lorentz 变换 $\Lambda$，把 $\Lambda p$ 重新拉回 $\bar p$ 所留下的复合变换会保持 $\bar p$ 不变，因此它不能改变参考动量，只能作用在内部指标上。所有保持 $\bar p$ 不变的 Lorentz 变换组成\emph{小群}。有质量粒子的自旋和无质量粒子的螺旋度，正是这个剩余对称群的不可约表示标签。


为使 Lorentz 变换与单粒子完备关系兼容，单粒子动量本征态采用协变归一化。对内部量子数 $r$，取
\begin{equation}
 \boxed{
 \langle \mathbf p',r'|\mathbf p,r\rangle_{\rm cov}
 =2p^0(2\pi\hbar)^3\delta_{r'r}
 \delta^{(3)}(\mathbf p'-\mathbf p),
 \qquad p^0=\frac{E_{\mathbf p}}{c},}
 \label{eq:rel-state-normalization-si-r9}
\end{equation}
并配套使用
\begin{equation}
 \boxed{
 \id(1)=\sum_r\int
 \frac{\dd^3 p}{(2\pi\hbar)^3\,2p^0}
 |\mathbf p,r\rangle_{\rm cov}
 {}_{\rm cov}\langle\mathbf p,r|.}
 \label{eq:rel-completeness-si-r9}
\end{equation}
这组归一化使测度 $\dd^3 p/(2p^0)$ 保持 Lorentz 不变，因此 Wigner 旋转的内部矩阵可以与动量测度分离，不会额外出现由非协变三维 $\delta$ 归一化产生的 Jacobian 平方根。

固定一条给定质量壳，并在这条轨道上选取标准动量 $\bar p^\mu$。对每个同轨道四动量 $p^\mu$，选一个标准 Lorentz 变换 $L(p)$ 满足
\begin{equation}
 L(p)\bar p=p.
 \label{eq:wigner-standard-boost-dp9}
\end{equation}
以 $|\bar p,\sigma\rangle$ 表示标准动量上的的内部基，定义
\begin{equation}
 |p,\sigma\rangle
 \equiv U[L(p)]|\bar p,\sigma\rangle.
 \label{eq:wigner-state-definition-dp9}
\end{equation}
$\sigma$ 标记固定标准动量时仍可变化的内部自由度，不增加时空坐标。对任意 Lorentz 变换 $\Lambda$，定义
\begin{equation}
 \boxed{
 W(\Lambda,p)
 \equiv L^{-1}(\Lambda p)\,\Lambda\,L(p),
 \qquad
 W(\Lambda,p)\bar p=\bar p.}
 \label{eq:wigner-little-group-element-dp68}
\end{equation}
因此 $W(\Lambda,p)$ 不再改变四动量，只能作用在内部指标上。若内部空间采用小群不可约表示 $D$，单粒子态的完整 Lorentz 变换为
\begin{equation}
 \boxed{
 U(\Lambda)|p,\sigma\rangle
 =\sum_{\sigma'}
 D_{\sigma'\sigma}[W(\Lambda,p)]
 |\Lambda p,\sigma'\rangle.}
 \label{eq:wigner-state-transform-dp68}
\end{equation}

\begin{figure}[htbp]
 \centering
 \resizebox{0.92\textwidth}{!}{\begin{tikzpicture}[x=1cm,y=1cm,>=Stealth,
 every node/.style={font=\large},
 box/.style={draw,rounded corners=2pt,inner xsep=8pt,inner ysep=5pt,align=center,line width=.75pt},
 line label/.style={fill=white,fill opacity=.96,text opacity=1,inner sep=1.2pt}]
\node[box] (k) at (-4.5,0) {standard momentum\\$\bar p$};
\node[box] (p) at (-1.2,2.0) {$p=L(p)\bar p$};
\node[box] (lp) at (3.1,2.0) {$\Lambda p$};
\node[box] (kb) at (3.1,-1.0) {$\bar p$\\same orbit reference};

\draw[->,line width=.85pt] (k)--node[line label,above left] {$L(p)$}(p);
\draw[->,line width=.85pt] (p)--node[line label,above] {$\Lambda$}(lp);
\draw[->,line width=.85pt] (lp)--node[line label,right] {$L^{-1}(\Lambda p)$}(kb);
\draw[->,line width=.95pt,dashed] (k)--node[line label,below] {$W(\Lambda,p)$}(kb);

\node[box,align=center,font=\normalsize] (stab) at (-0.7,-2.65)
  {$G_{\bar p}=\{\Lambda\mid \Lambda\bar p=\bar p\}$\\stabilizer of the reference momentum};
\node[box] (massive) at (-3.2,-4.65) {$\bar p^2=m^2c^2>0$\\$G_{\bar p}\simeq SO(3)$\\double cover $SU(2)$\\spin $s$};
\node[box] (massless) at (3.4,-4.65) {$\bar k^2=0$\\$G_{\bar k}\simeq ISO(2)$\\discrete-helicity sector\\helicity $\lambda$};
\draw[->,line width=.8pt] (stab.south west) --node[line label,left,font=\normalsize]{timelike orbit}(massive.north);
\draw[->,line width=.8pt] (stab.south east) --node[line label,right,font=\normalsize]{null orbit}(massless.north);
\end{tikzpicture}
}
 \caption{Wigner 小群构造。$L(p)$ 把标准动量推到 $p$，而 $W(\Lambda,p)=L^{-1}(\Lambda p)\Lambda L(p)$ 把复合变换拉回同一个标准动量，因此只作用在内部自旋/螺旋度空间。}
 \label{fig:wigner-little-group-dp9}
\end{figure}

对 $m>0$，最方便的标准动量就是静止四动量 $\bar p_m^\mu=(mc,\mathbf0)$。保持它不变的 Lorentz 变换不能混合时间与空间分量，只能在三维空间做转动，因此
\begin{equation}
 \boxed{
 G_{\bar p_m}\simeq SO(3),
 \qquad
 \widetilde G_{\bar p_m}\simeq SU(2).}
 \label{eq:wigner-massive-little-group-dp68}
\end{equation}
量子态使用双覆盖 $SU(2)$，所以其不可约表示由 $s=0,\tfrac12,1,\ldots$ 标记，并具有 $2s+1$ 个自旋态。这与前文 $W^2=-m^2c^2\hbar^2s(s+1)$ 的 Casimir 分类是同一件事的两种表述。

无质量轨道没有静止系，因此不能把有质量构造简单令 $m\to0$。选沿 $+z$ 传播的标准零质量动量
\begin{equation}
 \bar k^\mu=(\kappa_\star,0,0,\kappa_\star),
 \qquad \kappa_\star>0,
 \qquad \bar k^2=0,
 \label{eq:wigner-massless-standard-momentum-dp9}
\end{equation}
其中 $\kappa_\star$ 具有物理动量量纲，其数值只是在零质量轨道上选参考点。全部稳定子方向由满足 $X\bar k=0$ 的 Lorentz Lie 代数元给出。写一般元
\begin{equation*}
 X=\sum_{i=1}^3(a_iJ_i+b_iK_i)
\end{equation*}
出发，并要求它在 $P^\mu=\bar k^\mu$ 上满足 $[X,P^\mu]=0$。四个分量的线性条件化为
\begin{equation*}
 b_3=0,
 \qquad
 b_1=-a_2,
 \qquad
 b_2=a_1,
\end{equation*}
所以稳定子 Lie 代数恰为三维。取一组方便的基为
\begin{equation}
 \mathsf N_1\equiv K_1-J_2,
 \qquad
 \mathsf N_2\equiv K_2+J_1,
 \qquad
 J_3.
 \label{eq:wigner-massless-little-group-generators-dp68}
\end{equation}
它们满足
\begin{equation}
 \boxed{
 [J_3,\mathsf N_1]=\ii\hbar\mathsf N_2,
 \qquad
 [J_3,\mathsf N_2]=-\ii\hbar\mathsf N_1,
 \qquad
 [\mathsf N_1,\mathsf N_2]=0,}
 \label{eq:wigner-massless-little-group-algebra-dp68}
\end{equation}
这正是二维平面转动加平移的代数，因此在 $SO^+(1,3)$ 中零质量标准动量的小群同构于 $ISO(2)$。这里的“平移”作用在内部表示空间，不是把粒子的时空位置平移。量子态若由 Lorentz 群的双覆盖 $SL(2,\mathbb C)$ 实现，则实际作用的是 $ISO(2)$ 的相应双覆盖；这一区分对半整数螺旋度是必要的。

这个小群的量子不可约表示有两类。若 $\mathsf N_{1,2}$ 非平凡作用，会得到连续自旋表示；若这两个“平移型”生成元在物理态上平凡作用，则只剩绕运动方向的转动生成元 $J_3$。这种第二类表示由一个离散或半整数的螺旋度标签 $\lambda$ 区分，并满足
\begin{equation}
 \mathsf N_1|\bar k,\lambda\rangle=0,
 \qquad
 \mathsf N_2|\bar k,\lambda\rangle=0,
 \qquad
 J_3|\bar k,\lambda\rangle=\hbar\lambda|\bar k,\lambda\rangle.
 \label{eq:massless-helicity-irrep-dp9}
\end{equation}
此时 Pauli--Lubanski 矢量与四动量平行：
\begin{equation}
 \boxed{W^\mu|k,\lambda\rangle
 =\hbar\lambda\,P^\mu|k,\lambda\rangle.}
 \label{eq:massless-pauli-lubanski-helicity-dp9}
\end{equation}
因此无质量表示的 $W^2=0$ 对所有离散螺旋度都相同；真正区分内部态的是小群转动生成元本征值 $\lambda$。光子的物理 Hilbert 空间取 $\lambda=\pm1$，无质量自旋 $1/2$ 场取 $\lambda=\pm\tfrac12$。光子的 $\lambda=0$ 不属于物理横向单光子空间；后面 Maxwell 约束与规范商空间会从场变量层面再次得到同一结论。

三个常被混写的对象属于不同的表示层级：$(j_L,j_R)$ 是\emph{局域场}在有限维 Lorentz 复表示中的变换类型；$s$ 是有质量单粒子小群双覆盖 $SU(2)$ 的不可约标签；$\lambda$ 是无质量离散表示中稳定子转动方向的本征值。对整数自旋可直接把它看成 $SO(2)$ 角色标签；对半整数自旋必须在 Lorentz 双覆盖中的小群上理解它。一个四矢量场属于 $(\tfrac12,\tfrac12)$ 并不意味着物理光子有四个自旋态；规范约束最终只留下 $\lambda=\pm1$ 两个单粒子螺旋度态。


\begin{derivation}[从\eqnrefs{eq:wigner-state-definition-dp9}到\eqnrefs{eq:wigner-state-transform-dp68}]
对任意保向正时向 Lorentz 变换 $\Lambda$，定义

\begin{equation*}
 {
 W(\Lambda,p)
 \equiv L^{-1}(\Lambda p)\,\Lambda\,L(p).}
 \end{equation*}
直接作用在标准动量上：
\begin{align*}
 W(\Lambda,p)\bar p
 &=L^{-1}(\Lambda p)\Lambda L(p)\bar p,\\
 &=L^{-1}(\Lambda p)\Lambda p,\\
 &=\bar p.
 \end{align*}
因此 $W(\Lambda,p)$ 确实属于稳定 $\bar p$ 的小群
\begin{equation*}
 G_{\bar p}\equiv\{W\in SO^+(1,3)\mid W\bar p=\bar p\}.
 \end{equation*}
再把 \eqnrefs{eq:wigner-state-definition-dp9} 代入态变换：
\begin{align*}
 U(\Lambda)|p,\sigma\rangle
 &=U(\Lambda)U[L(p)]|\bar p,\sigma\rangle,\\
 &=U[L(\Lambda p)]U[W(\Lambda,p)]|\bar p,\sigma\rangle.
\end{align*}
因为 $W$ 不改变 $\bar p$，它只能在固定动量的内部空间中作用。若该内部空间取小群不可约表示 $D$，
\begin{equation*}
 U(W)|\bar p,\sigma\rangle
 =\sum_{\sigma'}D_{\sigma'\sigma}(W)|\bar p,\sigma'\rangle,
\end{equation*}
于是
\begin{equation*}
 {
 U(\Lambda)|p,\sigma\rangle
 =\sum_{\sigma'}
 D_{\sigma'\sigma}[W(\Lambda,p)]
 |\Lambda p,\sigma'\rangle.}
 \end{equation*}
这条式子把四动量的 Lorentz 变换与内部自旋/螺旋度变换严格分离。若改用非协变三维 $\delta$ 归一化，右端还会出现由雅可比行列式产生的 $\sqrt{(\Lambda p)^0/p^0}$ 因子；后续场量子化采用协变归一化，因此不再引入这一额外因子。
\end{derivation}
\begin{derivation}[从\eqnrefs{eq:wigner-little-group-element-dp68}到\eqnrefs{eq:wigner-massive-little-group-dp68}]
有质量粒子可选静止标准四动量
\begin{equation*}
 \bar p_m^\mu=(mc,0,0,0).
\end{equation*}
小群元素 $W$ 的定义是 $W\bar p_m=\bar p_m$。逐分量写出
\begin{equation*}
 W^\mu{}_{\nu}\bar p_m^\nu
 =mc\,W^\mu{}_0
 =mc\,\delta^\mu{}_0,
\end{equation*}
所以第一列被完全固定：
\begin{equation*}
 W^0{}_0=1,
 \qquad
 W^i{}_0=0.
\end{equation*}
还需要证明第一行也没有 boost 分量。Lorentz 条件
$W^{\mathsf T}\eta W=\eta$ 的 $(0,i)$ 分量给
\begin{equation*}
 \eta_{\alpha\beta}W^\alpha{}_0W^\beta{}_i=0.
\end{equation*}
代入 $W^\alpha{}_0=\delta^\alpha{}_0$ 后得到
$W^0{}_i=0$。因此 $W$ 必须具有块形式
\begin{equation*}
 W=
 \begin{pmatrix}
 1&\mathbf0^{\mathsf T}\\
 \mathbf0&R
 \end{pmatrix}.
\end{equation*}
再把这个块矩阵代回 Lorentz 条件，空间块满足
\begin{equation*}
 R^{\mathsf T}R=I_3.
\end{equation*}
若只考虑与恒等变换连续相连的 proper orthochronous Lorentz 群，还要求
$\det W=+1$，因此 $\det R=+1$。所以有质量标准动量的连续稳定子正是
$SO(3)$。

量子态的自旋不是用 $SO(3)$ 的普通表示分类，而是用其双覆盖 $SU(2)$ 的幺正不可约表示分类，因为 $2\pi$ 空间转动可以在态矢量上产生负号而不改变物理射线。$SU(2)$ 不可约表示由
$s=0,\tfrac12,1,\ldots$ 标记，维数为 $2s+1$；这正是正文\eqnrefs{eq:wigner-massive-little-group-dp68} 中出现的自旋量子数，也与 Pauli--Lubanski Casimir 的标签一致。
\end{derivation}
\begin{derivation}[从\eqnrefs{eq:wigner-massless-standard-momentum-dp9}到\eqnrefs{eq:wigner-massless-little-group-algebra-dp68}]

稳定子代数的完备性由一般 Lorentz 代数元的线性稳定条件决定。先解出全部稳定方向，再计算所得三维解空间的 Lie 括号。

写
\begin{equation*}
 X=a_1J_1+a_2J_2+a_3J_3
 +b_1K_1+b_2K_2+b_3K_3.
\end{equation*}
由正文\eqnrefs{eq:poincare-mp}，
\begin{equation*}
 [J_i,P^0]=0,
 \qquad
 [J_i,P^j]=\ii\hbar\epsilon_{ijk}P^k,
\end{equation*}
以及
\begin{equation*}
 [K_i,P^0]=-\ii\hbar P^i,
 \qquad
 [K_i,P^j]=-\ii\hbar\delta_{ij}P^0.
\end{equation*}
在
\begin{equation*}
 \bar k^\mu=(\kappa_\star,0,0,\kappa_\star)
\end{equation*}
上逐分量要求 $[X,P^\mu]_{P=\bar k}=0$。时间分量给
\begin{equation*}
 -\ii\hbar\kappa_\star b_3=0
 \quad\Longrightarrow\quad b_3=0.
\end{equation*}
第一空间分量给
\begin{equation*}
 -\ii\hbar\kappa_\star(a_2+b_1)=0
 \quad\Longrightarrow\quad b_1=-a_2,
\end{equation*}
第二空间分量给
\begin{equation*}
 \ii\hbar\kappa_\star(a_1-b_2)=0
 \quad\Longrightarrow\quad b_2=a_1.
\end{equation*}
第三空间分量在这些条件下自动为零。因此任意稳定子元都唯一写成
\begin{equation*}
 X=a_3J_3+a_1(K_2+J_1)-a_2(K_1-J_2),
\end{equation*}
证明稳定子空间恰为
\begin{equation*}
 \operatorname{span}\{J_3,\mathsf N_1,\mathsf N_2\}.
\end{equation*}

再用正文\eqnrefs{eq:JK-algebra}：
\begin{align*}
 [J_3,\mathsf N_1]
 &=[J_3,K_1]-[J_3,J_2]
 =\ii\hbar(K_2+J_1)
 =\ii\hbar\mathsf N_2,\\
 [J_3,\mathsf N_2]
 &=[J_3,K_2]+[J_3,J_1]
 =-\ii\hbar(K_1-J_2)
 =-\ii\hbar\mathsf N_1,\\
 [\mathsf N_1,\mathsf N_2]
 &=[K_1,K_2]+[K_1,J_1]-[J_2,K_2]-[J_2,J_1]\\
 &=-\ii\hbar J_3+\ii\hbar J_3\\
 &=0.
\end{align*}
这就是正文\eqnrefs{eq:wigner-massless-little-group-algebra-dp68}。由于三维解空间已经由一般稳定条件完整求出，所得 $\mathfrak{iso}(2)$ 不再只是三个候选生成元的子代数，而是零质量标准动量的全部连续稳定子 Lie 代数。
\end{derivation}




\begin{derivation}[从\eqnrefs{eq:pauli-lubanski}到\eqnrefs{eq:m-w}]
正文把 $W^\mu$ 称为赝四矢量，因此必须验证其四个分量确实按 Lorentz 四矢量混合。对连通保向 Lorentz 变换，$\det\Lambda=1$，完全反对称张量满足
\begin{equation*}
 \epsilon^{\mu\nu\rho\sigma}
 \Lambda_\nu{}^{\alpha}\Lambda_\rho{}^{\beta}\Lambda_\sigma{}^{\gamma}
 =\Lambda^\mu{}_{\delta}\epsilon^{\delta\alpha\beta\gamma}.
\end{equation*}
四动量与 Lorentz 生成元的有限伴随变换分别为
\begin{equation*}
 U^{-1}P_\nu U=\Lambda_\nu{}^{\alpha}P_\alpha,
 \qquad
 U^{-1}M_{\rho\sigma}U
 =\Lambda_\rho{}^{\beta}\Lambda_\sigma{}^{\gamma}M_{\beta\gamma}.
\end{equation*}
于是
\begin{align*}
 U^{-1}W^\mu U
 &=-\frac12\epsilon^{\mu\nu\rho\sigma}
 \bigl(U^{-1}P_\nu U\bigr)
 \bigl(U^{-1}M_{\rho\sigma}U\bigr),\\
 &=-\frac12\epsilon^{\mu\nu\rho\sigma}
 \Lambda_\nu{}^{\alpha}\Lambda_\rho{}^{\beta}\Lambda_\sigma{}^{\gamma}
 P_\alpha M_{\beta\gamma},\\
 &=\Lambda^\mu{}_{\delta}W^\delta.
\end{align*}
再令 $\Lambda^\mu{}_{\nu}=\delta^\mu{}_{\nu}+\omega^\mu{}_{\nu}+\order(\omega^2)$，并展开
\begin{equation*}
 U=\id()-\frac{\ii}{2\hbar}\omega_{\alpha\beta}M^{\alpha\beta}
 +\order(\omega^2),
\end{equation*}
比较任意反对称参数 $\omega_{\alpha\beta}$ 的一阶系数，得到
\begin{equation*}
 [M^{\alpha\beta},W^\mu]
 =\ii\hbar\left(\eta^{\beta\mu}W^\alpha-
 \eta^{\alpha\mu}W^\beta\right),
\end{equation*}
即正文\eqnrefs{eq:m-w}。若把宇称也包括进不连通的 Lorentz 群，\eqnrefs{eq:pauli-lubanski} 中的 $\epsilon^{\mu\nu\rho\sigma}$ 使 $W^\mu$ 具有赝矢量的额外宇称符号。
\end{derivation}


\begin{derivation}[从\eqnrefs{eq:w-p}到\eqnrefs{eq:w-w-dp34}]

直接从 $W^\mu$ 的定义出发，并使用 $[P_\alpha,W^\nu]=0$：
\begin{align*}
[W^\mu,W^\nu]
&=-\frac12\epsilon^{\mu\alpha\rho\sigma}
[P_\alpha M_{\rho\sigma},W^\nu]\\
&=-\frac12\epsilon^{\mu\alpha\rho\sigma}
P_\alpha[M_{\rho\sigma},W^\nu].
\end{align*}
由\eqnrefs{eq:m-w} 降低生成元指标，
\begin{equation*}
[M_{\rho\sigma},W^\nu]
=\ii\hbar(\delta_\sigma^\nu W_\rho-\delta_\rho^\nu W_\sigma).
\end{equation*}
代回并合并两个反对称项：
\begin{align*}
[W^\mu,W^\nu]
&=\ii\hbar\,\epsilon^{\mu\alpha\nu\rho}P_\alpha W_\rho\\
&=\ii\hbar\,\epsilon^{\mu\nu\rho\sigma}W_\rho P_\sigma,
\end{align*}
其中最后一步只重排 Levi--Civita 指标并使用 $[W_\rho,P_\sigma]=0$。这就是正文\eqnrefs{eq:w-w-dp34}，整个推导没有选择静止系或特定质量轨道。
\end{derivation}

\begin{derivation}[从\eqnrefs{eq:W-rest-r13}到\eqnrefs{eq:poincare-casimirs}]

Casimir 算符的定义不是“在某个参考系里数值不变”，而是它与整个 Lie 代数的所有生成元对易。先证明
\[
C_1\equiv P_\mu P^\mu
\]
满足这一要求。平移生成元彼此对易，所以
\[
[P^\lambda,C_1]=0.
\]
再用 Poincar\'e 代数
\[
[M^{\alpha\beta},P^\mu]
=\ii\hbar
\left(
\eta^{\beta\mu}P^\alpha
-\eta^{\alpha\mu}P^\beta
\right)
\]
计算
\begin{align*}
[M^{\alpha\beta},P_\mu P^\mu]
&=[M^{\alpha\beta},P_\mu]P^\mu
+P_\mu[M^{\alpha\beta},P^\mu]\\
&=\ii\hbar
\left(
P^\alpha P^\beta-P^\beta P^\alpha
\right)
+\ii\hbar
\left(
P^\beta P^\alpha-P^\alpha P^\beta
\right)\\
&=0.
\end{align*}
因此 \(P^2\) 与全部平移、Lorentz 生成元对易，是第一个 Casimir。

第二个 Casimir 从 Pauli--Lubanski 四矢量
\[
W^\mu
\equiv
\frac12\epsilon^{\mu\nu\rho\sigma}
P_\nu M_{\rho\sigma}
\]
构造。先看它与平移的对易子：
\begin{align*}
[P^\lambda,W^\mu]
&=\frac12\epsilon^{\mu\nu\rho\sigma}
P_\nu[P^\lambda,M_{\rho\sigma}]\\
&=-\frac{\ii\hbar}{2}
\epsilon^{\mu\nu\rho\sigma}
P_\nu
\left(
\delta^\lambda_{\ \sigma}P_\rho
-\delta^\lambda_{\ \rho}P_\sigma
\right).
\end{align*}
每一项都含对称乘积 \(P_\nu P_\rho\) 与完全反对称的
\(\epsilon^{\mu\nu\rho\sigma}\) 收缩，因此
\[
[P^\lambda,W^\mu]=0.
\]
同样，从 \([M,M]\) 与 \([M,P]\) 两部分逐项展开可得
\[
[M^{\alpha\beta},W^\mu]
=\ii\hbar
\left(
\eta^{\beta\mu}W^\alpha
-\eta^{\alpha\mu}W^\beta
\right),
\]
即 \(W^\mu\) 确实按 Lorentz 四矢量变换。于是
\begin{align*}
[M^{\alpha\beta},W_\mu W^\mu]
&=[M^{\alpha\beta},W_\mu]W^\mu
+W_\mu[M^{\alpha\beta},W^\mu]\\
&=0,
\end{align*}
而 \([P^\lambda,W^2]=0\) 已由 \([P,W]=0\) 立即得到。因此
\[
C_2\equiv W_\mu W^\mu
\]
也是 Poincar\'e 代数的 Casimir。

还可直接验证
\[
W^\mu P_\mu
=
\frac12
\epsilon^{\mu\nu\rho\sigma}
P_\mu P_\nu M_{\rho\sigma}
=0,
\]
因为 \(P_\mu P_\nu\) 对 \(\mu,\nu\) 对称。对有质量不可约表示 \(m>0\)，选择静止系
\[
P^\mu=(mc,\mathbf0).
\]
正文\eqnrefs{eq:W-rest-r13} 给出
\[
W^0=0,
\qquad
\mathbf W=mc\,\mathbf J_{\rm int},
\]
其中 \(\mathbf J_{\rm int}\) 是静止小群 \(SO(3)\) 的内部角动量生成元。于是
\[
W^2
=(W^0)^2-\mathbf W^2
=-m^2c^2\,\mathbf J_{\rm int}^2.
\]
不可约自旋 \(s\) 表示满足
\[
\mathbf J_{\rm int}^2
=\hbar^2s(s+1),
\]
所以
\[
\boxed{
P^2=m^2c^2,
\qquad
W^2=-m^2c^2\hbar^2s(s+1).
}
\]
这就是正文\eqnrefs{eq:poincare-casimirs}。质量与自旋之所以能标记 Poincar\'e 不可约表示，不是由静止系经验“读出来”的，而是因为对应算符先由代数证明为 Casimir，再在静止系评价其本征值。
\end{derivation}

\subsection{二分量与四分量 Lorentz 表示}

质量和自旋已经分类了单粒子态，但局域波动方程还需要回答另一个问题：场在每个时空点应有多少个分量，以及这些分量在换惯性系时怎样混合。这个问题属于 Lorentz 群的有限维表示，而不是单粒子 Hilbert 空间的不可约表示。局域场的内部分量必须按某个有限维矩阵表示线性变换，并保持群复合律；由此才能构造协变的一阶或二阶微分方程。

Lorentz 代数可重新组合成两套彼此独立的角动量代数，由此产生两种互不等价的二分量表示，即左手和右手 Weyl 表示。对有质量自旋 $1/2$ 粒子，质量项必须把这两种表示相互耦合，它们的直和构成四分量 Dirac 表示。于是 Weyl 旋量描述两个不可约手征表示，Dirac 旋量则把有质量理论所需的左右手分量组合在同一场中。

% Pauli/Weyl algebra targets. Detailed matrix multiplication is in the owner proof zone.
后面描述自旋 $1/2$ 场时会遇到一种具有两个复分量、并按 Lorentz 群二维复表示变换的对象；这种二分量对象称为 Weyl 旋量。Weyl 旋量所属的分量空间就是二维复线性空间。任意 $2\times2$ 复矩阵都可由单位矩阵与三个无迹厄米矩阵张成；要求这三个生成元闭合成 $SU(2)$ 转动代数后，自然选取 Pauli 矩阵
\begin{equation}
 \sigma_1=\begin{pmatrix}0&1\\1&0\end{pmatrix},\qquad
 \sigma_2=\begin{pmatrix}0&-\ii\\ \ii&0\end{pmatrix},\qquad
 \sigma_3=\begin{pmatrix}1&0\\0&-1\end{pmatrix}.
 \label{eq:pauli-explicit-r10}
\end{equation}
它们的乘法闭合为
\begin{equation}
 \boxed{
 \sigma_i\sigma_j
 =\delta_{ij}\id_2+\ii\epsilon_{ijk}\sigma_k,}
 \label{eq:pauli-product-proof-r10}
\end{equation}
因而对称与反对称部分分别满足
\begin{equation}
 \boxed{
 \{\sigma_i,\sigma_j\}=2\delta_{ij}\id_2,
 \qquad
 [\sigma_i,\sigma_j]=2\ii\epsilon_{ijk}\sigma_k.}
 \label{eq:pauli-comm-anticomm-r10}
\end{equation}

对单位向量 $\hat{\mathbf n}$，\eqnrefs{eq:pauli-product-proof-r10} 立即给出
$(\hat{\mathbf n}\cdot\boldsymbol\sigma)^2=\id_2$。因此有限转动的矩阵指数能够按奇偶幂精确求和：
\begin{equation}
 \boxed{
 \ee^{-\ii\theta\hat{\mathbf n}\cdot\boldsymbol\sigma/2}
 =\id_2\cos\frac{\theta}{2}
 -\ii\hat{\mathbf n}\cdot\boldsymbol\sigma\sin\frac{\theta}{2}.}
 \label{eq:pauli-exponential-r68}
\end{equation}
把纯虚系数 $-\ii\theta$ 换成实快度 $\pm\zeta$，同一代数同时给出 Weyl 推动中的双曲函数形式。

三维反对称张量的缩并恒等式取为
\begin{equation}
 \boxed{
 \epsilon_{ijk}\epsilon_{ilm}
 =\delta_{jl}\delta_{km}-\delta_{jm}\delta_{kl},
 \qquad
 \epsilon_{ijk}\epsilon_{ijm}=2\delta_{km}.}
 \label{eq:epsilon-contractions-r10}
\end{equation}
这些关系随后把三维自旋代数嵌入四维 Lorentz 结构。定义
\begin{equation}
 \sigma^\mu=(\id_2,\boldsymbol\sigma),
 \qquad
 \bar\sigma^\mu=(\id_2,-\boldsymbol\sigma),
 \label{eq:sigma-barsigma-foundation-r10}
\end{equation}
则 Pauli 乘法定律等价于四维交织恒等式
\begin{equation}
 \boxed{
 \sigma^\mu\bar\sigma^\nu+\sigma^\nu\bar\sigma^\mu
 =2\eta^{\mu\nu}\id_2,}
 \label{eq:sigmabar-proof-r10}
\end{equation}
交换 $\sigma$ 与 $\bar\sigma$ 的次序也成立。\eqnrefs{eq:sigmabar-proof-r10} 说明这组矩阵的对称反对易部分恰好重现 Minkowski 度规。这里出现这种关系并不是为了再引入一个新的对称群，而是为了完成一个具体的代数任务：相对论自由粒子的质量壳条件 $p^2=m^2c^2$ 对四动量是二次的；若希望以后把运动方程写成对四动量的一阶形式，就需要一组常矩阵 $\Gamma^\mu$，使一阶算符 $\Gamma^\mu p_\mu$ 的平方只留下 Lorentz 标量 $p^2$。这要求矩阵的对称乘积满足
\begin{equation*}
 \{\Gamma^\mu,\Gamma^\nu\}=2\eta^{\mu\nu}I.
\end{equation*}
由满足这组反对易关系的生成元及其有限乘积构成的结合代数，称为 Minkowski 度规对应的 \emph{Clifford 代数}。因此 Clifford 代数是一种把二次 Lorentz 不变量“线性化”的结合矩阵代数；时空对称变换仍由 Lorentz 群及其自旋表示描述。\eqnrefs{eq:sigmabar-proof-r10} 已经在两个二分量空间之间实现了同一种代数结构，后面把左右两个二分量空间合并后，就会得到作用在四分量空间上的 $\gamma^\mu$。


Lorentz 李代数的六个实生成元可以重新组合成两组彼此对易的三生成元代数。定义
\begin{equation}
 \mathsf J_i^{(+)}\equiv\frac12(J_i+\ii K_i),
 \qquad
 \mathsf J_i^{(-)}\equiv\frac12(J_i-\ii K_i),
 \label{eq:lorentz-su2-generators-r15}
\end{equation}
则
\begin{equation}
 [\mathsf J_i^{(\pm)},\mathsf J_j^{(\pm)}]
 =\ii\hbar\epsilon_{ijk}\mathsf J_k^{(\pm)},
 \qquad
 [\mathsf J_i^{(+)},\mathsf J_j^{(-)}]=0.
 \label{eq:lorentz-su2-algebra-dp34}
\end{equation}
\eqnrefs{eq:lorentz-su2-algebra-dp34} 中已经出现系数 $\ii$，所以现在允许用复数而不只是实数来线性组合原来的六个 Lorentz 生成方向。这个操作称为实 Lie 代数的\emph{复化}：若原代数是实向量空间 $\mathfrak g$，复化就是把标量域从 $\mathbb R$ 扩充到 $\mathbb C$，记作 $\mathfrak g_{\mathbb C}=\mathfrak g\otimes_{\mathbb R}\mathbb C$。它只是为了分类复矩阵表示而扩充线性组合系数，并没有增加新的物理 Lorentz 变换。于是这里分解的是复化后的代数，而不是 $\mathfrak{so}(1,3)$ 作为实 Lie 代数的直和。更准确地，
\begin{equation}
 \boxed{
 \mathfrak{so}(1,3)_\mathbb C
 \simeq
 \mathfrak{sl}(2,\mathbb C)\oplus\mathfrak{sl}(2,\mathbb C).}
 \label{eq:lorentz-complexified-decomposition-dp75}
\end{equation}
两套 $\mathsf J_i^{(\pm)}$ 具有角动量型括号，所以有限维\emph{复}不可约表示可以用 $(j_L,j_R)$ 标记。实 Lorentz 代数仍是非紧致的实简单代数；$(j_L,j_R)$ 分类来自它的复表示论。

四矢量表示给出时空坐标在 Lorentz 变换下的线性混合规律；自旋 $1/2$ 粒子的量子态还必须容纳一种四矢量表示看不见的性质：空间转动 $2\pi$ 后，所有四矢量已经回到原值，但量子态仍可以多出一个负号。因此仅用 $4\times4$ Lorentz 矩阵作用在四矢量上，还不足以描述旋量怎样随转动和推动变化。因此需要构造一组 $2\times2$ 复矩阵，使每个保向、正时向 Lorentz 变换对应其中两个相差整体负号的矩阵。这种“两个矩阵对应同一个 Lorentz 变换”的二对一关系称为\emph{双覆盖}，相应矩阵群就是 $SL(2,\mathbb C)$。


把实四矢量 $x^\mu=(x^0,\mathbf x)$ 与厄米 $2\times2$ 矩阵一一对应：
\begin{equation}
 X(x)\equiv x^0\id(2)+\mathbf x\cdot\boldsymbol\sigma_{\rm P}
 =\begin{pmatrix}
 x^0+x^3 & x^1-\ii x^2\\
 x^1+\ii x^2 & x^0-x^3
 \end{pmatrix}.
 \label{eq:sl2c-hermitian-map-dp8}
\end{equation}
这个对应之所以有用，是因为矩阵行列式恰好等于 Minkowski 范数：
\begin{equation}
 \boxed{\det X(x)=x_\mu x^\mu.}
 \label{eq:sl2c-det-minkowski-dp8}
\end{equation}
因此，只要找到保持 $\det X$ 的矩阵变换，就自动得到保持时空间隔的变换。

对任意 $A\in SL(2,\mathbb C)$，定义
\begin{equation}
 X' = A X A^\dagger.
 \label{eq:sl2c-action-X-dp8}
\end{equation}
$X'$ 仍为厄米矩阵，因而唯一对应某个实四矢量 $x'^\mu$；同时 $\det A=1$ 保证 $\det X'=\det X$。所以每个 $A$ 都诱导一个保向、正时向 Lorentz 变换。另一方面，$A$ 与 $-A$ 对 $X$ 的作用完全相同，于是得到二对一群同态
\begin{equation}
 \boxed{SL(2,\mathbb C)\longrightarrow SO^+(1,3),\qquad A\sim -A.}
 \label{eq:sl2c-double-cover-dp8}
\end{equation}
这就是所谓“双覆盖”：同一个 Lorentz 变换对应 $SL(2,\mathbb C)$ 中相差负号的两个元素。

这里必须区分实 Lie 代数与其复化。实 Lorentz 代数 $\mathfrak{so}(1,3)$ 本身是实简单 Lie 代数，并不等于两份 $\mathfrak{su}(2)$ 的实直和。为了分类有限维复表示，把生成元允许作复线性组合；此时
\begin{equation*}
 \mathfrak{so}(1,3)_\mathbb C
 \simeq \mathfrak{sl}(2,\mathbb C)\oplus\mathfrak{sl}(2,\mathbb C),
\end{equation*}
等价地可用 $\mathsf J_i^{(\pm)}=(J_i\pm\ii K_i)/2$ 得到两套彼此对易、具有 $\mathfrak{su}(2)$ 型结构常数的复子代数。有限维复不可约表示因此由两个角动量型标签 $(j_L,j_R)$ 标记，其复维数为
\begin{equation}
 \boxed{\dim(j_L,j_R)=(2j_L+1)(2j_R+1).}
 \label{eq:lorentz-irrep-dimension-dp8}
\end{equation}
最常用的表示由此统一写成
\begin{equation}
 \begin{array}{ccl}
 (0,0) &:& \text{标量},\\[2pt]
 (\tfrac12,0) &:& \text{左手 Weyl 旋量},\\[2pt]
 (0,\tfrac12) &:& \text{右手 Weyl 旋量},\\[2pt]
 (\tfrac12,\tfrac12) &:& \text{四矢量},\\[2pt]
 (1,0)\oplus(0,1) &:& \text{反对称二阶张量},\\[2pt]
 (\tfrac12,0)\oplus(0,\tfrac12) &:& \text{Dirac 旋量}.
 \end{array}
 \label{eq:lorentz-irrep-dictionary-dp8}
\end{equation}
这里的 $(j_L,j_R)$ 描述的是\emph{局域场怎样变换}；它与后面由 Poincar\'e 不可约表示得到的单粒子自旋标签相关，但不是同一个概念。


表示的张量积描述复合 Lorentz 指标怎样重新组合，而不是把两个粒子的自旋数简单相加。左右 Weyl 指标各取一次时，维数 $2\times2=4$，并且
\begin{equation}
 \boxed{
 \left(\frac12,0\right)\otimes\left(0,\frac12\right)
 =\left(\frac12,\frac12\right),
 \qquad
 V_{\alpha\dot\beta}=V_\mu(\sigma^\mu)_{\alpha\dot\beta}.}
 \label{eq:weyl-vector-product-r68}
\end{equation}
所以一个带左右各一个 Weyl 指标的 $2\times2$ 厄米对象正好按四矢量表示变换。若两个指标具有相同手征性，则
\begin{equation}
 \boxed{
 \left(\frac12,0\right)\otimes\left(\frac12,0\right)
 =(0,0)\oplus(1,0),
 \qquad
 \left(0,\frac12\right)\otimes\left(0,\frac12\right)
 =(0,0)\oplus(0,1).}
 \label{eq:weyl-same-chirality-product-r68}
\end{equation}
四维反对称二阶张量的六个分量因此分成两个三维手征部分，写成
\begin{equation}
 \boxed{F^{\mu\nu}\in(1,0)\oplus(0,1).}
 \label{eq:twoform-chiral-decomposition-r68}
\end{equation}
这就是电磁场强可以进一步组织成自对偶与反自对偶组合的表示论来源。

\begin{figure}[htbp]
 \centering
 \resizebox{0.93\textwidth}{!}{\begin{tikzpicture}[x=1cm,y=1cm,>=Stealth,every node/.style={font=\small}]
% Block-logic grid: one row per layer, orthogonal arrows only.
%   row 1 (top) : SO -- SL            (covering map)
%   row 2       : L , R               (chiral reps of SL(2,C))
%   row 3       : build               (combine L and R)
%   row 4 (bot) : D , V , F           (composite reps)
% Same-layer boxes share one minimum width; arrows are horizontal/vertical or
% L-shaped via -| .  All labels use blocklabel (= line-label, white backing).

% ---- row 1: covering map ---------------------------------------------
\node[blockbox,minimum width=48mm] (SO) at (-7.0,3.0)
  {$SO^+(1,3)$\\proper orthochronous\\Lorentz group};
\node[blockbox,minimum width=48mm] (SL) at (0,3.0)
  {$SL(2,\mathbb C)$};
\draw[layerarrow] (SL.west) -- node[blocklabel,above=4pt]{2:1 covering map} (SO.east);

% ---- row 2: chiral representations -----------------------------------
\node[blockbox,minimum width=46mm] (L) at (-3.6,0.0)
  {$(\tfrac12,0)$\\left-handed Weyl spinor $L$};
\node[blockbox,minimum width=46mm] (R) at (3.6,0.0)
  {$(0,\tfrac12)$\\right-handed Weyl spinor $R$};
% L-shaped: leave SL straight down, run out, then drop into each chiral box.
\draw[layerarrow] (SL.south) -| node[blocklabel,pos=0.80,left]{representation} (L.north);
\draw[layerarrow] (SL.south) -| node[blocklabel,pos=0.80,right]{representation} (R.north);

% ---- row 3: build node ------------------------------------------------
% Separate operation node prevents arrows from falsely suggesting conversion.
\node[blockbox,minimum width=58mm] (build) at (0,-2.2)
  {build new Lorentz representations\\from $L$ and $R$};
\draw[layerarrow] (L.south) -| (build.north west);
\draw[layerarrow] (R.south) -| (build.north east);

% ---- row 4: composite reps --------------------------------------------
\node[blockbox,minimum width=50mm] (D) at (-5.0,-4.6)
  {$L\oplus R$\\Dirac spinor};
\node[blockbox,minimum width=50mm] (V) at (0,-4.6)
  {$L\otimes R$\\four-vector $A^\mu$};
\node[blockbox,minimum width=50mm] (F) at (5.0,-4.6)
  {$\mathrm{Sym}^2L\oplus\mathrm{Sym}^2R$\\antisymmetric tensor $F^{\mu\nu}$};
\draw[layerarrow] (build.south west) -| node[blocklabel,pos=0.80,left]{$\oplus$} (D.north);
\draw[layerarrow] (build.south) -- node[blocklabel,right]{$\otimes$} (V.north);
\draw[layerarrow] (build.south east) -| node[blocklabel,pos=0.80,right]{$\mathrm{Sym}^2$} (F.north);
\end{tikzpicture}
}
 \caption{Lorentz 群、双覆盖 $SL(2,\mathbb C)$ 与常见有限维场表示之间的关系。箭头表示表示论依赖，不表示粒子转化。}
 \label{fig:lorentz-representation-map-dp8}
\end{figure}


最小的非标量有限维复 Lorentz 表示是 $(\tfrac12,0)$ 与 $(0,\tfrac12)$，分别由二分量左手和右手 Weyl 旋量承载。空间反演交换这两个表示。更重要的是，Lorentz 四矢量与 Weyl 旋量之间的交织结构允许四动量把一种手征表示映到另一种；非零 Lorentz 标量质量正是通过这种左右手耦合进入一阶局域方程。因此描述有质量自旋 $1/2$ 粒子的局域场取直和
\begin{equation}
 \boxed{
 \psi=\begin{pmatrix}\psi_L\\ \psi_R\end{pmatrix}
 \in \left(\frac12,0\right)\oplus\left(0,\frac12\right).}
 \label{eq:dirac-rep}
\end{equation}
Lorentz 变换至此规定了场的表示性质，但尚未给出动力学。动力学必须再回答：四动量如何把左、右手分量相互映射，以及为什么质量项会把两者耦合。

为此定义两组把四矢量写成 $2\times2$ 矩阵的 Weyl 交织矩阵
\begin{equation*}
 \sigma^\mu=(\id(2),\boldsymbol\sigma_{\rm P}),
 \qquad
 \bar\sigma^\mu=(\id(2),-\boldsymbol\sigma_{\rm P}),
\end{equation*}
它们满足
\begin{equation}
 \sigma^\mu\bar\sigma^\nu+\sigma^\nu\bar\sigma^\mu
 =2\eta^{\mu\nu}\id(2),
 \qquad
 \bar\sigma^\mu\sigma^\nu+\bar\sigma^\nu\sigma^\mu
 =2\eta^{\mu\nu}\id(2).
 \label{eq:sigmabar-identity}
\end{equation}
对两个手征表示，转动生成元相同而推动生成元符号相反：
\begin{equation}
 \mathbf J_L=\frac{\hbar}{2}\boldsymbol\sigma_{\rm P},
 \quad
 \mathbf K_L=-\frac{\ii\hbar}{2}\boldsymbol\sigma_{\rm P},
 \qquad
 \mathbf J_R=\frac{\hbar}{2}\boldsymbol\sigma_{\rm P},
 \quad
 \mathbf K_R=+\frac{\ii\hbar}{2}\boldsymbol\sigma_{\rm P}.
 \label{eq:weyl-generators}
\end{equation}
于是沿 $\hat{\mathbf n}$、快度为 $\zeta$ 的纯推动分别作用为
\begin{equation}
 S_L(\zeta)=\exp\!\left[-\frac{\zeta}{2}
 \hat{\mathbf n}\!\cdot\!\boldsymbol\sigma_{\rm P}\right],
 \qquad
 S_R(\zeta)=\exp\!\left[+\frac{\zeta}{2}
 \hat{\mathbf n}\!\cdot\!\boldsymbol\sigma_{\rm P}\right].
 \label{eq:weyl-boost-r7}
\end{equation}
四动量进入局域一阶方程时，$p_\mu\sigma^\mu$ 与 $p_\mu\bar\sigma^\mu$ 分别提供右手到左手、左手到右手的 Lorentz 协变交织映射。在不引入额外内部指标、系数只取 Lorentz 标量的自由局域线性一阶系统中，左右手耦合因而可写成
\begin{equation*}
 (p_\mu\sigma^\mu)\psi_R=\mu_L\psi_L,
 \qquad
 (p_\mu\bar\sigma^\mu)\psi_L=\mu_R\psi_R,
\end{equation*}
其中 $\mu_L,\mu_R$ 是 Lorentz 标量，量纲均为动量。连续作用两次交织矩阵，并使用\eqnrefs{eq:sigmabar-identity}，非零解要求
\begin{equation*}
 p^2=\mu_L\mu_R.
\end{equation*}
对质量为 $m_{\mathrm e}$ 的粒子，质量壳是 $p^2=m_{\mathrm e}^2c^2$。要求自由二次作用量厄米给出 $\mu_R=\mu_L^*$；再用左右手场之间的相对常数相位重定义，可把共同质量参数取成正实数 $m_{\mathrm e} c$。于是得到
\begin{equation}
 \boxed{
 (p_\mu\sigma^\mu)\psi_R=m_{\mathrm e} c\,\psi_L,
 \qquad
 (p_\mu\bar\sigma^\mu)\psi_L=m_{\mathrm e} c\,\psi_R.}
 \label{eq:massive-weyl-coupled-r7}
\end{equation}
这里一般协变耦合先于任何静止系选择。若随后把质量壳参数化为
\begin{equation}
 p^0=m_{\mathrm e} c\cosh\zeta,
 \qquad
 \mathbf p=m_{\mathrm e} c\,\hat{\mathbf n}\sinh\zeta,
 \label{eq:rapidity-momentum-r7}
\end{equation}
\eqnrefs{eq:weyl-boost-r7} 可作为这两条方程的显式一致性检验，而不是一般动力学的来源。$m_{\mathrm e}\neq0$ 时左右手 Weyl 表示由质量项耦合；$m_{\mathrm e}=0$ 时两式解耦。把两式合并得到四分量动量空间方程
\begin{equation}
 \begin{pmatrix}
 -m_{\mathrm e} c\,\id(2) & p_\mu\sigma^\mu\\
 p_\mu\bar\sigma^\mu & -m_{\mathrm e} c\,\id(2)
 \end{pmatrix}
 \begin{pmatrix}\psi_L\\\psi_R\end{pmatrix}=0.
 \label{eq:weyl-block-first-order-r7}
\end{equation}

这个块矩阵给出有质量自旋 $1/2$ 粒子的一阶结构。基底无关的写法需要四个作用在四分量旋量空间上的常矩阵 $\gamma^\mu$；同一个一阶算符平方后必须恢复相对论质量壳
\begin{equation}
 p_\mu p^\mu=m_{\mathrm e}^2c^2.
 \label{eq:mass-shell}
\end{equation}
这要求矩阵满足 Clifford 代数
\begin{equation}
 \boxed{\{\gamma^\mu,\gamma^\nu\}=2\eta^{\mu\nu}\id(4).}
 \label{eq:clifford-general}
\end{equation}
于是自由 Dirac 方程写为
\begin{equation}
 \boxed{
 \left(\ii\hbar\gamma^\mu\partial_\mu-m_{\mathrm e} c\right)\psi(x)=0.}
 \label{eq:general-first-order}
\end{equation}
在 Weyl 基底中，$2\times2$ 交织矩阵直接给出
\begin{equation}
 \boxed{
 \gamma^\mu_{\rm W}=
 \begin{pmatrix}
 0&\sigma^\mu\\
 \bar\sigma^\mu&0
 \end{pmatrix},}
 \label{eq:weyl-gamma}
\end{equation}
这个表示的 Clifford 关系不需要再引入另一条“质量壳开平方”路线；它就是前面 Weyl 交织恒等式的四分量写法。直接相乘有
\begin{align*}
 \gamma^\mu_{\rm W}\gamma^\nu_{\rm W}
 &=\begin{pmatrix}
 \sigma^\mu\bar\sigma^\nu&0\\
 0&\bar\sigma^\mu\sigma^\nu
 \end{pmatrix},\\
 \{\gamma^\mu_{\rm W},\gamma^\nu_{\rm W}\}
 &=\begin{pmatrix}
 \sigma^\mu\bar\sigma^\nu+\sigma^\nu\bar\sigma^\mu&0\\
 0&\bar\sigma^\mu\sigma^\nu+\bar\sigma^\nu\sigma^\mu
 \end{pmatrix}
 =2\eta^{\mu\nu}\id(4),
\end{align*}
因此\eqnrefs{eq:clifford-general} 是已经建立的两分量 Lorentz 交织结构在四分量空间中的统一记号。四分量并非任意增加的分量数。在 $3+1$ 维复 Clifford 代数中，$2\times2$ 矩阵无法同时容纳一个平方为 $+1$、三个平方为 $-1$ 且彼此反对易的非零生成元；另一方面，\eqnrefs{eq:weyl-gamma} 已显式给出四维表示。因此最小复旋量空间满足
\begin{equation}
 \boxed{\dim_{\mathbb C}\mathcal S_{\rm Dirac}=4.}
 \label{eq:dirac-minimal-spinor-dimension-dp68}
\end{equation}

Clifford 代数还必须正确实现旋量上的 Lorentz 变换。定义无量纲反对称矩阵
\begin{equation}
 \sigma^{\mu\nu}\equiv\frac{\ii}{2}[\gamma^\mu,\gamma^\nu],
 \qquad
 \Sigma^{\mu\nu}\equiv\frac12\sigma^{\mu\nu}
 =\frac{\ii}{4}[\gamma^\mu,\gamma^\nu].
 \label{eq:sigma-definition}
\end{equation}
$\hbar\Sigma^{\mu\nu}$ 就是总生成元 $M^{\mu\nu}=L^{\mu\nu}+S^{\mu\nu}$ 中作用在旋量内部指标上的 $S^{\mu\nu}$。Clifford 代数给出
\begin{align}
 [\Sigma^{\mu\nu},\gamma^\rho]
 &=\ii\left(\eta^{\nu\rho}\gamma^\mu
 -\eta^{\mu\rho}\gamma^\nu\right),
 \label{eq:sigma-gamma-comm}\\
 [\Sigma^{\mu\nu},\Sigma^{\rho\sigma}]
 &=\ii\left(
 \eta^{\nu\rho}\Sigma^{\mu\sigma}
 -\eta^{\mu\rho}\Sigma^{\nu\sigma}
 -\eta^{\nu\sigma}\Sigma^{\mu\rho}
 +\eta^{\mu\sigma}\Sigma^{\nu\rho}
 \right),
 \label{eq:spinor-lorentz-algebra}
\end{align}
所以 $\Sigma^{\mu\nu}$ 确实实现 Lorentz 李代数。有限旋量变换为
\begin{equation}
 S(\Lambda)=
 \exp\!\left[-\frac{\ii}{2}\omega_{\mu\nu}\Sigma^{\mu\nu}\right]
 =\exp\!\left[-\frac{\ii}{4}\omega_{\mu\nu}\sigma^{\mu\nu}\right],
 \label{eq:spinor-lorentz-transform}
\end{equation}
并满足
\begin{equation}
 \boxed{
 S^{-1}(\Lambda)\gamma^\rho S(\Lambda)
 =\Lambda^\rho{}_{\sigma}\gamma^\sigma.}
 \label{eq:gamma-covariance}
\end{equation}
因此 $\gamma^\mu\partial_\mu\psi$ 与 $\psi$ 具有相同的 Lorentz 变换类型，自由 Dirac 方程的协变性来自同一个表示结构。

普通的 $\psi^\dagger\psi$ 不是 Lorentz 标量，因为推动的有限维旋量矩阵并不对普通欧几里得内积幺正。定义 Dirac 共轭
\begin{equation}
 \boxed{\bar\psi\equiv\psi^\dagger\gamma^0,}
 \label{eq:dirac-adjoint}
\end{equation}
则 Dirac 共轭的变换律为
\begin{equation}
 \boxed{\bar\psi'(x')=\bar\psi(x)S^{-1}(\Lambda),}
 \label{eq:dirac-adjoint-covariance-dp68}
\end{equation}
从而 $\bar\psi\psi$ 是 Lorentz 标量、$\bar\psi\gamma^\mu\psi$ 是四矢量。

为了区分左右手分量以及构造赝标量、轴矢量双线性，定义
\begin{equation}
 \boxed{\gamma^5\equiv \ii\gamma^0\gamma^1\gamma^2\gamma^3,}
 \label{eq:gamma5-def}
\end{equation}
它满足
\begin{equation}
 \{\gamma^5,\gamma^\mu\}=0,
 \qquad
 (\gamma^5)^2=\id(4).
 \label{eq:gamma5-anticomm}
\end{equation}
因此左右手投影子为
\begin{equation}
 \boxed{
 P_R\equiv\frac{1+\gamma^5}{2},
 \qquad
 P_L\equiv\frac{1-\gamma^5}{2},}
 \label{eq:chiral-projectors}
\end{equation}
并定义 $\psi_{R,L}=P_{R,L}\psi$。由\eqnrefs{eq:gamma5-anticomm} 逐项反对易可得
\begin{equation}
 \boxed{
 P_R^2=P_R,\quad P_L^2=P_L,\quad P_RP_L=0,\quad P_R+P_L=\id(4),
 \qquad
 P_R\gamma^\mu=\gamma^\mu P_L,\quad
 P_L\gamma^\mu=\gamma^\mu P_R.}
 \label{eq:chiral-projector-algebra-dp68}
\end{equation}
前四个关系说明 $P_{R,L}$ 确实是互补投影，后两个关系说明 Dirac 动能项在左右手子空间之间交换手征性。当 $m_{\mathrm e}=0$ 时，Dirac 方程分别在两个手征子空间内闭合；非零质量项把两者重新耦合，这与
\eqnrefs{eq:massive-weyl-coupled-r7} 的表示论结论一致。

Clifford 代数生成整个 $4\times4$ 复矩阵空间。一组按 Lorentz 变换类型组织的基为
\begin{equation}
 \left\{
 \id(4),\ \gamma^\mu,\ \sigma^{\mu\nu},\
 \gamma^\mu\gamma^5,\ \gamma^5
 \right\},
 \qquad 1+4+6+4+1=16,
 \label{eq:dirac-bilinear-basis}
\end{equation}
相应双线性
\begin{equation}
 \bar\psi\psi,\qquad
 \bar\psi\gamma^\mu\psi,\qquad
 \bar\psi\sigma^{\mu\nu}\psi,\qquad
 \bar\psi\gamma^\mu\gamma^5\psi,\qquad
 \bar\psi\ii\gamma^5\psi
 \label{eq:dirac-bilinears}
\end{equation}
依次按标量、矢量、反对称张量、轴矢量与赝标量变换。这套基统一组织 QED 顶角、手征弱流和 Yukawa 相互作用中的 Dirac 双线性。

对任意四矢量或四维协变微分算符 $a_\mu$，定义
\begin{equation*}
 \boxed{\slashed a\equiv\gamma^\mu a_\mu.}
\end{equation*}
特别地，
\begin{equation*}
 \slashed\partial\equiv\gamma^\mu\partial_\mu,
 \qquad
 \slashed A\equiv\gamma^\mu A_\mu,
 \qquad
 \slashed D\equiv\gamma^\mu D_\mu,
 \qquad
 \slashed p\equiv\gamma^\mu p_\mu.
\end{equation*}
斜杠把 Lorentz 指标收缩与旋量矩阵结构写在同一个符号中，例如自由 Dirac 算符就是 $\ii\hbar\slashed\partial-m_{\mathrm e} c$。


时空表示至此确定了自由电子场的一阶动力学结构。带电电子还具有局域 $U(1)$ 相位自由度；Dirac 动能项若要在逐点相位变化下保持协变，普通导数 $\partial_\mu$ 必须由含规范联络的 $D_\mu$ 取代；这个联络就是电磁四势。

\begin{note}[一个常用显式表示]
需要写出分量时可以选择 Dirac 基底
\begin{equation*}
 \gamma^0=
 \begin{pmatrix}\id(2)&0\\0&-\id(2)\end{pmatrix},
 \qquad
 \gamma^i=
 \begin{pmatrix}0&\sigma_i\\-\sigma_i&0\end{pmatrix},
\end{equation*}
其中右端 $\sigma_i$ 是 Pauli 矩阵。三维 Pauli 矩阵 $\sigma_i$ 与反对称 Lorentz 张量 $\sigma^{\mu\nu}$ 是不同对象。
\end{note}

% -----------------------------------------------------------------------------
% Dirac/Clifford identities, proved after gamma matrices are constructed
% -----------------------------------------------------------------------------

Poincar\'e 表示与 Clifford 代数确定 $\gamma^\mu$ 以后，后续散射与圈图所需的 Dirac 矩阵代数不再是独立公式表，而都来自
$\{\gamma^\mu,\gamma^\nu\}=2\eta^{\mu\nu}\id(4)$、矩阵迹的循环性以及 $\gamma^5$ 的定义。本段只在正文列出会被反复引用的目标恒等式；逐次反对易和指标缩并统一放在本结构末尾的推导区。

最简单的 Lorentz 二阶张量迹只能正比于 $\eta^{\mu\nu}$。取 $\mu=\nu=0$ 并使用 $(\gamma^0)^2=\id(4)$ 固定系数后，得到
\begin{equation}
 \boxed{\Tr(\gamma^\mu\gamma^\nu)=4\eta^{\mu\nu}.}
 \label{eq:trace-two-gamma-r9}
\end{equation}
四个伽马矩阵的迹可把最左因子依次反对易到最右端，再利用\eqnrefs{eq:trace-two-gamma-r9} 关闭递推，结果为
\begin{equation}
 \boxed{
 \Tr(\gamma^\mu\gamma^\nu\gamma^\rho\gamma^\sigma)
 =4(\eta^{\mu\nu}\eta^{\rho\sigma}
 -\eta^{\mu\rho}\eta^{\nu\sigma}
 +\eta^{\mu\sigma}\eta^{\nu\rho}).}
 \label{eq:trace-four-gamma-r9}
\end{equation}
这两条迹恒等式是未偏振费米子散射把自旋求和转成 Lorentz 标量的基本工具。


当费米子链继续增长时，不需要为六个、八个乃至更多 $\gamma$ 矩阵分别记忆新公式。把最左端矩阵逐次反对易到最右端，再用迹的循环性关闭链条，可得到统一递推
\begin{equation}
 \boxed{
 T_{2n}^{\mu_1\cdots\mu_{2n}}
 =\sum_{j=2}^{2n}(-1)^j\eta^{\mu_1\mu_j}
 \Tr\!\left(\gamma^{\mu_2}\cdots\widehat{\gamma^{\mu_j}}\cdots\gamma^{\mu_{2n}}\right),}
 \qquad
 T_{2n}^{\mu_1\cdots\mu_{2n}}
 \equiv\Tr(\gamma^{\mu_1}\cdots\gamma^{\mu_{2n}}).
 \label{eq:trace-even-recursion-r68}
\end{equation}
这里帽号表示删除相应因子。另一方面，在严格四维且迹中不含额外 $\gamma^5$ 时，$\gamma^5$ 与每个普通 $\gamma^\mu$ 反对易，因此任意奇数长度的普通 Dirac 链满足
\begin{equation}
 \boxed{\Tr(\gamma^{\mu_1}\cdots\gamma^{\mu_{2n+1}})=0.}
 \label{eq:trace-odd-zero-r68}
\end{equation}
这两条式子分别负责把长偶数链递归缩短，以及在自旋求和中立即消去奇数链。涉及维数正规化中的真正手征迹时，$\gamma^5$ 的处理需按所采用的正规化方案另行定义。

对任意四矢量 $a^\mu,b^\mu,c^\mu$，定义斜线记号 $\slashed a\equiv\gamma^\mu a_\mu$。把外侧的 $\gamma^\mu$ 逐次穿过内部矩阵串，可得到长度为一、二、三的三条常用夹乘恒等式：
\begin{align}
 \boxed{\gamma^\mu\slashed a\gamma_\mu=-2\slashed a,}
 \label{eq:gamma-one-sandwich-r10}\\[2pt]
 \boxed{\gamma^\mu\slashed a\slashed b\gamma_\mu
 =4a\!\cdot b\,\id(4),}
 \label{eq:gamma-two-sandwich-r10}\\[2pt]
 \boxed{\gamma^\mu\slashed a\slashed b\slashed c\gamma_\mu
 =-2\slashed c\slashed b\slashed a.}
 \label{eq:gamma-three-sandwich-r10}
\end{align}
这些式子的输入只有四维 Clifford 代数；其中的数值系数 $-2,4,-2$ 会在维数正规化推广到 $d$ 维时发生相应改变，所以圈图计算若尚未取 $d\to4$，不能无条件使用四维版本。


在维数正规化中应保留 $d$ 维 Clifford 代数本身。由 $g^\mu{}_{\mu}=d$ 逐项收缩可得最常用的一阶夹乘
\begin{equation}
 \boxed{\gamma^\mu\gamma^\alpha\gamma_\mu=(2-d)\gamma^\alpha,}
 \label{eq:gamma-d-sandwich-r68}
\end{equation}
它在 $d=4$ 时才退化为 $-2\gamma^\alpha$。因此圈积分分子中的 $O(\epsilon)$ 代数不能在积分完成前被四维恒等式提前抹掉。

严格四维中，含一个 $\gamma^5$ 的最低阶非零迹由当前取向约定唯一固定为
\begin{equation}
 \boxed{\Tr(\gamma^5\gamma^\mu\gamma^\nu\gamma^\rho\gamma^\sigma)
 =-4\ii\epsilon^{\mu\nu\rho\sigma}.}
 \label{eq:gamma5-four-trace-r68}
\end{equation}
它把手征 Dirac 代数与 Levi--Civita 张量连接起来，后续轴矢量流与异常计算都从这一约定出发。

最后，复 $4\times4$ 矩阵空间有 $16$ 个复维度。由
$\{\id(4),\gamma^\mu,\sigma^{\mu\nu},\gamma^\mu\gamma^5,\gamma^5\}$
构造一组迹正交基 $\Gamma_A$，并以 $\Gamma^A$ 表示其对偶基，则矩阵指标的完备关系为
\begin{equation}
 \boxed{
 \frac14\sum_A
 (\Gamma_A)_{ij}(\Gamma^A)_{kl}
 =\delta_{il}\delta_{kj}.}
 \label{eq:dirac-matrix-completeness-r10}
\end{equation}
Fierz 重排只是把这条 $16$ 维矩阵完备关系插入两组旋量指标之间，再改变指标配对；因此 Fierz 重排完全由有限维 Dirac 矩阵空间的完备性与指标重配得到。


例如，对普通复数值外线旋量，标量双线性的一种常用重排为
\begin{align}
 (\bar u_1u_2)(\bar u_3u_4)
 =\frac14\Big[&
 (\bar u_1u_4)(\bar u_3u_2)
 +(\bar u_1\gamma^\mu u_4)(\bar u_3\gamma_\mu u_2)
 \nonumber\\
 &+\frac12(\bar u_1\sigma^{\mu\nu}u_4)(\bar u_3\sigma_{\mu\nu}u_2)
 -(\bar u_1\gamma^\mu\gamma^5u_4)(\bar u_3\gamma_\mu\gamma^5u_2)
 \nonumber\\
 &+(\bar u_1\gamma^5u_4)(\bar u_3\gamma^5u_2)\Big].
 \label{eq:fierz-scalar-r68}
\end{align}
若 $u_i$ 换成 Grassmann 场，还必须另外计入为改变场次序而产生的统计负号；矩阵完备系数与费米统计符号是两件不同的事。

同一组迹恒等式还给出 Lorentz 生成元矩阵的两个常用收缩：
\begin{equation}
 \boxed{\Tr(\sigma^{\mu\nu}\sigma^{\rho\sigma})
 =4(\eta^{\mu\rho}\eta^{\nu\sigma}-\eta^{\mu\sigma}\eta^{\nu\rho}),
 \qquad
 \sigma^{\mu\nu}\sigma_{\mu\nu}=12\id(4).}
 \label{eq:sigma-sigma-trace-r68}
\end{equation}
在当前 $\eta=\mathrm{diag}(1,-1,-1,-1)$、$\epsilon^{0123}=+1$ 约定下，$\gamma^5$ 把反对称张量映射到其 Hodge 对偶：
\begin{equation}
 \boxed{\sigma^{\mu\nu}\gamma^5
 =\frac{\ii}{2}\epsilon^{\mu\nu\rho\sigma}\sigma_{\rho\sigma}.}
 \label{eq:sigma-gamma5-dual-r68}
\end{equation}
这些结果分别用于四费米子算符换基、Pauli 形状因子以及手征张量结构。

% DeepPass21: detailed Pauli/Weyl identities at the end of the owner subsection.

\begin{equation}
\boxed{
\det(p_\mu\sigma^\mu)=p_\mu p^\mu,
\qquad
\det(p_\mu\bar\sigma^\mu)=p_\mu p^\mu.}
\label{eq:weyl-determinant-mass-shell-dp21}
\end{equation}





\begin{derivation}[从\eqnrefs{eq:pauli-product-proof-r10}到\eqnrefs{eq:epsilon-contractions-r10}]

Pauli 乘法\eqnrefs{eq:pauli-product-proof-r10} 的反对称部分由 $\epsilon_{ijk}$ 编码；两个这样的结构相乘时，需要把 Levi--Civita 的双重缩并化成 Kronecker $\delta$。这里从三维叉积的分量定义出发，补出正文\eqnrefs{eq:epsilon-contractions-r10} 的指标恒等式。



取标准正交基 $\mathbf e_i$，由
\begin{equation*}
 \epsilon_{ijk}=\mathbf e_i\cdot(\mathbf e_j\times\mathbf e_k)
\end{equation*}
以及完备关系 $\sum_i(\mathbf e_i\cdot\mathbf a)(\mathbf e_i\cdot\mathbf b)=\mathbf a\cdot\mathbf b$，有
\begin{align*}
 \epsilon_{ijk}\epsilon_{ilm}
 &= (\mathbf e_j\times\mathbf e_k)\cdot(\mathbf e_l\times\mathbf e_m)\\
 &= (\mathbf e_j\cdot\mathbf e_l)(\mathbf e_k\cdot\mathbf e_m)
 -(\mathbf e_j\cdot\mathbf e_m)(\mathbf e_k\cdot\mathbf e_l)\\
 &=\delta_{jl}\delta_{km}-\delta_{jm}\delta_{kl}.
\end{align*}
再令 $l=j$ 并对 $j$ 求和，
\begin{align*}
 \epsilon_{ijk}\epsilon_{ijm}
 &=\delta_{jj}\delta_{km}-\delta_{jm}\delta_{kj}\\
 &=3\delta_{km}-\delta_{km}=2\delta_{km},
\end{align*}
得到\eqnrefs{eq:epsilon-contractions-r10}。
\end{derivation}

\begin{derivation}[从\eqnrefs{eq:sigma-barsigma-foundation-r10}到\eqnrefs{eq:sigmabar-proof-r10}]

正文\eqnrefs{eq:sigma-barsigma-foundation-r10} 把四矢量嵌入两组 $2\times2$ Weyl 矩阵。为了使左右手一阶算符的乘积恢复 Minkowski 度规，逐分量检查时间--时间、时间--空间和空间--空间三类组合，便得到正文\eqnrefs{eq:sigmabar-proof-r10}。



由\eqnrefs{eq:sigma-barsigma-foundation-r10} 分三类指标检查。若 $\mu=\nu=0$，
\begin{equation*}
 \sigma^0\bar\sigma^0+\sigma^0\bar\sigma^0=2\id_2.
\end{equation*}
若一个指标为时间、另一个为空间，
\begin{equation*}
 \sigma^0\bar\sigma^i+\sigma^i\bar\sigma^0
 =-\sigma_i+\sigma_i=0.
\end{equation*}
若两个指标都为空间，\eqnrefs{eq:pauli-comm-anticomm-r10} 给
\begin{align*}
 \sigma^i\bar\sigma^j+\sigma^j\bar\sigma^i
 &=-(\sigma_i\sigma_j+\sigma_j\sigma_i)\\
 &=-2\delta_{ij}\id_2.
\end{align*}
三种情况分别对应 $\eta^{00}=1$、$\eta^{0i}=0$、$\eta^{ij}=-\delta_{ij}$，故合并为\eqnrefs{eq:sigmabar-proof-r10}。
\end{derivation}

\begin{derivation}[从\eqnrefs{eq:sl2c-hermitian-map-dp8}到\eqnrefs{eq:sl2c-double-cover-dp8}]
先计算

\begin{align*}
 \det X
 &=(x^0+x^3)(x^0-x^3)
 -(x^1-\ii x^2)(x^1+\ii x^2)\\
 &=(x^0)^2-(x^1)^2-(x^2)^2-(x^3)^2\\
 &=x_\mu x^\mu.
\end{align*}
对 $A\in SL(2,\mathbb C)$，$X'=AXA^\dagger$ 仍为厄米，且
\begin{align*}
 \det X'
 &=\det A\,\det X\,\det A^\dagger\\
 &=|\det A|^2\det X\\
 &=\det X.
\end{align*}
因此存在唯一实线性变换 $\Phi(A)$ 使 $X(x')=AX(x)A^\dagger$，并且 $\Phi(A)$ 保持 Minkowski 度规。矩阵乘法还给出
\begin{equation*}
 \Phi(A_2A_1)X
 =A_2A_1X A_1^\dagger A_2^\dagger
 =\Phi(A_2)\Phi(A_1)X,
\end{equation*}
所以 $\Phi$ 是群同态。由于 $SL(2,\mathbb C)$ 连通且 $\Phi(\id_2)=\id_4$，其像落在 Lorentz 群的 proper orthochronous 连通分支 $SO^+(1,3)$。

严格求核。若 $A\in\ker\Phi$，则对任意厄米 $X$ 都有
\begin{equation*}
 AXA^\dagger=X.
\end{equation*}
先取 $X=\id_2$，得到 $AA^\dagger=\id_2$，所以 $A$ 是幺正矩阵，$A^\dagger=A^{-1}$。于是上式改写成
\begin{equation*}
 AXA^{-1}=X
 \qquad\text{对任意厄米 }X.
\end{equation*}
厄米矩阵张成整个 $M_2(\mathbb C)$ 复矩阵空间，因此 $A$ 与所有 $2\times2$ 复矩阵对易，只能是标量矩阵 $A=\lambda\id_2$。再用 $\det A=\lambda^2=1$，得到
\begin{equation*}
 \ker\Phi=\{+\id_2,-\id_2\}.
\end{equation*}

最后证明满射，而不是只验证若干特殊变换。把 $\mathfrak{sl}(2,\mathbb C)$ 看成六维实 Lie 代数，其一组实基可取
\begin{equation*}
 -\frac{\ii}{2}\sigma_i,
 \qquad
 +\frac12\sigma_i,
 \qquad i=1,2,3.
\end{equation*}
对转动型无穷小元
\begin{equation*}
 A=\id_2-\frac{\ii}{2}\boldsymbol\theta\cdot\boldsymbol\sigma_{\rm P}
 +O(\theta^2),
\end{equation*}
有
\begin{align*}
 \delta X
 &=-\frac{\ii}{2}
 [\boldsymbol\theta\cdot\boldsymbol\sigma_{\rm P},
   \mathbf x\cdot\boldsymbol\sigma_{\rm P}]\\
 &=(\boldsymbol\theta\times\mathbf x)\cdot\boldsymbol\sigma_{\rm P},
\end{align*}
即 $\delta x^0=0$、$\delta\mathbf x=\boldsymbol\theta\times\mathbf x$，给出三个独立空间转动方向。对推动型无穷小元
\begin{equation*}
 A=\id_2+\frac12\boldsymbol\xi\cdot\boldsymbol\sigma_{\rm P}
 +O(\xi^2),
\end{equation*}
则
\begin{align*}
 \delta X
 &=\frac12\{\boldsymbol\xi\cdot\boldsymbol\sigma_{\rm P},X\}\\
 &=(\boldsymbol\xi\cdot\mathbf x)\id_2
 +x^0\boldsymbol\xi\cdot\boldsymbol\sigma_{\rm P},
\end{align*}
因此 $\delta x^0=\boldsymbol\xi\cdot\mathbf x$、$\delta\mathbf x=x^0\boldsymbol\xi$，给出三个独立 Lorentz 推动方向。故微分
\begin{equation*}
 \dd\Phi_e:\mathfrak{sl}(2,\mathbb C)_{\mathbb R}
 \longrightarrow\mathfrak{so}(1,3)
\end{equation*}
在六个实维度上满秩，因而是 Lie 代数同构。于是 $\Phi(SL(2,\mathbb C))$ 是 $SO^+(1,3)$ 的开子群；而 $SO^+(1,3)$ 连通，开子群的不同陪集若多于一个会把群分成互不相交的非空开集，所以只能有一个陪集。于是 $\Phi$ 满射。

综上，$\Phi$ 是满射群同态，核为 $\{\pm\id_2\}$，由第一同构定理得到
\begin{equation*}
 SL(2,\mathbb C)/\{\pm\id_2\}
 \cong SO^+(1,3),
\end{equation*}
即正文\eqnrefs{eq:sl2c-double-cover-dp8}。这里的“双覆盖”因此同时包含“恰有二元核”和“确实覆盖整个 $SO^+(1,3)$”两层结论。
\end{derivation}

\begin{derivation}[从\eqnrefs{eq:JK-algebra}到\eqnrefs{eq:lorentz-su2-algebra-dp34}]

正文\eqnrefs{eq:JK-algebra} 已把六个 Lorentz 生成元写成三维转动 $J_i$ 与推动 $K_i$。将它们组合成 $\mathsf J_i^{(\pm)}=(J_i\pm\ii K_i)/2$ 后，逐项代入三组 $J$--$K$ 对易关系即可证明两组生成元各自闭合且彼此对易，从而得到正文\eqnrefs{eq:lorentz-su2-algebra-dp34}。



由
\begin{equation*}
 \mathsf J_i^{(+)}=\frac12(J_i+\ii K_i),
 \qquad
 \mathsf J_i^{(-)}=\frac12(J_i-\ii K_i)
\end{equation*}
与
\begin{equation*}
 [J_i,J_j]=\ii\hbar\epsilon_{ijk}J_k,
 \qquad
 [J_i,K_j]=\ii\hbar\epsilon_{ijk}K_k,
 \qquad
 [K_i,K_j]=-\ii\hbar\epsilon_{ijk}J_k,
\end{equation*}
先计算第一组：
\begin{align*}
 [\mathsf J_i^{(+)},\mathsf J_j^{(+)}]
 &=\frac14\Bigl(
 [J_i,J_j]+\ii[J_i,K_j]
 +\ii[K_i,J_j]-[K_i,K_j]\Bigr).
\end{align*}
利用
\begin{equation*}
 [K_i,J_j]=-[J_j,K_i]
 =\ii\hbar\epsilon_{ijk}K_k,
\end{equation*}
得到
\begin{align*}
 [\mathsf J_i^{(+)},\mathsf J_j^{(+)}]
 &=\frac14\left(
 2\ii\hbar\epsilon_{ijk}J_k
 -2\hbar\epsilon_{ijk}K_k\right),\\
 &=\ii\hbar\epsilon_{ijk}
 \frac{J_k+\ii K_k}{2},\\
 &=\ii\hbar\epsilon_{ijk}\mathsf J_k^{(+)}.
\end{align*}
第二组从定义 $\mathsf J_i^{(-)}=(J_i-\ii K_i)/2$ 独立计算：
\begin{align*}
 [\mathsf J_i^{(-)},\mathsf J_j^{(-)}]
 &=\frac14\Bigl(
 [J_i,J_j]-\ii[J_i,K_j]
 -\ii[K_i,J_j]-[K_i,K_j]\Bigr),\\
 &=\frac14\Bigl(
 \ii\hbar\epsilon_{ijk}J_k
 +\hbar\epsilon_{ijk}K_k
 +\hbar\epsilon_{ijk}K_k
 +\ii\hbar\epsilon_{ijk}J_k\Bigr),\\
 &=\frac12\Bigl(
 \ii\hbar\epsilon_{ijk}J_k
 +\hbar\epsilon_{ijk}K_k\Bigr),\\
 &=\ii\hbar\epsilon_{ijk}\frac{J_k-\ii K_k}{2},\\
 &=\ii\hbar\epsilon_{ijk}\mathsf J_k^{(-)}.
\end{align*}
交叉对易子为
\begin{align*}
 [\mathsf J_i^{(+)},\mathsf J_j^{(-)}]
 &=\frac14\Bigl(
 [J_i,J_j]-\ii[J_i,K_j]
 +\ii[K_i,J_j]+[K_i,K_j]\Bigr),\\
 &=0.
\end{align*}
因此复化 Lorentz 李代数确实分解为两份彼此对易的 $\mathfrak{su}(2)$。
\end{derivation}

\begin{derivation}[从\eqnrefs{eq:sigmabar-identity}到\eqnrefs{eq:massive-weyl-coupled-r7}]

正文已指出 $p_\mu\sigma^\mu$ 与 $p_\mu\bar\sigma^\mu$ 是两个手征表示之间的一阶 Lorentz 交织映射。这里只需确定允许的标量系数，并由质量壳固定它们的乘积；静止系不参与一般式的建立。

先写最一般线性耦合
\begin{equation*}
 (p_\mu\sigma^\mu)\psi_R=\mu_L\psi_L,
 \qquad
 (p_\mu\bar\sigma^\mu)\psi_L=\mu_R\psi_R.
\end{equation*}
由正文\eqnrefs{eq:sigmabar-identity}，
\begin{align*}
 (p_\mu\bar\sigma^\mu)(p_\nu\sigma^\nu)
 &=\frac12p_\mu p_\nu
 \left(\bar\sigma^\mu\sigma^\nu+\bar\sigma^\nu\sigma^\mu\right)\\
 &=p_\mu p^\mu\,\id_2.
\end{align*}
把第一式左乘 $p_\mu\bar\sigma^\mu$，再代入第二式：
\begin{align*}
 p^2\psi_R
 &=(p\!\cdot\!\bar\sigma)(p\!\cdot\!\sigma)\psi_R\\
 &=\mu_L(p\!\cdot\!\bar\sigma)\psi_L\\
 &=\mu_L\mu_R\psi_R.
\end{align*}
对非零态得到
\begin{equation*}
 p^2=\mu_L\mu_R.
\end{equation*}
同样作用在 $\psi_L$ 上得到相同条件。对质量壳 $p^2=\me^2c^2$，
\begin{equation*}
 \mu_L\mu_R=\me^2c^2.
\end{equation*}

若自由二次作用量厄米，左右两条方程中的质量系数互为复共轭，
\begin{equation*}
 \mu_R=\mu_L^*.
\end{equation*}
于是 $|\mu_L|=|\mu_R|=\me c$。对两个手征场作相反的常数重相位可消去这一个共同复相位，因此可选
\begin{equation*}
 \mu_L=\mu_R=\me c>0.
\end{equation*}
这就得到正文\eqnrefs{eq:massive-weyl-coupled-r7}。

作为随后才做的检查，把质量壳写成正文\eqnrefs{eq:rapidity-momentum-r7}。由 Pauli 代数，
\begin{align*}
 p_\mu\sigma^\mu
 &=\me c\left(\cosh\zeta\,\id_2
 -\sinh\zeta\,\hat{\mathbf n}\cdot\boldsymbol\sigma_{\rm P}\right)\\
 &=\me c\,\ee^{-\zeta\hat{\mathbf n}\cdot\boldsymbol\sigma_{\rm P}},\\
 p_\mu\bar\sigma^\mu
 &=\me c\,\ee^{+\zeta\hat{\mathbf n}\cdot\boldsymbol\sigma_{\rm P}}.
\end{align*}
与正文\eqnrefs{eq:weyl-boost-r7} 的左右手有限推动一致。快度表示在这里提供一般耦合式的一种显式质量壳参数化检验；一般方程本身已经由 Lorentz 交织映射与质量壳条件确定。
\end{derivation}



\begin{derivation}[从\eqnrefs{eq:clifford-general}到\eqnrefs{eq:dirac-minimal-spinor-dimension-dp68}]
设 $A,B$ 是任意两个不同的 Clifford 生成元，因此

\begin{equation*}
 AB=-BA,
 \qquad
 \det A\neq0,
 \qquad
 \det B\neq0.
\end{equation*}
若表示维数为奇数 $n$，则
\begin{align*}
 \det(AB)
 &=\det(-BA)\\
 &=(-1)^n\det(BA)\\
 &=-\det(AB),
\end{align*}
与 $\det(AB)\neq0$ 矛盾。因此任何有限维复 Clifford 表示的维数在这里都不能为奇数；特别地，三维表示已经被排除。

还需排除唯一更小的偶数候选 $n=2$。任意 $2\times2$ 复矩阵可唯一写成
\begin{equation*}
 X=a\id_2+\mathbf b\cdot\boldsymbol\sigma_{\rm P}.
\end{equation*}
若三个空间 Clifford 生成元存在于二维空间，经相似变换与非零常数重标度后可取为三条 Pauli 方向。与三者都反对易的时间生成元 $X$ 必须满足
\begin{equation*}
 0=\{X,\sigma_i\}
 =2a\sigma_i+2b_i\id_2
 \qquad(i=1,2,3).
\end{equation*}
取迹给 $b_i=0$；再乘任意 $\sigma_j$ 后取迹给 $a=0$。所以唯一同时与三个 Pauli 方向反对易的二维矩阵是零矩阵，这与 $(\gamma^0)^2=\id_2$ 矛盾。

因此非零复表示的维数既不能为奇数，也不能为二，最小可能值至少为四。正文\eqnrefs{eq:weyl-gamma} 又给出显式 $4\times4$ 矩阵并满足\eqnrefs{eq:clifford-general}，所以
\begin{equation*}
 \dim_{\mathbb C}S_{\rm Dirac}^{\rm min}=4.
\end{equation*}
这才完成“最小维数为四”的存在性与不可更小性两部分证明。
\end{derivation}



\begin{derivation}[从\eqnrefs{eq:sigma-definition}到\eqnrefs{eq:spinor-lorentz-algebra}]
先计算
\begin{align*}
 [\gamma^\mu\gamma^\nu,\gamma^\rho]
 &=\gamma^\mu\gamma^\nu\gamma^\rho
 -\gamma^\rho\gamma^\mu\gamma^\nu\\
 &=\gamma^\mu(2\eta^{\nu\rho}-\gamma^\rho\gamma^\nu)
 -(2\eta^{\mu\rho}-\gamma^\mu\gamma^\rho)\gamma^\nu\\
 &=2\eta^{\nu\rho}\gamma^\mu-2\eta^{\mu\rho}\gamma^\nu.
\end{align*}
交换 $\mu\leftrightarrow\nu$ 后相减，再乘正文\eqnrefs{eq:sigma-definition} 中的 $\ii/4$，便得到
\begin{equation*}
 [\Sigma^{\mu\nu},\gamma^\rho]
 =\ii(\eta^{\nu\rho}\gamma^\mu-\eta^{\mu\rho}\gamma^\nu),
\end{equation*}
即正文\eqnrefs{eq:sigma-gamma-comm}。

由 $\Sigma^{\rho\sigma}=\frac{\ii}{4}[\gamma^\rho,\gamma^\sigma]$，有
\begin{align*}
 [\Sigma^{\mu\nu},\Sigma^{\rho\sigma}]
 &=\frac{\ii}{4}\Bigl(
 [\Sigma^{\mu\nu},\gamma^\rho]\gamma^\sigma
 +\gamma^\rho[\Sigma^{\mu\nu},\gamma^\sigma]\\
 &\qquad-[\Sigma^{\mu\nu},\gamma^\sigma]\gamma^\rho
 -\gamma^\sigma[\Sigma^{\mu\nu},\gamma^\rho]\Bigr).
\end{align*}
代入正文\eqnrefs{eq:sigma-gamma-comm}，再把出现的 $\gamma^\alpha\gamma^\beta-\gamma^\beta\gamma^\alpha$ 重新写回 $\Sigma^{\alpha\beta}$，得到
\begin{equation*}
 [\Sigma^{\mu\nu},\Sigma^{\rho\sigma}]
 =\ii\left(
 \eta^{\nu\rho}\Sigma^{\mu\sigma}
 -\eta^{\mu\rho}\Sigma^{\nu\sigma}
 -\eta^{\nu\sigma}\Sigma^{\mu\rho}
 +\eta^{\mu\sigma}\Sigma^{\nu\rho}
 \right),
\end{equation*}
即正文\eqnrefs{eq:spinor-lorentz-algebra}。
\end{derivation}

\begin{derivation}[从\eqnrefs{eq:spinor-lorentz-transform}到\eqnrefs{eq:gamma-covariance}]

正文\eqnrefs{eq:spinor-lorentz-transform} 给出有限旋量变换的指数形式。其无穷小相似变换由 $[\Sigma^{\mu\nu},\gamma^\rho]$ 完全决定；正文\eqnrefs{eq:sigma-gamma-comm} 恰好把这一对易子写成四矢量表示的生成元作用，因此两种表示具有同一 Lie 代数微分，指数化后得到 $\gamma^\rho$ 的有限协变变换律。

写成一阶参数展开，
\begin{equation*}
 S=\id-\frac{\ii}{2}\omega_{\mu\nu}\Sigma^{\mu\nu}+O(\omega^2),
 \qquad
 S^{-1}=\id+\frac{\ii}{2}\omega_{\mu\nu}\Sigma^{\mu\nu}+O(\omega^2).
\end{equation*}
于是
\begin{align*}
 S^{-1}\gamma^\rho S
 &=\gamma^\rho+\frac{\ii}{2}\omega_{\mu\nu}
 [\Sigma^{\mu\nu},\gamma^\rho]+O(\omega^2)\\
 &=\gamma^\rho+\omega^\rho{}_{\sigma}\gamma^\sigma+O(\omega^2).
\end{align*}
四矢量表示同样满足
\begin{equation*}
 \Lambda^\rho{}_{\sigma}
 =\delta^\rho{}_{\sigma}+\omega^\rho{}_{\sigma}+O(\omega^2).
\end{equation*}
两种表示由同一 Lie 代数参数生成，因此沿群参数连续积分后得到
\begin{equation*}
 S^{-1}(\Lambda)\gamma^\rho S(\Lambda)
 =\Lambda^\rho{}_{\sigma}\gamma^\sigma,
\end{equation*}
即正文\eqnrefs{eq:gamma-covariance}。
\end{derivation}

\begin{derivation}[从\eqnrefs{eq:dirac-adjoint}到\eqnrefs{eq:dirac-adjoint-covariance-dp68}]

普通内积 $\psi^\dagger\psi$ 在推动下并不保持不变，因此正文\eqnrefs{eq:dirac-adjoint} 引入 $\gamma^0$ 来构造适合 Lorentz 群的双线性。要验证这个定义确实有效，只需证明 $S^\dagger\gamma^0=\gamma^0S^{-1}$。

在所采用的度规约定下，
\begin{equation*}
 (\gamma^0)^\dagger=\gamma^0,
 \qquad
 (\gamma^i)^\dagger=-\gamma^i,
\end{equation*}
等价于
\begin{equation*}
 (\gamma^\mu)^\dagger=\gamma^0\gamma^\mu\gamma^0.
\end{equation*}
因此
\begin{equation*}
 (\Sigma^{\mu\nu})^\dagger=\gamma^0\Sigma^{\mu\nu}\gamma^0.
\end{equation*}
对正文\eqnrefs{eq:spinor-lorentz-transform} 取厄米共轭便有
\begin{align*}
 S^\dagger
 &=\exp\!\left[+\frac{\ii}{2}\omega_{\mu\nu}(\Sigma^{\mu\nu})^\dagger\right]\\
 &=\gamma^0\exp\!\left[+\frac{\ii}{2}\omega_{\mu\nu}\Sigma^{\mu\nu}\right]\gamma^0
 =\gamma^0S^{-1}\gamma^0.
\end{align*}
于是
\begin{align*}
 \bar\psi'
 &=(S\psi)^\dagger\gamma^0
 =\psi^\dagger S^\dagger\gamma^0
 =\psi^\dagger\gamma^0S^{-1}
 =\bar\psi S^{-1},
\end{align*}
这就是正文\eqnrefs{eq:dirac-adjoint-covariance-dp68}。
\end{derivation}

\begin{derivation}[从\eqnrefs{eq:gamma5-def}到\eqnrefs{eq:chiral-projector-algebra-dp68}]
以 $\mu=0$ 为例，
\begin{align*}
 \gamma^5\gamma^0
 &=\ii\gamma^0\gamma^1\gamma^2\gamma^3\gamma^0\\
 &=-\ii\gamma^0\gamma^0\gamma^1\gamma^2\gamma^3
 =-\gamma^0\gamma^5.
\end{align*}
同理对任意 $\mu$ 都有 $\{\gamma^5,\gamma^\mu\}=0$。平方时，右半组四个 $\gamma$ 移到左半组同名矩阵旁共需 $3+2+1=6$ 次交换：
\begin{align*}
 (\gamma^5)^2
 &=-\gamma^0\gamma^1\gamma^2\gamma^3
 \gamma^0\gamma^1\gamma^2\gamma^3\\
 &=-(-1)^6(\gamma^0)^2(\gamma^1)^2(\gamma^2)^2(\gamma^3)^2
 =\id_4.
\end{align*}
两条关系合起来就是正文\eqnrefs{eq:gamma5-anticomm}。

先计算
\begin{align*}
 P_R^2&=\frac14(1+2\gamma^5+(\gamma^5)^2)=P_R,\\
 P_L^2&=\frac14(1-2\gamma^5+(\gamma^5)^2)=P_L,\\
 P_RP_L&=\frac14(1-(\gamma^5)^2)=0,\\
 P_R+P_L&=\id_4.
\end{align*}
再用 $\gamma^5\gamma^\mu=-\gamma^\mu\gamma^5$，
\begin{align*}
 P_R\gamma^\mu
 &=\frac12(1+\gamma^5)\gamma^\mu
 =\frac12\gamma^\mu(1-\gamma^5)
 =\gamma^\mu P_L,\\
 P_L\gamma^\mu&=\gamma^\mu P_R.
\end{align*}
这就得到正文\eqnrefs{eq:chiral-projector-algebra-dp68}。
\end{derivation}

\begin{derivation}[从\eqnrefs{eq:lorentz-irrep-dictionary-dp8}到\eqnrefs{eq:weyl-vector-product-r68}]
只写维数 $2\times2=4$ 不能证明一个双旋量就是四矢量；关键是它的变换律必须与正文已经建立的
$X(x)\mapsto AX(x)A^\dagger$ 完全相同。设左手 Weyl 旋量按
\begin{equation*}
 \xi_\alpha\longmapsto \xi'_\alpha
 =A_\alpha{}^\beta\xi_\beta,
 \qquad A\in SL(2,\mathbb C),
\end{equation*}
变换。其厄米共轭带点指标按
\begin{equation*}
 \bar\chi_{\dot\alpha}\longmapsto
 \bar\chi'_{\dot\alpha}
 =(A^\dagger)_{\dot\alpha}{}^{\dot\beta}
 \bar\chi_{\dot\beta}
\end{equation*}
变换。构造一个左右各带一个指标的对象
\begin{equation*}
 V_{\alpha\dot\beta}=\xi_\alpha\bar\chi_{\dot\beta}.
\end{equation*}
把 $V_{\alpha\dot\beta}$ 看成 $2\times2$ 矩阵，则
\begin{align*}
 V'_{\alpha\dot\beta}
 &=A_\alpha{}^\gamma
 V_{\gamma\dot\delta}
 (A^\dagger)^{\dot\delta}{}_{\dot\beta},\\
 \Longleftrightarrow\qquad
 V'&=AVA^\dagger.
\end{align*}
这已经与正文\eqnrefs{eq:sl2c-action-X-dp8} 的四矢量矩阵变换完全相同。

为了把矩阵分量显式还原成四矢量，使用
$\sigma^\mu=(I_2,\boldsymbol\sigma_{\rm P})$、
$\bar\sigma^\mu=(I_2,-\boldsymbol\sigma_{\rm P})$。Weyl 交织恒等式给
\begin{equation*}
 \operatorname{Tr}(\bar\sigma^\mu\sigma^\nu)=2\eta^{\mu\nu}.
\end{equation*}
因此展开
\begin{equation*}
 V=V_\nu\sigma^\nu
\end{equation*}
的逆变换不是一个符号约定，而由迹正交性唯一固定：
\begin{equation*}
 V^\mu=\frac12\operatorname{Tr}(\bar\sigma^\mu V).
\end{equation*}
对变换后的矩阵，
\begin{align*}
 V'^\mu
 &=\frac12\operatorname{Tr}(\bar\sigma^\mu A\sigma^\nu A^\dagger)V_\nu\\
 &\equiv\Lambda^\mu{}_{\nu}(A)V^\nu .
\end{align*}
这里
\begin{equation*}
 \Lambda^\mu{}_{\nu}(A)
 =\frac12\operatorname{Tr}
 \left(\bar\sigma^\mu A\sigma_\nu A^\dagger\right)
\end{equation*}
正是 $SL(2,\mathbb C)$ 双覆盖诱导出的 Lorentz 矩阵。进一步，
\begin{equation*}
 \det V=V_\mu V^\mu,
 \qquad
 \det(AVA^\dagger)=|\det A|^2\det V=\det V,
\end{equation*}
所以
$V'_\mu V'^\mu=V_\mu V^\mu$；这再次验证 $\Lambda(A)$ 保持 Minkowski 度规。

因此
\begin{equation*}
 V_{\alpha\dot\beta}=V_\mu(\sigma^\mu)_{\alpha\dot\beta}
\end{equation*}
不是靠“维数碰巧都是四”猜出的字典，而是因为一左一右 Weyl 指标的矩阵变换律与四矢量的 $SL(2,\mathbb C)$ 实现逐项相同。这才证明
$(\tfrac12,0)\otimes(0,\tfrac12)=(\tfrac12,\tfrac12)$，即正文\eqnrefs{eq:weyl-vector-product-r68}。
\end{derivation}

\begin{derivation}[从\eqnrefs{eq:weyl-same-chirality-product-r68}到\eqnrefs{eq:twoform-chiral-decomposition-r68}]
对两个左手 Weyl 指标，任意二阶对象 $T_{\alpha\beta}$ 都有唯一的对称--反对称分解
\begin{equation*}
 T_{\alpha\beta}=T_{(\alpha\beta)}+T_{[\alpha\beta]}.
\end{equation*}
二维旋量空间中反对称二阶张量只有一个独立分量，因此必可写成
\begin{equation*}
 T_{[\alpha\beta]}=\phi\,\epsilon_{\alpha\beta}.
\end{equation*}
为什么这个系数 $\phi$ 是 Lorentz 标量？因为 $A\in SL(2,\mathbb C)$ 满足
\begin{equation*}
 A_\alpha{}^\gamma A_\beta{}^\delta
 \epsilon_{\gamma\delta}
 =(\det A)\epsilon_{\alpha\beta}
 =\epsilon_{\alpha\beta}.
\end{equation*}
所以反对称的一维子空间保持不变，正是 $(0,0)$。剩下的对称矩阵满足
$T_{(11)},T_{(12)},T_{(22)}$ 三个独立复分量，并在
\begin{equation*}
 T_{(\alpha\beta)}\longmapsto
 A_\alpha{}^\gamma A_\beta{}^\delta
 T_{(\gamma\delta)}
\end{equation*}
下封闭，因此构成三维不可约表示 $(1,0)$。于是从实际指标分解得到
\begin{equation*}
 (\tfrac12,0)\otimes(\tfrac12,0)
 =(0,0)\oplus(1,0).
\end{equation*}
右手带点指标完全平行，给出
\begin{equation*}
 (0,\tfrac12)\otimes(0,\tfrac12)
 =(0,0)\oplus(0,1).
\end{equation*}

现在转向四维反对称张量 $F_{\mu\nu}=-F_{\nu\mu}$。它有
$\binom42=6$ 个独立分量。复化以后定义 Hodge 算符
\begin{equation*}
 (\star F)^{\mu\nu}
 =\frac12\epsilon^{\mu\nu\rho\sigma}F_{\rho\sigma}.
\end{equation*}
在当前 Lorentz 度规和 $\epsilon^{0123}=+1$ 约定下，二形式满足
$\star^2=-I$，所以其本征值是 $\pm\ii$。相应投影算符为
\begin{equation*}
 \mathcal P_\pm=\frac12(I\mp\ii\star),
 \qquad
 F_\pm\equiv\mathcal P_\pm F,
 \qquad
 \star F_\pm=\pm\ii F_\pm.
\end{equation*}
因为
$\mathcal P_+\mathcal P_-=0$、
$\mathcal P_++\mathcal P_-=I$，六维复二形式空间严格分成两个互补的三维子空间：
\begin{equation*}
 \Lambda^2_\mathbb C
 =\Lambda^2_+\oplus\Lambda^2_-.
\end{equation*}
Lorentz 变换与 Hodge 算符对易，所以两个三维子空间各自保持不变。与上面同手征双旋量的两个三维不可约表示比较，分别对应 $(1,0)$ 与 $(0,1)$；因此
\begin{equation*}
 F^{\mu\nu}\in(1,0)\oplus(0,1),
\end{equation*}
即正文\eqnrefs{eq:twoform-chiral-decomposition-r68}。对电磁场，$F_\pm$ 的三个独立复分量可等价组织成 $\mathbf E\pm\ii c\mathbf B$；这不是另外引入的场，而是同一个六分量 $F_{\mu\nu}$ 在 Hodge 本征子空间中的两种手征投影。
\end{derivation}

\begin{derivation}[从\eqnrefs{eq:clifford-general}到\eqnrefs{eq:trace-four-gamma-r9}]

未偏振费米子振幅在完成自旋求和后会变成 Dirac 矩阵迹。正文已经把 Clifford 代数固定为\eqnrefs{eq:clifford-general}；这里仅用反对易关系与迹循环性，把它化成正文\eqnrefs{eq:trace-two-gamma-r9,eq:trace-four-gamma-r9}，从而得到后续树级截面最常用的两个迹公式。



由已经建立的关系
\begin{equation*}
 \{\gamma^\mu,\gamma^\nu\}=2\eta^{\mu\nu}\id,
 \qquad
 \Tr\id=4.
\end{equation*}
对 Clifford 关系取迹：
\begin{align*}
 \Tr(\gamma^\mu\gamma^\nu)
 +\Tr(\gamma^\nu\gamma^\mu)
 &=2\eta^{\mu\nu}\Tr\id,\\
 2\Tr(\gamma^\mu\gamma^\nu)
 &=8\eta^{\mu\nu},
\end{align*}
其中第二行使用迹的循环性。故
\begin{equation*}
 {\Tr(\gamma^\mu\gamma^\nu)=4\eta^{\mu\nu}.}
 \end{equation*}
对四个 gamma，先把最左边 $\gamma^\mu$ 逐次移过后面的矩阵：
\begin{align*}
 T^{\mu\nu\rho\sigma}
 &\equiv\Tr(\gamma^\mu\gamma^\nu\gamma^\rho\gamma^\sigma),\\
 &=\Tr[(2\eta^{\mu\nu}-\gamma^\nu\gamma^\mu)
 \gamma^\rho\gamma^\sigma],\\
 &=2\eta^{\mu\nu}\Tr(\gamma^\rho\gamma^\sigma)
 -\Tr(\gamma^\nu\gamma^\mu\gamma^\rho\gamma^\sigma).
\end{align*}
再在第二项中把 $\gamma^\mu$ 移过 $\gamma^\rho$，
\begin{align*}
 T^{\mu\nu\rho\sigma}
 &=8\eta^{\mu\nu}\eta^{\rho\sigma}
 -2\eta^{\mu\rho}\Tr(\gamma^\nu\gamma^\sigma)
 +\Tr(\gamma^\nu\gamma^\rho\gamma^\mu\gamma^\sigma).
\end{align*}
最后用迹的循环性
\begin{equation*}
 \Tr(\gamma^\nu\gamma^\rho\gamma^\mu\gamma^\sigma)
 =\Tr(\gamma^\mu\gamma^\sigma\gamma^\nu\gamma^\rho)
\end{equation*}
并重复一次反对易，得到
\begin{equation*}
 {
 \Tr(\gamma^\mu\gamma^\nu\gamma^\rho\gamma^\sigma)
 =4(\eta^{\mu\nu}\eta^{\rho\sigma}
 -\eta^{\mu\rho}\eta^{\nu\sigma}
 +\eta^{\mu\sigma}\eta^{\nu\rho}).}
 \end{equation*}
这条式子是后面所有未偏振树级费米子截面的核心；后续未偏振费米子散射的自旋求和将直接使用这一由 Clifford 代数固定的结果。
\end{derivation}

\begin{derivation}[从\eqnrefs{eq:clifford-general}到\eqnrefs{eq:trace-even-recursion-r68}]
正文\eqnrefs{eq:trace-even-recursion-r68} 已经把 $T_{2n}^{\mu_1\cdots\mu_{2n}}$ 记为长度 $2n$ 的 Dirac 迹。把最左端 $\gamma^{\mu_1}$ 逐次移到右端。第一次交换给

\begin{equation*}
 \gamma^{\mu_1}\gamma^{\mu_2}
 =2\eta^{\mu_1\mu_2}-\gamma^{\mu_2}\gamma^{\mu_1},
\end{equation*}
第二次再对残余的 $\gamma^{\mu_1}\gamma^{\mu_3}$ 使用同一关系。一直进行到第 $2n$ 个矩阵，得到
\begin{align*}
 \gamma^{\mu_1}\cdots\gamma^{\mu_{2n}}
 =&\sum_{j=2}^{2n}(-1)^j 2\eta^{\mu_1\mu_j}
 \gamma^{\mu_2}\cdots\widehat{\gamma^{\mu_j}}\cdots\gamma^{\mu_{2n}}
 \nonumber\\
 &-\gamma^{\mu_2}\cdots\gamma^{\mu_{2n}}\gamma^{\mu_1},
 \end{align*}
其中帽号表示删除该因子；最后一项前的负号来自一共经过 $2n-1$ 次交换。取迹并用迹的循环性，最后一项等于 $-T_{2n}$，因此移到左边得到 $2T_{2n}$。两边除以 $2$：
\begin{equation*}
 {
 T_{2n}^{\mu_1\cdots\mu_{2n}}
 =\sum_{j=2}^{2n}(-1)^j
 \eta^{\mu_1\mu_j}
 \Tr\!\left(
 \gamma^{\mu_2}\cdots\widehat{\gamma^{\mu_j}}\cdots\gamma^{\mu_{2n}}
 \right).}
 \end{equation*}
$n=2$ 时立刻恢复四伽马矩阵迹；$n=3$ 时得到六个 $\gamma$ 矩阵的迹的十五个度规配对。Compton、M\o ller、Bhabha 与圈积分分子中的长矩阵链都可按这一递推逐级缩短。
\end{derivation}



\begin{derivation}[从\eqnrefs{eq:clifford-general}到\eqnrefs{eq:gamma-d-sandwich-r68}]

第一章的维数正规化已经采用对复维数 $d$ 作解析延拓的约定。若后续积分采用这一调节方法，矩阵代数也必须先保持在同一个 $d$ 维参数下；这里从正文 Clifford 母式直接缩并一个指标，得到正文\eqnrefs{eq:gamma-d-sandwich-r68}。



由 Clifford 代数
\begin{align*}
 \gamma^\mu\gamma^\alpha\gamma_\mu
 &=(2\eta^{\mu\alpha}-\gamma^\alpha\gamma^\mu)\gamma_\mu,\\
 &=2\gamma^\alpha-\gamma^\alpha(\gamma^\mu\gamma_\mu).
\end{align*}
在 $d$ 维中
\begin{equation*}
 \gamma^\mu\gamma_\mu=d\id,
\end{equation*}
因此
\begin{equation*}
 {
 \gamma^\mu\gamma^\alpha\gamma_\mu=(2-d)\gamma^\alpha.}
 \end{equation*}
四维时就是 $-2\gamma^\alpha$。自能分子的约化直接使用这一恒等式。
\end{derivation}

\begin{derivation}[从\eqnrefs{eq:clifford-general}到\eqnrefs{eq:gamma-three-sandwich-r10}]
M\o ller、Bhabha、Compton 与圈分子中经常出现把一串 $\gamma$ 矩阵夹在 $\gamma^\mu$ 与 $\gamma_\mu$ 之间的结构。它们都可逐次用

\begin{equation*}
 \gamma^\mu\gamma^\nu=2\eta^{\mu\nu}-\gamma^\nu\gamma^\mu
 \end{equation*}
把左端的 $\gamma^\mu$ 向右移动得到。四维中首先有
\begin{align*}
 \gamma^\mu\slashed a\gamma_\mu
 &=a_\alpha\gamma^\mu\gamma^\alpha\gamma_\mu,\\
 &=a_\alpha(2\gamma^\alpha-\gamma^\alpha\gamma^\mu\gamma_\mu),\\
 &=(2-4)\slashed a=-2\slashed a.
 \end{align*}
对两个斜线记号，
\begin{align*}
 \gamma^\mu\slashed a\slashed b\gamma_\mu
 &=\gamma^\mu\slashed a(2b_\mu-\gamma_\mu\slashed b),\\
 &=2\slashed b\slashed a-
 (\gamma^\mu\slashed a\gamma_\mu)\slashed b,\\
 &=2\slashed b\slashed a+2\slashed a\slashed b,\\
 &=2\{\slashed a,\slashed b\}=4a\!\cdot b\,\id_4.
 \end{align*}
对三个斜线记号，先把最右的 $\slashed c$ 留在外面：
\begin{align*}
 \gamma^\mu\slashed a\slashed b\slashed c\gamma_\mu
 &=2\slashed c\slashed a\slashed b
 -\gamma^\mu\slashed a\slashed b\gamma_\mu\slashed c,\\
 &=2\slashed c\slashed a\slashed b-4(a\!\cdot b)\slashed c.
\end{align*}
再把第一项中的斜线记号重新排序，得到等价形式
\begin{equation*}
 {
 \gamma^\mu\slashed a\slashed b\slashed c\gamma_\mu
 =-2\slashed c\slashed b\slashed a.}
 \end{equation*}
把 $a,b,c$ 换成实际外腿四动量，就得到后续干涉迹中使用的短矩阵链化简式。
\end{derivation}

\begin{derivation}[从\eqnrefs{eq:gamma5-def}到\eqnrefs{eq:gamma5-four-trace-r68}]

轴矢量流和手征结构需要一个与普通四伽马迹不同的对象。正文\eqnrefs{eq:gamma5-def} 已固定 $\gamma^5$ 与取向约定；先证明相应秩四迹完全反对称，再用 $0123$ 分量固定系数，从而得到正文\eqnrefs{eq:gamma5-four-trace-r68}。



固定
\begin{equation*}
 \gamma^5=\ii\gamma^0\gamma^1\gamma^2\gamma^3,
 \qquad \epsilon^{0123}=+1.
\end{equation*}
张量
\begin{equation*}
 T^{\mu\nu\rho\sigma}
 \equiv\Tr(\gamma^5\gamma^\mu\gamma^\nu\gamma^\rho\gamma^\sigma)
\end{equation*}
先证明 $T^{\mu\nu\rho\sigma}$ 的完全反对称性。交换相邻的 $\mu,\nu$ 两个指标时，Clifford 代数给出
\begin{align*}
T^{\nu\mu\rho\sigma}
&=\Tr\!\left(\gamma^5\gamma^\nu\gamma^\mu\gamma^\rho\gamma^\sigma\right),\\
&=-\Tr\!\left(\gamma^5\gamma^\mu\gamma^\nu\gamma^\rho\gamma^\sigma\right)
+2\eta^{\mu\nu}\Tr\!\left(\gamma^5\gamma^\rho\gamma^\sigma\right).
\end{align*}
第二项为零。为看清这一点，利用 $\{\gamma^5,\gamma^\rho\}=0$ 和迹的循环性：
\begin{align*}
\Tr(\gamma^5\gamma^\rho\gamma^\sigma)
&=-\Tr(\gamma^\rho\gamma^5\gamma^\sigma),\\
&=-\Tr(\gamma^5\gamma^\sigma\gamma^\rho),\\
&=-\Tr\!\left[\gamma^5(2\eta^{\rho\sigma}-\gamma^\rho\gamma^\sigma)\right],\\
&=-2\eta^{\rho\sigma}\Tr(\gamma^5)+\Tr(\gamma^5\gamma^\rho\gamma^\sigma).
\end{align*}
而 $\Tr\gamma^5=0$，所以这个两 $\gamma$ 迹为零。于是
$T^{\nu\mu\rho\sigma}=-T^{\mu\nu\rho\sigma}$。对其余任意相邻指标交换重复同一计算，得到 $T$ 对四个 Lorentz 指标完全反对称。四维中任何秩四完全反对称张量都与 $\epsilon^{\mu\nu\rho\sigma}$ 成正比，因此写成
\begin{equation*}
T^{\mu\nu\rho\sigma}=C_5\epsilon^{\mu\nu\rho\sigma}.
\end{equation*}
比例常数 $C_5$ 由 $0123$ 分量确定：
\begin{align*}
 T^{0123}
 &=\Tr(\gamma^5\gamma^0\gamma^1\gamma^2\gamma^3),\\
 &=\Tr\left[
 \ii(\gamma^0\gamma^1\gamma^2\gamma^3)^2
 \right].
\end{align*}
把第二组 $\gamma$ 矩阵逐次穿过第一组并使用
$(\gamma^0)^2=+1$、$(\gamma^i)^2=-1$，得到
\begin{equation*}
 (\gamma^0\gamma^1\gamma^2\gamma^3)^2=-\id_4,
\end{equation*}
所以
\begin{equation*}
 T^{0123}=-4\ii.
\end{equation*}
因此在上述度规与 $\epsilon^{0123}$ 约定下
\begin{equation*}
 {
 \Tr(\gamma^5\gamma^\mu\gamma^\nu\gamma^\rho\gamma^\sigma)
 =-4\ii\epsilon^{\mu\nu\rho\sigma}.}
 \end{equation*}
若改变度规或 $\epsilon^{0123}$ 的符号，右端也必须同步改变；这就是为何约定必须在全书开头冻结。
\end{derivation}

\begin{derivation}[从\eqnrefs{eq:dirac-bilinear-basis}到\eqnrefs{eq:dirac-matrix-completeness-r10}]
任意复 $4\times4$ 矩阵形成 $16$ 维线性空间。由 Clifford 代数构造

\begin{equation*}
 \mathcal B=
 \{\id,\gamma^\mu,\sigma^{\mu\nu},\gamma^\mu\gamma^5,\gamma^5\},
 \end{equation*}
其独立分量数为
$1+4+6+4+1=16$。迹恒等式说明不同 Lorentz 类型之间在迹内积
\begin{equation*}
 (A,B)\equiv\frac14\Tr(AB)
\end{equation*}
下彼此正交；例如 $\Tr\gamma^\mu=0$、$\Tr(\gamma^\mu\gamma^5)=0$、$\Tr\sigma^{\mu\nu}=0$。因此任意矩阵 $M$ 都可唯一展开为
\begin{equation*}
 M=\frac14\sum_A \Gamma_A\Tr(\Gamma^A M),
 \end{equation*}
其中 $A$ 遍历上述 16 个基元素，$\Gamma^A$ 表示按迹度规取对偶基。把 $M$ 取成只有一个元素非零的矩阵单位 $(E_{lk})_{ij}=\delta_{il}\delta_{jk}$，有
\begin{align*}
 \delta_{il}\delta_{jk}
 &=\frac14\sum_A(\Gamma_A)_{ij}
 \Tr(\Gamma^A E_{lk}),\\
 &=\frac14\sum_A(\Gamma_A)_{ij}(\Gamma^A)_{kl}.
\end{align*}
所以得到矩阵指标层面的完备恒等式
\begin{equation*}
 {
 \frac14\sum_A
 (\Gamma_A)_{ij}(\Gamma^A)_{kl}
 =\delta_{il}\delta_{kj}.}
 \end{equation*}
Fierz 重排来自完备矩阵基上的分量展开；它把
\eqnrefs{eq:dirac-matrix-completeness-r10}
插到两个旋量双线性型之间并重新配对旋量指标。若需要某一个具体 Fierz 关系，将从这条母完备关系逐项投影，而不直接引用公式表。
\end{derivation}



\begin{derivation}[从\eqnrefs{eq:sigma-barsigma-foundation-r10}到\eqnrefs{eq:weyl-determinant-mass-shell-dp21}]

由正文\eqnrefs{eq:sigma-barsigma-foundation-r10}，四动量被写成 $p\cdot\sigma$ 与 $p\cdot\bar\sigma$ 两个厄米矩阵。这里求它们的两个本征值并相乘，证明行列式正是 Lorentz 不变量 $p^2$，从而得到正文\eqnrefs{eq:weyl-determinant-mass-shell-dp21} 的质量壳条件。



写
\begin{equation*}
 p_\mu\sigma^\mu=p^0I-\mathbf p\cdot\boldsymbol\sigma.
\end{equation*}
对任意 $2\times2$ 矩阵 $M=aI+\mathbf b\cdot\boldsymbol\sigma$，Pauli 代数给
\begin{equation*}
 (\mathbf b\cdot\boldsymbol\sigma)^2=\mathbf b^2I.
\end{equation*}
因此 $\mathbf b\cdot\boldsymbol\sigma$ 的两个本征值是 $\pm|\mathbf b|$，而 $M$ 的本征值是 $a\pm|\mathbf b|$。行列式是本征值乘积：
\begin{equation*}
 \det M=(a+|\mathbf b|)(a-|\mathbf b|)=a^2-\mathbf b^2.
\end{equation*}
取 $a=p^0$、$\mathbf b=-\mathbf p$，得到
\begin{equation*}
 \det(p_\mu\sigma^\mu)
 =(p^0)^2-\mathbf p^2=p_\mu p^\mu.
\end{equation*}
对 $p_\mu\bar\sigma^\mu=p^0I+\mathbf p\cdot\boldsymbol\sigma$ 只把 $\mathbf b$ 改号，行列式不变。

质量壳 $p^2=m^2c^2$ 因而等价于
\begin{equation*}
 \det(p\cdot\sigma)=m^2c^2.
\end{equation*}
这解释了为何有质量 Weyl 方程中 $p\cdot\sigma$ 与 $p\cdot\bar\sigma$ 的乘积恢复二阶质量壳算符 $p^2-m^2c^2$；对电子而言，一阶 Dirac 方程仍是基础动力学方程。
\end{derivation}


\begin{derivation}[从\eqnrefs{eq:dirac-matrix-completeness-r10}到\eqnrefs{eq:fierz-scalar-r68}]

正文\eqnrefs{eq:dirac-matrix-completeness-r10} 是 $4\times4$ 矩阵空间的恒等分解。把它插入四个外线旋量的指标之间并改变配对顺序，就能把 $(1,2)(3,4)$ 的标量双线性改写为 $(1,4)(3,2)$ 的五类 Dirac 双线性，得到正文\eqnrefs{eq:fierz-scalar-r68}。



从矩阵完备关系
\begin{equation*}
\delta_{il}\delta_{kj}
=\frac14\sum_A(\Gamma_A)_{ij}(\Gamma^A)_{kl}
\end{equation*}
出发。对四个普通复数值 Dirac 旋量 $u_1,u_2,u_3,u_4$，标量双线性乘积写成显式矩阵指标
\begin{equation*}
(\bar u_1u_2)(\bar u_3u_4)
=(\bar u_1)_i(u_2)_j(\bar u_3)_k(u_4)_l\,
\delta_{ij}\delta_{kl}.
\end{equation*}
配对从 $(1,2)(3,4)$ 交换为 $(1,4)(3,2)$ 时， $\delta_{ij}\delta_{kl}$ 写成 $i$ 与 $l$、$k$ 与 $j$ 相连的基展开。由完备式交换哑指标得
\begin{equation*}
\delta_{ij}\delta_{kl}
=\frac14\sum_A(\Gamma_A)_{il}(\Gamma^A)_{kj}.
\end{equation*}
代入后
\begin{align*}
(\bar u_1u_2)(\bar u_3u_4)
&=\frac14\sum_A
[(\bar u_1)_i(\Gamma_A)_{il}(u_4)_l]
[(\bar u_3)_k(\Gamma^A)_{kj}(u_2)_j]\\
&=\frac14\sum_A
(\bar u_1\Gamma_Au_4)(\bar u_3\Gamma^Au_2).
\end{align*}
现在把 16 维基分成标量、矢量、张量、轴矢量、赝标量五类。按
\begin{equation*}
\sigma^{\mu\nu}=\frac{i}{2}[\gamma^\mu,\gamma^\nu]
\end{equation*}
的归一化，对偶基在张量扇区带因子 $1/2$；轴矢量扇区的迹度规带负号。于是
\begin{align*}
(\bar u_1u_2)(\bar u_3u_4)
=\frac14\Big[&
(\bar u_1u_4)(\bar u_3u_2)
+(\bar u_1\gamma^\mu u_4)(\bar u_3\gamma_\mu u_2)\\
&+\frac12(\bar u_1\sigma^{\mu\nu}u_4)
(\bar u_3\sigma_{\mu\nu}u_2)\\
&-(\bar u_1\gamma^\mu\gamma^5u_4)
(\bar u_3\gamma_\mu\gamma^5u_2)\\
&+(\bar u_1\gamma^5u_4)(\bar u_3\gamma^5u_2)
\Big].
\end{align*}
若四个对象是 Grassmann 量子场，重排 $\psi_2$ 与 $\bar\psi_3$ 到新配对顺序还会产生反对易负号。Fierz 系数由矩阵完备性决定，统计符号则由场算符次序决定。
\end{derivation}

\begin{derivation}[从\eqnrefs{eq:sigma-definition}到\eqnrefs{eq:sigma-sigma-trace-r68}]

正文\eqnrefs{eq:sigma-definition} 已给出 $\sigma^{\mu\nu}$ 与两个 $\gamma$ 矩阵对易子的关系。将两个这样的对象相乘并展开成四伽马迹后，所有张量结构都能由已知迹公式缩并，最终得到正文\eqnrefs{eq:sigma-sigma-trace-r68}。



定义
\begin{equation*}
\sigma^{\mu\nu}=\frac{i}{2}(\gamma^\mu\gamma^\nu-\gamma^\nu\gamma^\mu).
\end{equation*}
所以
\begin{align*}
\Tr(\sigma^{\mu\nu}\sigma^{\rho\sigma})
=-\frac14\Tr\Big[&
\gamma^\mu\gamma^\nu\gamma^\rho\gamma^\sigma
-\gamma^\mu\gamma^\nu\gamma^\sigma\gamma^\rho\\
&-\gamma^\nu\gamma^\mu\gamma^\rho\gamma^\sigma
+\gamma^\nu\gamma^\mu\gamma^\sigma\gamma^\rho
\Big].
\end{align*}
对四项分别代入
\(
\Tr(\gamma^a\gamma^b\gamma^c\gamma^d)
=4(\eta^{ab}\eta^{cd}-\eta^{ac}\eta^{bd}+\eta^{ad}\eta^{bc})
\)。四个迹依次为
\begin{align*}
T_1&=4(\eta^{\mu\nu}\eta^{\rho\sigma}
-\eta^{\mu\rho}\eta^{\nu\sigma}
+\eta^{\mu\sigma}\eta^{\nu\rho}),\\
T_2&=4(\eta^{\mu\nu}\eta^{\rho\sigma}
-\eta^{\mu\sigma}\eta^{\nu\rho}
+\eta^{\mu\rho}\eta^{\nu\sigma}),\\
T_3&=4(\eta^{\mu\nu}\eta^{\rho\sigma}
-\eta^{\nu\rho}\eta^{\mu\sigma}
+\eta^{\nu\sigma}\eta^{\mu\rho}),\\
T_4&=4(\eta^{\mu\nu}\eta^{\rho\sigma}
-\eta^{\nu\sigma}\eta^{\mu\rho}
+\eta^{\nu\rho}\eta^{\mu\sigma}).
\end{align*}
原组合是 $-(T_1-T_2-T_3+T_4)/4$。代入后，四个
$\eta^{\mu\nu}\eta^{\rho\sigma}$ 项的系数为 $1-1-1+1=0$；其余两种张量结构分别收集为
\begin{align*}
-\frac14(T_1-T_2-T_3+T_4)
&=4\eta^{\mu\rho}\eta^{\nu\sigma}
-4\eta^{\mu\sigma}\eta^{\nu\rho}.
\end{align*}
因此
\begin{equation*}
{
\Tr(\sigma^{\mu\nu}\sigma^{\rho\sigma})
=4(\eta^{\mu\rho}\eta^{\nu\sigma}
-\eta^{\mu\sigma}\eta^{\nu\rho}).}
\end{equation*}
进一步令 $\rho=\mu$、$\sigma=\nu$ 并求和，得到
\begin{align*}
\Tr(\sigma^{\mu\nu}\sigma_{\mu\nu})
&=4(\delta^\mu_\mu\delta^\nu_\nu-\delta^\mu_\nu\delta^\nu_\mu)\\
&=4(16-4)=48.
\end{align*}
由于 $\sigma^{\mu\nu}\sigma_{\mu\nu}$ 是 Lorentz 标量并在不可约 Dirac 矩阵空间中与恒等式成正比，而 $\Tr I_4=4$，因此
\begin{equation*}
\sigma^{\mu\nu}\sigma_{\mu\nu}=12I_4.
\end{equation*}
这两个式子是磁矩、Pauli 形状因子和张量算符中最常用的 Clifford 收缩。
\end{derivation}

\begin{derivation}[从\eqnrefs{eq:sigma-definition}到\eqnrefs{eq:sigma-gamma5-dual-r68}]
在四维和当前约定 $\epsilon^{0123}=+1$ 下，考虑两个秩二反对称张量

\begin{equation*}
\sigma^{\mu\nu}\gamma^5,
\qquad
\epsilon^{\mu\nu\rho\sigma}\sigma_{\rho\sigma}.
\end{equation*}
两者的自由指标都是 $\mu,\nu$，并且都在 Lorentz 变换下按反对称二阶张量变换，因此它们之间只能相差一个与指标无关的常数。把关系写成
\begin{equation*}
\sigma^{\mu\nu}\gamma^5
=C_D\epsilon^{\mu\nu\rho\sigma}\sigma_{\rho\sigma}.
\end{equation*}
取 $(\mu,\nu)=(0,1)$ 来求 $C_D$。由于
$\gamma^1\gamma^0=-\gamma^0\gamma^1$，有
\begin{align*}
\sigma^{01}
&=\frac{i}{2}(\gamma^0\gamma^1-\gamma^1\gamma^0),\\
&=i\gamma^0\gamma^1.
\end{align*}
所以
\begin{align*}
\sigma^{01}\gamma^5
&=(i\gamma^0\gamma^1)
(i\gamma^0\gamma^1\gamma^2\gamma^3),\\
&=-\gamma^0\gamma^1\gamma^0\gamma^1\gamma^2\gamma^3,\\
&=+(\gamma^0)^2(\gamma^1)^2\gamma^2\gamma^3,\\
&=-\gamma^2\gamma^3.
\end{align*}
另一方面，$\epsilon^{01\rho\sigma}$ 只有 $(\rho,\sigma)=(2,3)$ 和 $(3,2)$ 两项：
\begin{align*}
\epsilon^{01\rho\sigma}\sigma_{\rho\sigma}
&=\epsilon^{0123}\sigma_{23}+\epsilon^{0132}\sigma_{32},\\
&=\sigma_{23}+(-1)(-\sigma_{23}),\\
&=2\sigma_{23}.
\end{align*}
对空间指标降指标不改变 $\sigma_{23}$ 的符号，而
\begin{align*}
\sigma_{23}=\sigma^{23}
&=\frac{i}{2}(\gamma^2\gamma^3-\gamma^3\gamma^2),\\
&=i\gamma^2\gamma^3.
\end{align*}
因此右边等于 $2i\gamma^2\gamma^3$。比较
$-\gamma^2\gamma^3=C_D(2i\gamma^2\gamma^3)$ 得到
$C_D=i/2$，于是
\begin{equation*}
{
\sigma^{\mu\nu}\gamma^5
=\frac{i}{2}\epsilon^{\mu\nu\rho\sigma}\sigma_{\rho\sigma}.}
\end{equation*}
若翻转度规或 $\epsilon^{0123}$ 约定，右端符号必须同步改变。因此该对偶恒等式的符号由度规与 $\epsilon^{0123}$ 的取向约定共同固定。
\end{derivation}


% between spacetime representation theory and local gauge symmetry.

\section{作用量、场方程与 Noether 守恒律}\label{sec:action-noether-dp61}

表示论告诉我们场可以怎样变换，却还没有规定哪一种运动由自然界实现。作用量把“允许的场表示”与“实际的动力学”连接起来，而 Noether 定理进一步把连续对称性转化为守恒量，为随后把整体相位对称局域化做好准备。

\subsection{从局域拉氏量到场方程}

质点力学里，Euler--Lagrange 方程来自“真实轨迹使作用量驻值”。场论没有改变这个原则，只是把有限个坐标 \(q_a(t)\) 换成每个时空点都有取值的场 \(\Phi_A(x)\)。因此这里首先要回答的不是“场方程长什么样”，而是怎样把普通变分法推广到局域场，并看清边界项为何决定运动方程与守恒流。


令 $\Phi_A(x)$ 表示一组经典场，$A$ 只用于枚举不同场分量。全文取 $x^0=ct$，因此 $\dd^4 x=c\,\dd t\,\dd^3 r$。若拉氏量密度只依赖同一点的场及其一阶导数，作用量定义为
\begin{equation}
 \boxed{
 S[\Phi]=\frac1c\int_{\Omega}\dd^4 x\,
 \Lag\!\left(\Phi_A,\partial_\mu\Phi_A\right).}
 \label{eq:field-action-general-dp61}
\end{equation}
这里 $\Omega$ 是所考察的时空区域；$[\Lag]=\mathrm{J\,m^{-3}}$，所以 $[S]=\mathrm{J\,s}$。作用量 $S$ 是把整段时空历史映射到一个具有作用量量纲标量的泛函；其驻值条件决定经典运动方程。

对场作任意小变分 $\Phi_A\to\Phi_A+\delta\Phi_A$。只要在边界 $\partial\Omega$ 上固定场，即 $\delta\Phi_A|_{\partial\Omega}=0$，驻值条件 $\delta S=0$ 就要求
\begin{equation}
 \boxed{
 \partial_\mu\frac{\partial\Lag}{\partial(\partial_\mu\Phi_A)}
 -\frac{\partial\Lag}{\partial\Phi_A}=0.}
 \label{eq:field-euler-lagrange-dp61}
\end{equation}
这是场的 Euler--Lagrange 方程。它与质点力学中的 Euler--Lagrange 方程完全同源；差别只在于离散自由度指标被时空点和场分量共同替代。



\subsection{连续对称怎样变成守恒流}

场方程告诉我们给定初态以后怎样演化，但还没有解释为什么某些组合在整个演化中保持不变。连续对称性的作用正是把“变换后物理不变”转化成一个局域连续性方程。与只记住“对称性对应守恒量”相比，Noether 定理更强：它直接给出守恒流的表达式，并说明守恒荷只能通过边界流动而不能在体内凭空产生。

先考虑不移动时空坐标的连续内部变换。令 $\epsilon$ 是一个与 $x$ 无关的无穷小常数，并写
\begin{equation*}
 \delta\Phi_A=\epsilon\,\Delta_A(\Phi).
\end{equation*}
若拉氏量的变化最多是一个四维散度，
\begin{equation}
 \delta\Lag=\epsilon\,\partial_\mu K^\mu,
 \label{eq:noether-total-derivative-dp61}
\end{equation}
则作用量在固定边界条件下不变。利用 Euler--Lagrange 方程把体项消去以后，剩余变化必须写成局域连续性方程
\begin{equation}
 \boxed{
 \partial_\mu j^\mu=0,
 \qquad
 j^\mu
 =\sum_A
 \frac{\partial\Lag}{\partial(\partial_\mu\Phi_A)}\,\Delta_A
 -K^\mu.}
 \label{eq:noether-current-general-dp61}
\end{equation}
这就是固定坐标内部连续对称性的 Noether 定理。\eqnrefs{eq:noether-current-general-dp61} 比“对称性对应守恒量”更具体：它告诉我们守恒流怎样从拉氏量逐项计算出来。

若把该流写成 $j^\mu=(c\rho_j,\mathbf j)$，则
\begin{equation*}
 \partial_\mu j^\mu=0
 \quad\Longleftrightarrow\quad
 \partial_t\rho_j+\boldsymbol\nabla\!\cdot\mathbf j=0.
\end{equation*}
对任意固定空间体积 $V$ 积分并用 Gauss 定理，得到
\begin{equation}
 \boxed{
 \frac{\dd}{\dd t}\int_V\dd^3 r\,\rho_j
 =-\oint_{\partial V}\mathbf j\cdot\dd\mathbf S.}
 \label{eq:noether-flux-balance-dp61}
\end{equation}
左边是区域内守恒荷的变化率，右边是穿过边界的净流出。若取全空间且场在无穷远处衰减得足够快，边界通量消失，于是总荷
\begin{equation*}
 Q_j=\int\dd^3 r\,\rho_j
\end{equation*}
不随时间变化。局域连续性方程因此比“总量守恒”更强：它还规定守恒量只能通过边界输运，不能在体内凭空产生或消失。

\begin{exbox}[自由 Dirac 场的整体 $U(1)$]
自由 Dirac 拉氏量为
\begin{equation*}
 \Lag_D
 =\bar\psi\left(\ii\hbar c\gamma^\mu\partial_\mu-\me c^2\right)\psi.
\end{equation*}
它在常数相位变换 $\psi\to \ee^{\ii\theta}\psi$、$\bar\psi\to\bar\psi \ee^{-\ii\theta}$ 下不变。Noether 流的整体倍数取决于生成元的归一化；选择粒子数流的常用归一化后，
\begin{equation}
 \boxed{n^\mu=c\,\bar\psi\gamma^\mu\psi,
 \qquad \partial_\mu n^\mu=0.}
 \label{eq:dirac-number-current-dp61}
\end{equation}
于是
\begin{equation*}
 n^0=c\psi^\dagger\psi,
 \qquad
 N=\int\dd^3 r\,\psi^\dagger\psi.
\end{equation*}
若该 Dirac 场携带电荷 $\qe$，物理电流就是
\begin{equation}
 \boxed{j^\mu=\qe n^\mu=\qe c\,\bar\psi\gamma^\mu\psi.}
 \label{eq:dirac-electric-current-from-noether-dp61}
\end{equation}
整体 $U(1)$ 对称取 $\theta$ 为时空常数，因此普通导数不产生额外项。令相位依赖时空位置 $\theta=\theta(x)$ 后，$\partial_\mu\psi$ 会额外产生与 $\partial_\mu\theta$ 成正比的项。抵消该项需要引入电磁四势，并用协变导数取代普通导数。对应关系为
\[
 \text{整体相位对称性}
 \longrightarrow \text{Noether 守恒流}
 \longrightarrow \text{局域化相位}
 \longrightarrow \text{规范联络 }A_\mu.
\]


\medskip

一般 Noether 流~\eqnrefs{eq:noether-current-general-dp61} 已经建立。对自由 Dirac 拉氏量，只需把整体相位变换 $\delta\psi=\ii\epsilon\psi$、$\delta\bar\psi=-\ii\epsilon\bar\psi$ 代入该公式。质量项与动能项的相位变化彼此抵消，所以 $K^\mu=0$；再按“每个粒子携带单位数荷”的约定重标度生成元，就得到\eqnrefs{eq:dirac-number-current-dp61}。

取无穷小整体变换
\begin{equation*}
 \delta\psi=\ii\epsilon\psi,
 \qquad
 \delta\bar\psi=-\ii\epsilon\bar\psi,
 \qquad \partial_\mu\epsilon=0.
\end{equation*}
质量项的变化为
\begin{equation*}
 \delta(\bar\psi\psi)
 =(-\ii\epsilon\bar\psi)\psi
 +\bar\psi(\ii\epsilon\psi)=0.
\end{equation*}
动能项中
\begin{align*}
 \delta\!\left(\bar\psi\gamma^\mu\partial_\mu\psi\right)
 &=(-\ii\epsilon\bar\psi)\gamma^\mu\partial_\mu\psi
 +\bar\psi\gamma^\mu\partial_\mu(\ii\epsilon\psi),\\
 &=-\ii\epsilon\bar\psi\gamma^\mu\partial_\mu\psi
 +\ii\epsilon\bar\psi\gamma^\mu\partial_\mu\psi=0,
\end{align*}
所以 $K^\mu=0$。又因为
\begin{equation*}
 \frac{\partial\Lag_D}{\partial(\partial_\mu\psi)}
 =\ii\hbar c\,\bar\psi\gamma^\mu,
 \qquad
 \frac{\partial\Lag_D}{\partial(\partial_\mu\bar\psi)}=0,
\end{equation*}
Noether 公式给出的原始流为
\begin{equation*}
 j_{\rm raw}^\mu
 =(\ii\hbar c\bar\psi\gamma^\mu)(\ii\psi)
 =-\hbar c\bar\psi\gamma^\mu\psi.
\end{equation*}
连续群生成元可以整体重标度；相应 Noether 流也按同一常数重标度。把生成元归一化成“每个粒子携带单位数荷”，定义
\begin{equation*}
 n^\mu\equiv-\frac1\hbar j_{\rm raw}^\mu
 =c\bar\psi\gamma^\mu\psi.
\end{equation*}
于是 $\partial_\mu n^\mu=0$。时间分量使用 $\bar\psi=\psi^\dagger\gamma^0$：
\begin{equation*}
 n^0=c\bar\psi\gamma^0\psi
 =c\psi^\dagger\psi.
\end{equation*}
故数荷为 $N=c^{-1}\int\dd^3 r\,n^0=\int\dd^3 r\,\psi^\dagger\psi$。
\end{exbox}

\begin{derivation}[从\eqnrefs{eq:field-action-general-dp61}到\eqnrefs{eq:field-euler-lagrange-dp61}]
正文\eqnrefs{eq:field-action-general-dp61} 把场的整段时空历史映射成作用量。对场做任意小变分时，唯一需要处理的是含 $\partial_\mu\delta\Phi_A$ 的项；把它分部积分后，边界条件 $\delta\Phi_A|_{\partial\Omega}=0$ 消去表面项，而区域内部的 $\delta\Phi_A(x)$ 又可以独立任取，因此体积分系数必须逐点为零。这正给出正文\eqnrefs{eq:field-euler-lagrange-dp61}。

从定义出发，
\begin{align*}
 \delta S
 &=\frac1c\int_\Omega\dd^4 x\,
 \left[
 \frac{\partial\Lag}{\partial\Phi_A}\delta\Phi_A
 +\frac{\partial\Lag}{\partial(\partial_\mu\Phi_A)}
 \delta(\partial_\mu\Phi_A)
 \right],\\
 &=\frac1c\int_\Omega\dd^4 x\,
 \left[
 \frac{\partial\Lag}{\partial\Phi_A}\delta\Phi_A
 +\frac{\partial\Lag}{\partial(\partial_\mu\Phi_A)}
 \partial_\mu\delta\Phi_A
 \right].
\end{align*}
第二项使用乘积求导反解：
\begin{align*}
 \frac{\partial\Lag}{\partial(\partial_\mu\Phi_A)}
 \partial_\mu\delta\Phi_A
 &=\partial_\mu\!\left[
 \frac{\partial\Lag}{\partial(\partial_\mu\Phi_A)}
 \delta\Phi_A\right]
 -\partial_\mu\!\left[
 \frac{\partial\Lag}{\partial(\partial_\mu\Phi_A)}\right]
 \delta\Phi_A.
\end{align*}
代回以后，
\begin{align*}
 \delta S
 &=\frac1c\int_\Omega\dd^4 x\,
 \left[
 \frac{\partial\Lag}{\partial\Phi_A}
 -\partial_\mu\frac{\partial\Lag}{\partial(\partial_\mu\Phi_A)}
 \right]\delta\Phi_A\\
 &\quad+\frac1c\int_\Omega\dd^4 x\,
 \partial_\mu\!\left[
 \frac{\partial\Lag}{\partial(\partial_\mu\Phi_A)}
 \delta\Phi_A\right].
\end{align*}
四维 Gauss 定理把最后一项变成 $\partial\Omega$ 上的面积分；由于边界上 $\delta\Phi_A=0$，该项消失。于是
\begin{equation*}
 \delta S
 =\frac1c\int_\Omega\dd^4 x\,
 \mathfrak E_A(x)\,\delta\Phi_A(x),
\end{equation*}
其中
\begin{equation*}
 \mathfrak E_A
 \equiv
 \frac{\partial\Lag}{\partial\Phi_A}
 -\partial_\mu\frac{\partial\Lag}{\partial(\partial_\mu\Phi_A)}.
\end{equation*}
区域内部的 $\delta\Phi_A(x)$ 可以独立任取；要使 $\delta S=0$ 对所有变分都成立，只能有 $\mathfrak E_A(x)=0$。移项即得到\eqnrefs{eq:field-euler-lagrange-dp61}。
\end{derivation}

\begin{derivation}[从\eqnrefs{eq:noether-total-derivative-dp61}到\eqnrefs{eq:noether-current-general-dp61}]
链式法则给

\begin{align*}
 \delta\Lag
 &=\frac{\partial\Lag}{\partial\Phi_A}\delta\Phi_A
 +\frac{\partial\Lag}{\partial(\partial_\mu\Phi_A)}
 \partial_\mu\delta\Phi_A.
\end{align*}
对第二项做同样的乘积拆分，
\begin{align*}
 \delta\Lag
 &=\left[
 \frac{\partial\Lag}{\partial\Phi_A}
 -\partial_\mu\frac{\partial\Lag}{\partial(\partial_\mu\Phi_A)}
 \right]\delta\Phi_A\\
 &\quad+\partial_\mu\!\left[
 \frac{\partial\Lag}{\partial(\partial_\mu\Phi_A)}
 \delta\Phi_A\right].
\end{align*}
在满足 Euler--Lagrange 方程的经典解上，第一行的方括号为零。再代入 $\delta\Phi_A=\epsilon\Delta_A$：
\begin{equation*}
 \delta\Lag
 =\epsilon\,\partial_\mu\!\left[
 \sum_A\frac{\partial\Lag}{\partial(\partial_\mu\Phi_A)}
 \Delta_A\right].
\end{equation*}
另一方面，对称性假设给 $\delta\Lag=\epsilon\partial_\mu K^\mu$。两式相减，并利用常数 $\epsilon$ 可约去，得到
\begin{equation*}
 \partial_\mu\!\left[
 \sum_A\frac{\partial\Lag}{\partial(\partial_\mu\Phi_A)}
 \Delta_A-K^\mu\right]=0,
\end{equation*}
即\eqnrefs{eq:noether-current-general-dp61}。
\end{derivation}





\section{局域 $U(1)$ 与 Maxwell--Dirac 理论}

时空表示已经确定场怎样随坐标变换；带电场还具有独立的内部相位自由度。Noether 分析说明整体 $U(1)$ 对称对应守恒电荷，但若相位基准允许随时空位置变化，普通导数会额外作用在局域相位上而失去协变性。补偿这一项所需的规范联络就是电磁四势。把协变导数写入作用量以后，Maxwell 方程与带电 Dirac 方程从同一个局域 $U(1)$ 原则得到；哈密顿量形式随后自然区分正则动量与动力学动量。

\subsection{协变导数}

自由 Dirac 场具有全局 $U(1)$ 相位对称性。若希望相位选择可以逐点改变，普通导数就会感受到局域相位的梯度，因此必须引入一个能够补偿这项变化的联络。采用全书约定，带电场在局域变换下满足
\begin{equation*}
\psi'(x)=\exp\!\left(\frac{i q_{\mathrm e}}{\hbar}\chi(x)\right)\psi(x),
\qquad q_{\mathrm e}=-e,
\end{equation*}
其中 $\chi$ 是任意实函数。要求导数项与 $\psi$ 具有相同的局域相位变换律，就必须定义
\begin{equation}
\boxed{D_\mu\equiv\partial_\mu+\frac{i q_{\mathrm e}}{\hbar}A_\mu,}
\label{eq:covariant-derivative}
\end{equation}
并要求规范势变换为
\begin{equation}
A'_\mu=A_\mu-\partial_\mu\chi.
\label{eq:gauge-A-lower}
\end{equation}
于是 $D'_\mu\psi'=\ee^{\ii q_{\mathrm e}\chi/\hbar}D_\mu\psi$。这就是“协变”二字的具体含义：$D_\mu\psi$ 与 $\psi$ 属于同一个局域表示，因此可以把自由 Dirac 项中的 $\partial_\mu$ 替换为 $D_\mu$ 而不破坏局域 $U(1)$ 对称性。

在三维记号中
\begin{equation*}
A^\mu=(\phi/c,\mathbf A),
\qquad
A_\mu=(\phi/c,-\mathbf A),
\end{equation*}
\eqnrefs{eq:gauge-A-lower} 等价于 $\mathbf A' =\mathbf A+\nabla\chi$、$\phi'=\phi-\partial_t\chi$。因此标势和矢势不是两个独立的“物理场”，而是同一个四维联络的分量。

联络的局域曲率由协变导数的对易子测量：
\begin{equation}
\boxed{
[D_\mu,D_\nu]=\frac{i q_{\mathrm e}}{\hbar}F_{\mu\nu},
\qquad
F_{\mu\nu}=\partial_\mu A_\nu-\partial_\nu A_\mu.}
\label{eq:covariant-commutator}
\end{equation}
$F_{\mu\nu}$ 在\eqnrefs{eq:gauge-A-lower} 下不变，其时空分量分别给电场与磁场。局域 $U(1)$ 因而把“为什么需要电磁势”和“为什么物理场强是势的反对称导数”连接成同一个几何结构。

\begin{derivation}[从\eqnrefs{eq:covariant-derivative}到\eqnrefs{eq:covariant-commutator}]
先把协变性要求写成对任意测试旋量都成立的算符等式
\[
D'_\mu\,U_\chi=U_\chi D_\mu,
\qquad
U_\chi(x)=\ee^{\ii q_{\mathrm e}\chi(x)/\hbar}.
\]
若局域联络仍取一阶形式
\[
D'_\mu=\partial_\mu+\frac{\ii q_{\mathrm e}}{\hbar}A'_\mu,
\]
则作用在任意 $\psi$ 上给
\begin{align*}
D'_\mu(U_\chi\psi)
={}&U_\chi\partial_\mu\psi
+\frac{\ii q_{\mathrm e}}{\hbar}U_\chi(\partial_\mu\chi)\psi
+\frac{\ii q_{\mathrm e}}{\hbar}A'_\mu U_\chi\psi,\\
U_\chi D_\mu\psi
={}&U_\chi\partial_\mu\psi
+\frac{\ii q_{\mathrm e}}{\hbar}U_\chi A_\mu\psi.
\end{align*}
由于这两式必须对任意 $\psi$ 相等，乘法算符的系数必须逐点相等：
\[
A'_\mu+\partial_\mu\chi=A_\mu,
\]
即正文\eqnrefs{eq:gauge-A-lower}。因此在给定电荷约定和一阶局域导数结构后，$A_\mu$ 的变换律不是额外选择，而是协变性唯一固定的。

再计算曲率时，不能只把 $D_\mu$ 当成普通数相乘；应让对易子作用在任意测试场上：
\begin{align*}
[D_\mu,D_\nu]\psi
={}&\frac{\ii q_{\mathrm e}}{\hbar}
\bigl[(\partial_\mu A_\nu)-(\partial_\nu A_\mu)\bigr]\psi\\
&+\left(\frac{\ii q_{\mathrm e}}{\hbar}\right)^2
(A_\mu A_\nu-A_\nu A_\mu)\psi.
\end{align*}
所有 $A\,\partial\psi$ 交叉项成对消失。对阿贝尔 $U(1)$，$A_\mu$ 是普通乘法函数，所以第二行也消失，留下
\[
[D_\mu,D_\nu]\psi
=\frac{\ii q_{\mathrm e}}{\hbar}F_{\mu\nu}\psi,
\qquad
F_{\mu\nu}=\partial_\mu A_\nu-\partial_\nu A_\mu.
\]
最后代入 $A'_\mu=A_\mu-\partial_\mu\chi$，有
\[
F'_{\mu\nu}-F_{\mu\nu}
=-\partial_\mu\partial_\nu\chi
+\partial_\nu\partial_\mu\chi=0,
\]
其中只用了光滑标量函数的混合偏导可交换。这同时解释了为什么阿贝尔场强是联络的一阶反对称导数，以及为什么非阿贝尔推广中会多出联络自身的对易子项。
\end{derivation}
\subsection{Maxwell--Dirac 方程}

局域 $U(1)$ 已经确定带电 Dirac 场必须通过 $D_\mu=\partial_\mu+\ii q_{\mathrm e} A_\mu/\hbar$ 与四势耦合。由同一作用量变分，可以同时得到物质方程、规范场方程和守恒流；再给波动算符指定因果边界条件，就得到 Green 张量与 QED 传播子共同使用的逆算符结构。

在平直时空中，由 $F_{\mu\nu}=\partial_\mu A_\nu-\partial_\nu A_\mu$ 构造的一阶导数、Lorentz 与局域 $U(1)$ 不变的最低阶规范动能是 $F_{\mu\nu}F^{\mu\nu}$。另一个二次标量 $F_{\mu\nu}{}^\star\!F^{\mu\nu}$ 是全导数，因此不改变这里讨论的局域 Euler--Lagrange 方程。电子与电磁场的最小耦合作用量于是写成
\begin{equation}
 \boxed{
 S[\psi,\bar\psi,A]
 =\frac1c\int\dd^4 x\,
 \left[
 \bar\psi\left(\ii\hbar c\gamma^\mu D_\mu-m_{\mathrm e} c^2\right)\psi
 -\frac{1}{4\mu_0}F_{\mu\nu}F^{\mu\nu}
 \right].}
 \label{eq:maxwell-dirac-action}
\end{equation}
展开协变导数，物质与规范场之间唯一的最小三点耦合为
\begin{equation}
 \boxed{
 \Lag_{\rm int}
 =-q_{\mathrm e} c\,\bar\psi\gamma^\mu\psi A_\mu.}
 \label{eq:qed-interaction-lagrangian}
\end{equation}
这个顶角随后既生成 QED 的微扰图，也生成自由电子--经典光场和自由电子--量子光场的低能约化。

把第一章已经建立的场论 Euler--Lagrange 方程分别用于 $\bar\psi$ 和 $A_\nu$，得到
\begin{equation}
 \boxed{
 \left(\ii\hbar c\slashed D-m_{\mathrm e} c^2\right)\psi=0,}
 \label{eq:dirac-em}
\end{equation}
以及
\begin{equation}
 \boxed{
 \partial_\mu F^{\mu\nu}=\mu_0 j^\nu,
 \qquad
 j^\nu\equiv q_{\mathrm e} c\,\bar\psi\gamma^\nu\psi.}
 \label{eq:maxwell-source}
\end{equation}
这里 $j^\mu$ 是物质场的电磁 Noether 流。全局相位对称性，等价地 Dirac 方程和伴随方程，给出
\begin{equation}
 \boxed{\partial_\mu j^\mu=0.}
 \label{eq:current-conservation}
\end{equation}
因此 Maxwell 源方程与物质动力学并不是两套独立输入：反对称性使 $\partial_\nu\partial_\mu F^{\mu\nu}=0$，而物质方程恰好保证相同的连续性条件。

后续介质和开放几何真正需要的是“给定源以后怎样选择因果解”。在 Lorenz 规范 $\partial_\mu A^\mu=0$ 下，熟悉的 Maxwell 方程压缩为
\begin{equation*}
 \Box A^\mu=\mu_0j^\mu,
 \qquad
 \Box\equiv\partial_\mu\partial^\mu.
\end{equation*}
给波动算符指定推迟边界条件，就定义其推迟逆核 $G_R^{(0)}$；真空中该逆核把源卷积成标准推迟势。对电子--场问题真正需要继续保留的是同一个波动算符、源项与因果逆核之间的关系：
\begin{equation*}
 A^\mu(x)=\mu_0\int\dd^4 x'\,G_R^{(0)}(x-x')j^\mu(x'),
 \qquad
 G_R^{(0)}(x-x')=0\quad(t<t').
\end{equation*}
Maxwell Green 张量把同一“带因果边界条件的逆算符”推广到色散介质、边界与开放几何；QED 传播子则把这一逆核思想提升到量子相关函数。

\subsection{正则动量与动力学动量}

从 \eqnrefs{eq:dirac-em} 出发，把时间导数单独提出。定义
\begin{equation*}
 \beta_D\equiv\gamma^0,
 \qquad
 \alpha^i\equiv\gamma^0\gamma^i,
 \qquad
 \hat{\mathbf p}\equiv-\ii\hbar\nabla.
\end{equation*}
逐项展开
\begin{align*}
 \ii\hbar c\gamma^\mu D_\mu
 &=\ii\hbar\gamma^0\partial_t
 +\ii\hbar c\gamma^i\partial_i
 -q_{\mathrm e}\gamma^0\phi
 +q_{\mathrm e} c\gamma^i A_i,
\end{align*}
将空间项与质量项移到右侧并左乘 $\gamma^0$，得到
\begin{equation}
 \boxed{
 \ii\hbar\partial_t\psi
 =\hat H_D\psi,
 \qquad
 \hat H_D
 =c\vb{\alpha}\cdot\hat{\vb{\Pi}}
 +\beta_D m_{\mathrm e} c^2+q_{\mathrm e}\phi,
 }
 \label{eq:dirac-hamiltonian-em}
\end{equation}
其中
\begin{equation}
 \hat{\vb{\Pi}}\equiv\hat{\mathbf p}-q_{\mathrm e}\mathbf A
 \label{eq:kinetic-momentum}
\end{equation}
是动力学动量。$\hat{\mathbf p}$ 是与坐标正则共轭的正则动量；在有矢势时二者不能混同。

若外场依赖时间，$\hat H_D(t)$ 在不同时刻一般不对易。定义演化算符
\begin{equation*}
 \ii\hbar\partial_t\hat U_D(t,t_0)=\hat H_D(t)\hat U_D(t,t_0),
 \qquad \hat U_D(t_0,t_0)=\id().
\end{equation*}
逐次积分得到
\begin{equation}
 \hat U_D(t,t_0)=
 \Tord\exp\!\left[-\frac{\ii}{\hbar}
 \int_{t_0}^{t}\dd t'\,\hat H_D(t')\right].
 \label{eq:dirac-time-evolution}
\end{equation}
因此在进入自由极限之前，单电子外场问题已经由 $\hat H_D[A_\mu]$ 与 \eqnrefs{eq:dirac-time-evolution} 完整定义。相互作用绘景、程函相位和 PINEM 相位调制都可看成这个演化问题在不同表示、投影和尺度极限下的约化。


动力学动量包含位置依赖的矢势，因此不同方向的动力学平移一般不再对易。它们的对易子由磁场唯一确定：
\begin{equation}
 \boxed{
 [\hat\Pi_i,\hat\Pi_j]
 =\ii q_{\mathrm e}\hbar\epsilon_{ijk}B_k.}
 \label{eq:pi-commutator}
\end{equation}
这里 $i,j,k\in\{1,2,3\}$，$\epsilon_{ijk}$ 是三维 Levi--Civita 符号。\eqnrefs{eq:pi-commutator} 是协变导数曲率关系在空间分量上的单粒子形式：磁场同时控制经典 Lorentz 力与量子动力学平移的非对易性。

同一个曲率还出现在 Dirac 算符的平方中。将一阶方程 \eqnrefs{eq:dirac-em} 与共轭的一阶因子组合后，二阶方程为
\begin{equation}
 \boxed{
 \left[
 \hbar^2D_\mu D^\mu+m_{\mathrm e}^2c^2
 +\frac{q_{\mathrm e}\hbar}{2}\sigma^{\mu\nu}F_{\mu\nu}
 \right]\psi=0.
 }
 \label{eq:dirac-squared}
\end{equation}
前三个对象的来源彼此独立而在这一式中汇合：$D_\mu D^\mu$ 是规范协变的二阶平移，$m_{\mathrm e}^2c^2$ 来自相对论质量壳，$\sigma^{\mu\nu}F_{\mu\nu}$ 则把旋量 Lorentz 生成元与电磁曲率收缩。其低能展开将产生 Pauli 磁矩、Darwin 项和自旋--轨道耦合。

四分量 Dirac 理论的低能正能约化首先写成前章建立的 $P/Q$ 消元问题；级数展开只在谱隙和算符范数给出控制参数之后进行。定义
\begin{equation}
 P_+\equiv\frac{\id()+\beta_D}{2},
 \qquad
 Q_-\equiv\frac{\id()-\beta_D}{2},
 \qquad
 P_++Q_-=\id(),
 \label{eq:dirac-positive-negative-projectors-r16}
\end{equation}
其中 $P_+$ 与 $Q_-$ 分别投影到 $\beta_D=+1$ 与 $\beta_D=-1$ 的静止能量二分量块。除 $\mathbf p=0$ 外，它们并不是自由 Dirac 哈密顿量的精确正、负能谱投影，因为 $c\boldsymbol\alpha\cdot\mathbf p$ 会混合两个块；低能正能电子支是在消去 $Q_-$ 块之后得到的有效谱支。对静态外场写
\begin{equation}
 \hat H_D=\beta_Dm_{\mathrm e} c^2+\hat{\mathcal V}+\hat{\mathcal O},
 \qquad
 \hat{\mathcal V}\equiv q_{\mathrm e}\phi,
 \qquad
 \hat{\mathcal O}\equiv c\boldsymbol\alpha\cdot\hat{\boldsymbol\Pi},
 \label{eq:dirac-even-odd-decomposition-r16}
\end{equation}
其中 $[\beta_D,\hat{\mathcal V}]=0$、$\{\beta_D,\hat{\mathcal O}\}=0$，所以 $\hat{\mathcal V}$ 保持 $P_+/Q_-$ 两个块，而 $\hat{\mathcal O}$ 只连接两个块。若研究的正能能窗满足
\begin{equation*}
 |z-m_{\mathrm e} c^2|\le\delta_\star,
 \qquad
 \|\hat{\mathcal V}\|\le V_\star,
\end{equation*}
其中 $\delta_\star$ 与 $V_\star$ 都具有能量单位 $\mathrm J$，则 $\beta_D=-1$ 块的谱与该能窗之间至少相隔
\begin{equation*}
 \Delta_-
 \equiv2m_{\mathrm e} c^2-\delta_\star-V_\star.
\end{equation*}
只要 $\Delta_->0$，谱定理对预解式给出
\begin{align}
 R_-(z)
 &\equiv\left[z-Q_-\hat H_DQ_-\right]^{-1},\\
 \|R_-(z)\|
 &=\frac{1}{\operatorname{dist}\!\left(z,\sigma(Q_-\hat H_DQ_-)\right)},\\
 &\le\frac1{\Delta_-}.
 \label{eq:dirac-negative-resolvent-bound-r16}
\end{align}
这里 $\|\cdot\|$ 是算符范数，$\sigma(\cdot)$ 表示算符谱。因此两个静止能量块之间混合的直接无量纲控制量可以取为
\begin{equation}
 \boxed{
 \epsilon_{\rm NR}
 \equiv
 \frac{\|P_+\hat{\mathcal O}Q_-\|}{\Delta_-}\ll1.}
 \label{eq:dirac-pauli-control-r16}
\end{equation}
在弱外场且典型动量为 $p_\star$ 时，$\|\hat{\mathcal O}\|\sim cp_\star$、$\Delta_-\simeq2m_{\mathrm e} c^2$，所以 $\epsilon_{\rm NR}\sim p_\star/(2m_{\mathrm e} c)$。该比值把正能投影与真实的正负能谱隙联系起来，从而给出消去负能子空间的定量条件。

受控低能分支的静态电磁哈密顿量可由 Foldy--Wouthuysen 变换逐阶分块对角化。时变外场还会从幺正变换的 $-\ii\hbar U\partial_tU^\dagger$ 项产生含 $\partial_t\mathbf E$ 的附加项，因此静态展开只适用于外场时间变化满足相应低频条件的情形。

Foldy--Wouthuysen 变换用幺正变换逐阶消去关于 $\beta_D$ 的奇算符。对静态外场展开到 $1/(m_{\mathrm e} c^2)$ 的相对论低能修正阶，分块对角的偶算符结构为
\begin{equation}
 \boxed{
 H_{\rm FW}
 =\beta_Dm_{\mathrm e} c^2+\hat{\mathcal V}
 +\frac{\beta_D\hat{\mathcal O}^2}{2m_{\mathrm e} c^2}
 -\frac{\beta_D\hat{\mathcal O}^4}{8m_{\mathrm e}^3c^6}
 -\frac{1}{8m_{\mathrm e}^2c^4}[\hat{\mathcal O},[\hat{\mathcal O},\hat{\mathcal V}]]
 +\cdots.}
 \label{eq:fw-even-operator-form-r16}
\end{equation}
这条式子先保持抽象的奇、偶算符，因此与具体电磁场分量无关；电磁结构只进入 $\hat{\mathcal O}^2$ 与双重对易子。定义无量纲四分量自旋矩阵
\begin{equation*}
 \boldsymbol\Sigma
 \equiv
 \begin{pmatrix}
 \boldsymbol\sigma_{\rm P}&0\\
 0&\boldsymbol\sigma_{\rm P}
 \end{pmatrix},
 \qquad
 \mathbf S=\frac{\hbar}{2}\boldsymbol\Sigma,
\end{equation*}
其中 $\boldsymbol\sigma_{\rm P}$ 是 Pauli 矩阵向量。利用动力学动量曲率，奇算符平方化为
\begin{equation}
 \boxed{
 \hat{\mathcal O}^2
 =c^2\left(
 \boldsymbol\Pi^2-q_{\mathrm e}\hbar\boldsymbol\Sigma\cdot\mathbf B
 \right).}
 \label{eq:fw-O2-r16}
\end{equation}
第一项给出动能，第二项已经包含自旋--磁场耦合。对静态标势 $\hat{\mathcal V}=q_{\mathrm e}\phi$，双重对易子进一步给出
\begin{equation}
 \boxed{
 \begin{aligned}
 [\hat{\mathcal O},[\hat{\mathcal O},\hat{\mathcal V}]]
 ={}&q_{\mathrm e}\hbar^2c^2\,\boldsymbol\nabla\cdot\mathbf E\\
 &+q_{\mathrm e}\hbar c^2\,
 \boldsymbol\Sigma\cdot
 \left(
 \mathbf E\times\boldsymbol\Pi
 -\boldsymbol\Pi\times\mathbf E
 \right).
 \end{aligned}}
 \label{eq:fw-double-commutator-r16}
\end{equation}
其中散度项来自 $[\Pi_i,E_i]$，叉乘项来自 Pauli 矩阵乘法中的反对称部分。把这两个电磁恒等式代回 \eqnrefs{eq:fw-even-operator-form-r16} 并投影到低能正能块，得到
\begin{equation}
 \boxed{
 \begin{aligned}
 \hat H_{+}
 ={}&m_{\mathrm e} c^2+q_{\mathrm e}\phi
 +\frac{\hat{\boldsymbol\Pi}^2}{2m_{\mathrm e}}
 -\frac{q_{\mathrm e}\hbar}{2m_{\mathrm e}}\boldsymbol\sigma_{\rm P}\cdot\mathbf B
 -\frac{\left(\hat{\boldsymbol\Pi}^2-q_{\mathrm e}\hbar\boldsymbol\sigma_{\rm P}\cdot\mathbf B\right)^2}{8m_{\mathrm e}^3c^2}\\
 &-\frac{q_{\mathrm e}\hbar^2}{8m_{\mathrm e}^2c^2}\boldsymbol\nabla\cdot\mathbf E
 -\frac{q_{\mathrm e}\hbar}{8m_{\mathrm e}^2c^2}
 \boldsymbol\sigma_{\rm P}\cdot
 \left(\mathbf E\times\hat{\boldsymbol\Pi}
 -\hat{\boldsymbol\Pi}\times\mathbf E\right)+\cdots.
 \end{aligned}}
 \label{eq:fw-pauli}
\end{equation}
依次可辨认出静质量、标势、非相对论动能、磁矩、相对论动能与弱场联合修正、Darwin 项和自旋--轨道项。这里四次奇算符项保留为完整算符平方；由于 $\hat{\boldsymbol\Pi}^2$ 与位置依赖的 $\mathbf B$ 一般不对易，不能在没有额外弱场计数时把它无条件替换成 $\hat{\boldsymbol\Pi}^4$。若进一步满足 $|q_{\mathrm e}|\hbar B_\star\ll p_\star^2$，才可在这一项内舍去磁场交叉项并恢复常见的 $-\hat{\boldsymbol\Pi}^4/(8m_{\mathrm e}^3c^2)$ 形式。所有这些项仍来自同一个 Dirac Hamiltonian 的分块对角化。
为把上述低能展开写成定量适用条件，令 $p_\star$ 是电子典型动力学动量，$\phi_\star,E_\star,B_\star$ 是外场在作用区的典型幅值，$L_\star$ 与 $\omega_\star$ 分别是其最短空间变化尺度和最高相关角频率，并定义 $\lambda_C\equiv\hbar/(m_{\mathrm e} c)$。一组足够的无量纲控制量是
\begin{equation}
 \max\!\left\{
 \frac{p_\star}{m_{\mathrm e} c},
 \frac{|q_{\mathrm e}|\,|\phi_\star|}{m_{\mathrm e} c^2},
 \frac{\hbar\omega_\star}{m_{\mathrm e} c^2},
 \frac{\lambda_C}{L_\star},
 \frac{|q_{\mathrm e}|\hbar B_\star}{m_{\mathrm e}^2c^2},
 \frac{|q_{\mathrm e}|E_\star\lambda_C}{m_{\mathrm e} c^2}
 \right\}\ll1.
 \label{eq:fw-control-set-r16}
\end{equation}
其中各项分别控制动量、标势、时间变化、空间梯度、磁场能标与电场在 Compton 长度上的势降。对数十至数百 keV 电子，第一项通常并不小，因此自由电子动力学必须保留完整相对论色散；\eqnrefs{eq:fw-pauli} 只在满足 \eqnrefs{eq:dirac-pauli-control-r16,eq:fw-control-set-r16} 时作为低能特例。


\begin{derivation}[从\eqnrefs{eq:kinetic-momentum}到\eqnrefs{eq:pi-commutator}]

正文\eqnrefs{eq:kinetic-momentum} 已写出动力学动量与正则动量、矢势之间的关系。不同空间分量的非对易性因此完全来自 $[\hat p_i,A_j]$；把位置表象的基本对易关系逐项展开，再用 $\mathbf B=\boldsymbol\nabla\times\mathbf A$ 收集反对称导数，就得到正文\eqnrefs{eq:pi-commutator}。


由 $\hat\Pi_i=\hat p_i-q_{\mathrm e} A_i(\mathbf r,t)$ 与 $[\hat p_i,f(\mathbf r,t)]=-\ii\hbar\partial_i f$，逐项展开
\begin{align*}
 [\hat\Pi_i,\hat\Pi_j]
 &=[\hat p_i-q_{\mathrm e} A_i,\hat p_j-q_{\mathrm e} A_j]\\
 &=[\hat p_i,\hat p_j]
 -q_{\mathrm e}[\hat p_i,A_j]
 -q_{\mathrm e}[A_i,\hat p_j]
 +q_{\mathrm e}^2[A_i,A_j].
\end{align*}
坐标表象中的普通阿贝尔势是乘法算符，因此 $[A_i,A_j]=0$，而正则动量满足 $[\hat p_i,\hat p_j]=0$。剩余两项为
\begin{align*}
 [\hat\Pi_i,\hat\Pi_j]
 &=-q_{\mathrm e}(-\ii\hbar\partial_iA_j)
 -q_{\mathrm e}(+\ii\hbar\partial_jA_i)\\
 &=\ii q_{\mathrm e}\hbar(\partial_iA_j-\partial_jA_i).
\end{align*}
利用 $B_k=\epsilon_{k\ell m}\partial_\ell A_m$ 以及
\begin{align*}
 \epsilon_{ijk}B_k
 &=\epsilon_{ijk}\epsilon_{k\ell m}\partial_\ell A_m\\
 &=(\delta_{i\ell}\delta_{jm}-\delta_{im}\delta_{j\ell})
 \partial_\ell A_m\\
 &=\partial_iA_j-\partial_jA_i,
\end{align*}
得到主文公式。
\end{derivation}

\begin{derivation}[从\eqnrefs{eq:dirac-em}到\eqnrefs{eq:dirac-squared}]

正文\eqnrefs{eq:dirac-em} 是含任意时空依赖电磁场的一阶 Dirac 方程。左乘符号相反的一阶因子后，协变导数乘积的对称部分给出 $D_\mu D^\mu$，反对称部分则通过 $[D_\mu,D_\nu]$ 化成场强并与 $\sigma^{\mu\nu}$ 收缩，从而得到正文\eqnrefs{eq:dirac-squared}。


将 \eqnrefs{eq:dirac-em} 除以 $c$，
\begin{equation*}
 (\ii\hbar\slashed D-m_{\mathrm e} c)\psi=0.
\end{equation*}
左乘共轭的一阶因子 $\ii\hbar\slashed D+m_{\mathrm e} c$：
\begin{align*}
 0
 &=(\ii\hbar\slashed D+m_{\mathrm e} c)(\ii\hbar\slashed D-m_{\mathrm e} c)\psi\\
 &=\left[(\ii\hbar)^2\gamma^\mu\gamma^\nu D_\mu D_\nu-m_{\mathrm e}^2c^2\right]\psi\\
 &=\left[-\hbar^2\gamma^\mu\gamma^\nu D_\mu D_\nu-m_{\mathrm e}^2c^2\right]\psi.
\end{align*}
由 $\sigma^{\mu\nu}=\ii[\gamma^\mu,\gamma^\nu]/2$ 与 Clifford 代数，
\begin{align*}
 \gamma^\mu\gamma^\nu
 &=\frac12\{\gamma^\mu,\gamma^\nu\}
 +\frac12[\gamma^\mu,\gamma^\nu]\\
 &=\eta^{\mu\nu}\id()-\ii\sigma^{\mu\nu}.
\end{align*}
因此
\begin{align*}
 \gamma^\mu\gamma^\nu D_\mu D_\nu
 &=D_\mu D^\mu-\ii\sigma^{\mu\nu}D_\mu D_\nu.
\end{align*}
由于 $\sigma^{\mu\nu}$ 对 $\mu,\nu$ 反对称，$D_\mu D_\nu$ 的对称部分与它收缩为零：
\begin{align*}
 \sigma^{\mu\nu}D_\mu D_\nu
 &=\frac12\sigma^{\mu\nu}
 (D_\mu D_\nu-D_\nu D_\mu)\\
 &=\frac12\sigma^{\mu\nu}[D_\mu,D_\nu]\\
 &=\frac{\ii q_{\mathrm e}}{2\hbar}\sigma^{\mu\nu}F_{\mu\nu}.
\end{align*}
于是
\begin{align*}
 \gamma^\mu\gamma^\nu D_\mu D_\nu
 &=D_\mu D^\mu
 +\frac{q_{\mathrm e}}{2\hbar}\sigma^{\mu\nu}F_{\mu\nu}.
\end{align*}
代回二阶方程并整体乘以 $-1$：
\begin{align*}
 0
 &=\left[
 \hbar^2D_\mu D^\mu+m_{\mathrm e}^2c^2
 +\frac{q_{\mathrm e}\hbar}{2}\sigma^{\mu\nu}F_{\mu\nu}
 \right]\psi,
\end{align*}
即\eqnrefs{eq:dirac-squared}。
\end{derivation}

\begin{derivation}[从\eqnrefs{eq:dirac-even-odd-decomposition-r16}到\eqnrefs{eq:fw-even-operator-form-r16}]
取第一步反厄米生成元

\begin{equation*}
 A_1\equiv\ii S_1=\frac{\beta_D\hat{\mathcal O}}{2m_{\mathrm e} c^2},
 \qquad U_1=\ee^{A_1}.
\end{equation*}
由于 $\{\beta_D,\hat{\mathcal O}\}=0$，有 $(\beta_D\hat{\mathcal O})^\dagger=-\beta_D\hat{\mathcal O}$，所以 $U_1$ 幺正。Baker--Campbell--Hausdorff 展开为
\begin{equation*}
 H^{(1)}=H_D+[A_1,H_D]
 +\frac12[A_1,[A_1,H_D]]
 +\frac1{3!}[A_1,[A_1,[A_1,H_D]]]+\cdots .
\end{equation*}
对静质量与奇算符的嵌套对易子逐阶计算，
\begin{align*}
 [A_1,\beta_Dm_{\mathrm e} c^2]&=-\hat{\mathcal O},\\
 [A_1,\hat{\mathcal O}]&=\frac{\beta_D\hat{\mathcal O}^2}{m_{\mathrm e} c^2},\\
 [A_1,[A_1,\beta_Dm_{\mathrm e} c^2]]&=-\frac{\beta_D\hat{\mathcal O}^2}{m_{\mathrm e} c^2},\\
 [A_1,[A_1,\hat{\mathcal O}]]&=-\frac{\hat{\mathcal O}^3}{m_{\mathrm e}^2c^4}.
\end{align*}
继续到四阶并按 BCH 系数收集，原来的 $+\hat{\mathcal O}$ 被一阶项抵消，偶部分分别留下
\begin{equation*}
 \frac{\beta_D\hat{\mathcal O}^2}{2m_{\mathrm e} c^2},
 \qquad
 -\frac{\beta_D\hat{\mathcal O}^4}{8m_{\mathrm e}^3c^6}.
\end{equation*}
偶势满足 $[\beta_D,\hat{\mathcal V}]=0$。第一层对易子为奇算符，第二层对易子为
\begin{equation*}
 [A_1,[A_1,\hat{\mathcal V}]]
 =-\frac{1}{4m_{\mathrm e}^2c^4}[\hat{\mathcal O},[\hat{\mathcal O},\hat{\mathcal V}]],
\end{equation*}
乘上 BCH 的 $1/2!$ 后得到
\begin{equation*}
 -\frac{1}{8m_{\mathrm e}^2c^4}[\hat{\mathcal O},[\hat{\mathcal O},\hat{\mathcal V}]].
\end{equation*}
第一次变换仍留下三阶奇算符
\begin{equation*}
 \hat{\mathcal O}_1
 =\frac{\beta_D[\hat{\mathcal O},\hat{\mathcal V}]}{2m_{\mathrm e} c^2}
 -\frac{\hat{\mathcal O}^3}{3m_{\mathrm e}^2c^4}+\cdots .
\end{equation*}
再取 $A_2=\beta_D\hat{\mathcal O}_1/(2m_{\mathrm e} c^2)$，其与静质量的对易子正好给 $-\hat{\mathcal O}_1$，因此这一奇块被下一步变换消去。按 $\|\hat{\mathcal O}\|/(m_{\mathrm e} c^2)=O(\epsilon)$、$\|\hat{\mathcal V}\|/(m_{\mathrm e} c^2)=O(\epsilon^2)$ 计数，第二步对偶部分的首个新修正高于当前保留阶数。于是得到正文\eqnrefs{eq:fw-even-operator-form-r16}。
\end{derivation}

\begin{derivation}[从\eqnrefs{eq:pi-commutator}到\eqnrefs{eq:fw-O2-r16}]

正文\eqnrefs{eq:pi-commutator} 已说明动力学动量的反对称部分由磁场控制。对最小耦合奇算符 $\hat{\mathcal O}=c\boldsymbol\alpha\cdot\boldsymbol\Pi$ 求平方时，Pauli 代数的对称部分给普通动能，反对称部分恰好读取这个动量对易子，因此磁矩项由最小耦合动量的非对易部分与 Dirac--Pauli 代数共同生成。

利用
\begin{equation*}
 \alpha_i\alpha_j
 =\delta_{ij}\id()+\ii\epsilon_{ijk}\Sigma_k,
\end{equation*}
有
\begin{align*}
 \hat{\mathcal O}^2
 &=c^2\alpha_i\alpha_j\Pi_i\Pi_j,\\
 &=c^2\left[
 \boldsymbol\Pi^2
 +\frac{\ii}{2}\epsilon_{ijk}\Sigma_k[\Pi_i,\Pi_j]
 \right].
\end{align*}
代入正文\eqnrefs{eq:pi-commutator}，
\begin{equation*}
 \frac{\ii}{2}\epsilon_{ijk}\Sigma_k[\Pi_i,\Pi_j]
 =-q_{\mathrm e}\hbar\boldsymbol\Sigma\cdot\mathbf B,
\end{equation*}
从而得到正文\eqnrefs{eq:fw-O2-r16}。
\end{derivation}

\begin{derivation}[从\eqnrefs{eq:dirac-even-odd-decomposition-r16}到\eqnrefs{eq:fw-double-commutator-r16}]

正文\eqnrefs{eq:dirac-even-odd-decomposition-r16} 中 $\hat{\mathcal V}=q_{\mathrm e}\phi$、$\hat{\mathcal O}=c\boldsymbol\alpha\cdot\boldsymbol\Pi$。对静态标势，第一层对易子只依赖 $[\Pi_i,\phi]$；再与 $\hat{\mathcal O}$ 对易一次后，Pauli 代数的对称部分产生 $\nabla\cdot\mathbf E$，反对称部分产生自旋与 $\mathbf E\times\boldsymbol\Pi$ 的耦合。

因为 $\mathbf E=-\boldsymbol\nabla\phi$，
\begin{align*}
 [\hat{\mathcal O},\hat{\mathcal V}]
 &=q_{\mathrm e} c\alpha_i[\Pi_i,\phi],\\
 &=\ii q_{\mathrm e}\hbar c\,\boldsymbol\alpha\cdot\mathbf E.
\end{align*}
于是
\begin{align*}
 [\hat{\mathcal O},[\hat{\mathcal O},\hat{\mathcal V}]]
 &=\ii q_{\mathrm e}\hbar c^2[\alpha_i\Pi_i,\alpha_jE_j],\\
 &=\ii q_{\mathrm e}\hbar c^2
 \left
 \{[\Pi_i,E_i]
 +\ii\epsilon_{ijk}\Sigma_k(\Pi_iE_j+E_j\Pi_i)
 \right\}.
\end{align*}
再用 $[\Pi_i,E_i]=-\ii\hbar\partial_iE_i$ 并整理 Levi--Civita 收缩，得到
\begin{equation*}
 [\hat{\mathcal O},[\hat{\mathcal O},\hat{\mathcal V}]]
 =q_{\mathrm e}\hbar^2c^2\boldsymbol\nabla\cdot\mathbf E
 +q_{\mathrm e}\hbar c^2\boldsymbol\Sigma\cdot
 \left(\mathbf E\times\boldsymbol\Pi-\boldsymbol\Pi\times\mathbf E\right),
\end{equation*}
即正文\eqnrefs{eq:fw-double-commutator-r16}。
\end{derivation}


\begin{derivation}[从\eqnrefs{eq:fw-even-operator-form-r16}到\eqnrefs{eq:fw-pauli}]
令
\[
P_+\equiv\frac{\id()+\beta_D}{2}.
\]
因为 $P_+\beta_D P_+=P_+$，而 $\boldsymbol\Sigma$ 在正、负能静止块内分别退化成同一组 Pauli 矩阵，所以
\[
P_+\boldsymbol\Sigma P_+
=\boldsymbol\sigma_{\rm P}P_+.
\]
对任意只作用在正能块内的态，正文\eqnrefs{eq:fw-O2-r16} 因而给出
\[
P_+\hat{\mathcal O}^2P_+
=c^2\Bigl(
\hat{\boldsymbol\Pi}^2
-q_{\mathrm e}\hbar\boldsymbol\sigma_{\rm P}\cdot\mathbf B
\Bigr)P_+.
\]
四次奇算符不能再单独把磁场项丢掉。由 $\hat{\mathcal O}^4=(\hat{\mathcal O}^2)^2$，有严格的算符恒等式
\begin{align*}
P_+\hat{\mathcal O}^4P_+
={}&c^4
\Bigl(
\hat{\boldsymbol\Pi}^2
-q_{\mathrm e}\hbar\boldsymbol\sigma_{\rm P}\cdot\mathbf B
\Bigr)^2P_+\\
={}&c^4\Bigl[
\hat{\boldsymbol\Pi}^4
-q_{\mathrm e}\hbar\bigl\{
\hat{\boldsymbol\Pi}^2,
\boldsymbol\sigma_{\rm P}\cdot\mathbf B
\bigr\}
+q_{\mathrm e}^2\hbar^2
(\boldsymbol\sigma_{\rm P}\cdot\mathbf B)^2
\Bigr]P_+,
\end{align*}
其中反对易子保留了 $\hat{\boldsymbol\Pi}$ 与非均匀磁场不对易的信息。把这两式以及正文\eqnrefs{eq:fw-double-commutator-r16} 代入\eqnrefs{eq:fw-even-operator-form-r16}，得到
\begin{align*}
P_+H_{\rm FW}P_+
={}&\Biggl[
m_{\mathrm e} c^2+q_{\mathrm e}\phi
+\frac{\hat{\boldsymbol\Pi}^2-q_{\mathrm e}\hbar\boldsymbol\sigma_{\rm P}\cdot\mathbf B}{2m_{\mathrm e}}\\
&-\frac{\left(\hat{\boldsymbol\Pi}^2-q_{\mathrm e}\hbar\boldsymbol\sigma_{\rm P}\cdot\mathbf B\right)^2}{8m_{\mathrm e}^3c^2}
-\frac{q_{\mathrm e}\hbar^2}{8m_{\mathrm e}^2c^2}\boldsymbol\nabla\cdot\mathbf E\\
&-\frac{q_{\mathrm e}\hbar}{8m_{\mathrm e}^2c^2}
\boldsymbol\sigma_{\rm P}\cdot
\left(
\mathbf E\times\hat{\boldsymbol\Pi}
-\hat{\boldsymbol\Pi}\times\mathbf E
\right)
\Biggr]P_+ + P_+R_{\rm FW}P_+ .
\end{align*}
$R_{\rm FW}$ 收集前一推导中没有保留的更高阶 Foldy--Wouthuysen 项。去掉末尾的投影算符，就是正文\eqnrefs{eq:fw-pauli}。若还要使用更常见的纯 $\Pi^4$ 相对论动能修正，则必须另外控制
\[
\epsilon_{B/p}
\equiv
\frac{|q_{\mathrm e}|\hbar B_\star}{p_\star^2}\ll1;
\]
这时四次项中的反对易子与 $B^2$ 项才相对于 $\Pi^4$ 受控变小。这个附加条件与正负能块解耦条件\eqnrefs{eq:dirac-pauli-control-r16} 是不同的近似。
\end{derivation}


\subsection{自由场与相对论运动学}

把电磁联络置零以后，Maxwell--Dirac 理论的电子部分退化为自由 Dirac 方程
\begin{equation}
 (\ii\hbar\slashed\partial-m_{\mathrm e} c)\psi=0.
 \label{eq:free-dirac}
\end{equation}
令 $A_\mu=0$ 得到含场方程的自由极限；它同时是自由传播、散射外线态与反冲运动学的共同母方程。对正频平面波
\begin{equation*}
 \psi(x)=u_s(p)\ee^{-\ii p\cdot x/\hbar},
 \qquad p^0=E_{\mathbf p}/c>0,
\end{equation*}
\eqnrefs{eq:free-dirac} 给出
\begin{equation}
 (\slashed p-m_{\mathrm e} c)u_s(p)=0,
 \label{eq:u-dirac}
\end{equation}
而负频分支写成 $v_s(p)\ee^{+\ii p\cdot x/\hbar}$，满足
\begin{equation}
 (\slashed p+m_{\mathrm e} c)v_s(p)=0.
 \label{eq:v-dirac}
\end{equation}
两类旋量都用未来指向的 $p^0>0$ 参数化；量子化后，$u$ 与 $v$ 分支分别乘电子湮灭算符与正电子产生算符。

Clifford 代数使一阶 Dirac 算符成为相对论质量壳的因式分解，因此非零自由解必须满足
\begin{equation}
 \boxed{
 p^2=m_{\mathrm e}^2c^2,
 \qquad
 E_{\mathbf p}=\sqrt{m_{\mathrm e}^2c^4+c^2\mathbf p^2}.}
 \label{eq:rel-dispersion}
\end{equation}
这条精确色散关系是传播与反冲近似的共同母式；任何近轴、低反冲或二阶色散公式都必须从它展开，而不能另行假设抛物线能谱。

为固定外线旋量归一化，取两组二分量正交基 $\xi_s^{(+)}$ 与 $\xi_s^{(-)}$，分别对应正能与负能支，并满足
\[
(\xi_s^{(+)})^\dagger\xi_{s'}^{(+)}=(\xi_s^{(-)})^\dagger\xi_{s'}^{(-)}=\delta_{ss'}.
\]
在 Dirac 表示中可选
\begin{equation}
 u_s(p)=
 \sqrt{\frac{E}{c}+m_{\mathrm e} c}
 \begin{pmatrix}
 \xi_s^{(+)}\\[2mm]
 \dfrac{\boldsymbol\sigma_{\rm P}\cdot\mathbf p}{E/c+m_{\mathrm e} c}\,\xi_s^{(+)}
 \end{pmatrix},
 \label{eq:explicit-u}
\end{equation}
负能支由同一分块方程得到，只是大、小分量的位置互换：
\begin{equation}
 v_s(p)=
 \sqrt{\frac{E}{c}+m_{\mathrm e} c}
 \begin{pmatrix}
 \dfrac{\boldsymbol\sigma_{\rm P}\cdot\mathbf p}{E/c+m_{\mathrm e} c}\,\xi_s^{(-)}\\[2mm]
 \xi_s^{(-)}
 \end{pmatrix},
 \label{eq:explicit-v}
\end{equation}
并采用
$\bar u_su_{s'}=2m_{\mathrm e} c\,\delta_{ss'}$、
$\bar v_sv_{s'}=-2m_{\mathrm e} c\,\delta_{ss'}$。
比显式分量更常用的是在任意质量壳四动量上直接构造能量投影算符。由 $\slashed p^{\,2}=m_{\mathrm e}^2c^2$ 定义
\begin{equation*}
\Lambda_+(p)\equiv\frac{\slashed p+m_{\mathrm e} c}{2m_{\mathrm e} c},
\qquad
\Lambda_-(p)\equiv\frac{m_{\mathrm e} c-\slashed p}{2m_{\mathrm e} c}.
\end{equation*}
它们满足 $\Lambda_\pm^2=\Lambda_\pm$、$\Lambda_+\Lambda_-=0$ 且迹均为 $2$；$\Lambda_+$ 的像空间正是 $(\slashed p-m_{\mathrm e} c)u=0$ 的二维正能解空间，$\Lambda_-$ 的像空间正是 $(\slashed p+m_{\mathrm e} c)v=0$ 的二维负频解空间。结合上面的 Dirac 归一化，得到协变完备关系
\begin{align}
 \boxed{\sum_su_s(p)\bar u_s(p)=\slashed p+m_{\mathrm e} c,}
 &\qquad \label{eq:u-complete}\\
 \boxed{\sum_sv_s(p)\bar v_s(p)=\slashed p-m_{\mathrm e} c.}
 &\qquad \label{eq:v-complete}
\end{align}
它们由任意在壳 $p$ 上的 Dirac 投影算符直接给出；在散射振幅里，外部自旋求和因此可以化成伽马矩阵迹，同时固定 Dirac 传播子的分子结构。

\begin{exbox}[沿 $z$ 方向运动的电子]
取 $\mathbf p=p\,\mathbf e_z$、$p\ge0$，并在 $z$ 轴自旋基中选
\begin{equation*}
 \xi_\uparrow=\begin{pmatrix}1\\0\end{pmatrix},
 \qquad
 \xi_\downarrow=\begin{pmatrix}0\\1\end{pmatrix}.
\end{equation*}
记
\begin{equation*}
 A\equiv\sqrt{\frac{E}{c}+m_{\mathrm e} c},
 \qquad
 \kappa\equiv\frac{p}{E/c+m_{\mathrm e} c},
 \qquad
 E=\sqrt{m_{\mathrm e}^2c^4+c^2p^2}.
\end{equation*}
由\eqnrefs{eq:explicit-u}，两个正能旋量就是
\begin{equation}
 u_\uparrow(p)=A
 \begin{pmatrix}1\\0\\\kappa\\0\end{pmatrix},
 \qquad
 u_\downarrow(p)=A
 \begin{pmatrix}0\\1\\0\\-\kappa\end{pmatrix}.
 \label{eq:explicit-z-spinors-dp62}
\end{equation}
这里小分量与大分量之比恰为 $\kappa$。低速时 $p\ll m_{\mathrm e} c$，有 $\kappa\simeq p/(2m_{\mathrm e} c)\ll1$；超相对论时 $E\gg m_{\mathrm e} c^2$，则 $\kappa\to1$。因此所谓 Dirac 旋量的“大、小分量”不是绝对分类，而由动量尺度控制。

以 $u_\uparrow$ 为例验证归一化。Dirac 表示中 $\gamma^0=\operatorname{diag}(1,1,-1,-1)$，所以
\begin{align*}
 \bar u_\uparrow u_\uparrow
 &=A^2(1-\kappa^2)\\
 &=\left(\frac Ec+m_{\mathrm e} c\right)
 \left[1-\frac{p^2}{(E/c+m_{\mathrm e} c)^2}\right]\\
 &=\frac{(E/c+m_{\mathrm e} c)^2-p^2}{E/c+m_{\mathrm e} c}\\
 &=2m_{\mathrm e} c,
\end{align*}
最后一步使用 $E^2/c^2-p^2=m_{\mathrm e}^2c^2$。类似地
\begin{align*}
 u_\uparrow^\dagger u_\uparrow
 &=A^2(1+\kappa^2)
 =\frac{2E}{c}.
\end{align*}
因此时间分量的 Dirac 流为
\begin{equation*}
 \bar u_\uparrow\gamma^0u_\uparrow
 =u_\uparrow^\dagger u_\uparrow
 =\frac{2E}{c}.
\end{equation*}
空间 $z$ 分量用 $\alpha^3\equiv\gamma^0\gamma^3$ 计算：
\begin{align*}
 \bar u_\uparrow\gamma^3u_\uparrow
 &=u_\uparrow^\dagger\alpha^3u_\uparrow\\
 &=2A^2\kappa\\
 &=2p.
\end{align*}
$x,y$ 分量为零，故
\begin{equation}
 \boxed{\bar u_s(p)\gamma^\mu u_s(p)=2p^\mu}
 \qquad(\mathbf p=p\mathbf e_z),
 \label{eq:explicit-dirac-current-z-dp62}
\end{equation}
且结果与选取 $\uparrow$ 或 $\downarrow$ 无关。这个算例把三个彼此相关的抽象公式连在一起：旋量归一化 $\bar uu=2m_{\mathrm e} c$、概率/粒子数密度 $u^\dagger u=2E/c$，以及协变四流与四动量同向。散射迹中的 $\bar u\gamma^\mu u$ 使用已经建立的一般协变恒等式；这里的显式分量计算提供对旋量归一化、号规和四流方向的独立核验。
\end{exbox}

自由电子的群速度和波包二阶色散由同一个 $E(\mathbf p)$ 的一阶、二阶导数决定。定义 Hessian（二阶导数）矩阵
\begin{equation}
 \boxed{
 H_{ij}(\mathbf p)
 \equiv\frac{\partial^2E}{\partial p_i\partial p_j}
 =\frac{c^2}{E}
 \left(\delta_{ij}-\frac{c^2p_ip_j}{E^2}\right),}
 \label{eq:exact_dispersion_hessian}
\end{equation}
其 SI 单位为 $\mathrm{kg^{-1}}$。若中心动量沿 $z$ 轴，先由精确色散定义中心能量和群速度，
\begin{equation*}
 \mathbf p_0=p_0\mathbf e_z,
 \qquad
 E_0=\sqrt{m_{\mathrm e}^2c^4+p_0^2c^2},
 \qquad
 v_0=\frac{c^2p_0}{E_0}.
\end{equation*}
再定义中心速度比与其 Lorentz 因子
\begin{equation*}
 \beta_0\equiv\frac{v_0}{c},
 \qquad
 \gamma_0\equiv\gamma_{\rm L}(\beta_0)
 =\frac1{\sqrt{1-\beta_0^2}}
 =\frac{E_0}{m_{\mathrm e} c^2},
 \qquad
 p_0=\gamma_0m_{\mathrm e} v_0.
\end{equation*}
则 Hessian 矩阵的两个主曲率对应两个曲率质量
\begin{equation}
 \boxed{
 m_\perp\equiv\gamma_0m_{\mathrm e},
 \qquad
 m_\parallel\equiv\gamma_0^3m_{\mathrm e},}
 \qquad
 H_{xx}=H_{yy}=\frac1{m_\perp},
 \quad
 H_{zz}=\frac1{m_\parallel}.
 \label{eq:dispersion-hessian-principal-r15}
\end{equation}
$m_\perp,m_\parallel$ 定义为色散在横向与纵向方向的曲率质量；静质量仍记为 $m_{\mathrm e}$。其中纵向关系也可写成
\begin{equation}
 \left.\frac{\dd^2 E}{\dd p^2}\right|_{p_0}
 =\frac1{\gamma_0^3m_{\mathrm e}}
 =\frac1{m_\parallel}.
 \label{eq:longitudinal-mass}
\end{equation}

任意自由纵向波包仍由精确 Fourier 积分传播
\begin{equation}
 \psi(z,t)=\frac{1}{\sqrt{2\pi\hbar}}
 \int_{-\infty}^{\infty}\dd p\,
 \widetilde\psi(p,0)
 \exp\!\left[\frac{\ii}{\hbar}\bigl(pz-E(p)t\bigr)\right].
 \label{eq:exact-wavepacket}
\end{equation}
只有当动量支持集中在 $p_0$ 附近时，才令 $\delta p=p-p_0$ 并展开
\begin{equation}
 E(p)=E_0+v_0\delta p+\frac{(\delta p)^2}{2m_\parallel}
 +\frac{E^{(3)}(p_0)}{3!}(\delta p)^3+\cdots.
 \label{eq:dispersion-taylor}
\end{equation}
若波包典型宽度为 $\sigma_p$，二阶近似可由无量纲比
\begin{equation}
 \boxed{
 \epsilon_{\rm curv}\equiv
 \frac{|E^{(3)}(p_0)|\sigma_p^3/6}
 {|E''(p_0)|\sigma_p^2/2}
 =\beta_0^2\frac{\sigma_p}{p_0}\ll1},
 \label{eq:third-order-control}
\end{equation}
直接控制；这比只写 $\sigma_p\ll p_0$ 更接近真正被忽略的三阶色散误差。

\begin{figure}[htbp]
 \centering
 \includegraphics[width=.78\linewidth]{dispersion_control.png}
 \caption{动能 $K_0=200\,\mathrm{keV}$ 时的精确相对论色散与在 $p_0$ 处的一、二阶展开。横轴为 $\delta p/p_0$，纵轴为 $(E-E_0)/(m_{\mathrm e} c^2)$；二阶曲线偏离精确曲线的位置直接显示高阶色散开始重要的区域。}
 \label{fig:dispersion-control}
\end{figure}

传播之外，\eqnrefs{eq:rel-dispersion} 还固定自由电子交换能量--动量时允许的运动学区域。为把吸收与发射放在同一个母式中，先把环境中一个\emph{正频模式}的四波矢定义为
\begin{equation}
 K^\mu=\left(\frac{\omega}{c},\mathbf k\right),
 \qquad \omega>0,
\end{equation}
并用符号 $\sigma=+1$ 表示电子吸收该模式量子、$\sigma=-1$ 表示电子发射该模式量子。电子四动量的有限交换统一写成
\begin{equation}
 \boxed{P_f^\mu=P_i^\mu+\sigma\hbar K^\mu.}
 \label{eq:four-transfer}
\end{equation}
这里 $K^\mu$ 固定描述正频模式本身，$\sigma$ 单独记录能量--动量流向，因此吸收与发射共享同一组 $\omega>0$、$\mathbf k$ 模式参数。初末电子都在同一质量壳 $P_i^2=P_f^2=m_{\mathrm e}^2c^2$，因此
\begin{equation}
 \boxed{2\sigma P_i\cdot K+\hbar K^2=0.}
 \label{eq:exact-recoil-four}
\end{equation}
这是一切单量子吸收/发射运动学的母恒等式。真空实光子满足 $K^2=0$；对孤立自由电子，未来指向的类时 $P_i$ 与非零未来指向类光 $K$ 有 $P_i\!\cdot K>0$，故无论 $\sigma=+1$ 还是 $-1$ 都不能满足上式。近场、周期结构、介质或散射体能够提供额外空间动量，或者使参与交换的模式具有 $K^2\neq0$，从而解除这一运动学禁阻。

对一维纵向几何，模式波数记为 $\kappa_z$。同一母条件等价于
\begin{equation}
 \boxed{E(p_0+\sigma\hbar \kappa_z)-E(p_0)=\sigma\hbar\omega.}
 \label{eq:exact-longitudinal-match}
\end{equation}
只有在 $|\hbar\kappa_z|$ 落入正文已经建立的局部 Taylor 控制窗以后，才把精确色散展开。为使截断误差本身也是一个定义明确的量，记二阶 Taylor 余项为
\begin{equation*}
 \mathcal R_3^{(\sigma)}(\kappa_z)
 \equiv
 \frac{(\sigma\hbar\kappa_z)^3}{2}
 \int_0^1\dd s\,(1-s)^2
 E^{(3)}\!\left(p_0+s\sigma\hbar\kappa_z\right).
\end{equation*}
于是精确关系可写成
\begin{equation}
 \boxed{
 \omega
 =v_0\kappa_z+\sigma\frac{\hbar \kappa_z^2}{2m_\parallel}
 +\frac{\mathcal R_3^{(\sigma)}(\kappa_z)}{\sigma\hbar}.}
 \label{eq:recoil-phase-match}
\end{equation}
余项满足
\begin{equation*}
 \left|
 \frac{\mathcal R_3^{(\sigma)}(\kappa_z)}{\hbar}
 \right|
 \le
 \frac{\hbar^2|\kappa_z|^3}{6}
 \sup_{|p-p_0|\le\hbar|\kappa_z|}|E^{(3)}(p)|.
\end{equation*}
二阶反冲项的符号随吸收/发射改变，而最低阶同步项 $v_0\kappa_z$ 与流向无关。定义曲率项相对线性项的控制量
\begin{equation*}
 \chi_r^{(1)}\equiv
 \left|\frac{\hbar \kappa_z}{2m_\parallel v_0}\right|,
\end{equation*}
只有在 $\chi_r^{(1)}\ll1$ 时才可进一步写成常用同步条件
\begin{equation}
 \omega\simeq v_0\kappa_z.
 \label{eq:linear-sync}
\end{equation}
若连续吸收同一正频模式量子，取 $\sigma=+1$；此时已占据能量梯上相邻跃迁的精确失谐必须逐级计算：
\begin{equation}
 \boxed{
 \Delta_\ell\equiv
 \frac{E(p_0+(\ell+1)\hbar \kappa_z)-E(p_0+\ell\hbar \kappa_z)}{\hbar}-\omega.}
 \label{eq:exact-ladder-detuning}
\end{equation}
从宽能量梯约化到少能级反冲模型时，真正需要比较的是被占据区间内的 $\Delta_\ell$ 变化，而不是只检查单次交换的一个失谐。

自由 Dirac 方程同时固定三种共同结构：在壳旋量与完备关系固定 QED 外态；精确相对论色散及其 Hessian 固定自由传播；同一质量壳条件固定反冲与相位匹配。低反冲、近轴和同步近似都由这些精确母式在显式小参数下约化。


\begin{derivation}[从\eqnrefs{eq:u-dirac}到\eqnrefs{eq:rel-dispersion}]
对正频旋量左乘 $\slashed p+m_{\mathrm e} c$：

\begin{align*}
 (\slashed p+m_{\mathrm e} c)(\slashed p-m_{\mathrm e} c)u_s
 &=\left(\gamma^\mu\gamma^\nu p_\mu p_\nu-m_{\mathrm e}^2c^2\right)u_s.
\end{align*}
把矩阵乘积分成对称与反对称两部分，
\begin{align*}
 \gamma^\mu\gamma^\nu p_\mu p_\nu
 &=\frac12\{\gamma^\mu,\gamma^\nu\}p_\mu p_\nu
 +\frac12[\gamma^\mu,\gamma^\nu]p_\mu p_\nu\\
 &=\eta^{\mu\nu}p_\mu p_\nu\,\id(4)
 =p^2\id(4),
\end{align*}
因为 $p_\mu p_\nu$ 对称，而 $[\gamma^\mu,\gamma^\nu]$ 反对称。于是
\begin{align*}
 (p^2-m_{\mathrm e}^2c^2)u_s&=0.
\end{align*}
非零旋量要求
\begin{align*}
 p^2&=m_{\mathrm e}^2c^2,\\
 \frac{E_{\mathbf p}^2}{c^2}-\mathbf p^2&=m_{\mathrm e}^2c^2,\\
 E_{\mathbf p}&=\sqrt{m_{\mathrm e}^2c^4+c^2\mathbf p^2},
\end{align*}
其中最后选择正能分支。对正文\eqnrefs{eq:v-dirac} 左乘 $\slashed p-m_{\mathrm e} c$ 得到同一二阶条件，因此电子与正电子外线共享同一正能质量壳参数化。
\end{derivation}

\begin{derivation}[从\eqnrefs{eq:u-dirac}到\eqnrefs{eq:explicit-v}]
在 Dirac 表示中把四分量旋量写成两个二分量块。对正频支令
\[
u_s(p)=\begin{pmatrix}\phi_s\\ \chi_s\end{pmatrix}.
\]
正文\eqnrefs{eq:u-dirac} 给出
\begin{align*}
(E-m_{\mathrm e} c^2)\phi_s
-c\boldsymbol\sigma_{\rm P}\!\cdot\!\mathbf p\,\chi_s&=0,\\
c\boldsymbol\sigma_{\rm P}\!\cdot\!\mathbf p\,\phi_s
-(E+m_{\mathrm e} c^2)\chi_s&=0.
\end{align*}
第二式先确定下分量：
\[
\chi_s
=\frac{c\boldsymbol\sigma_{\rm P}\cdot\mathbf p}{E+m_{\mathrm e} c^2}\phi_s
=\frac{\boldsymbol\sigma_{\rm P}\cdot\mathbf p}{E/c+m_{\mathrm e} c}\phi_s.
\]
取
$\phi_s=\sqrt{E/c+m_{\mathrm e} c}\,\xi_s^{(+)}$，便得到正文\eqnrefs{eq:explicit-u}。用质量壳 $E^2=m_{\mathrm e}^2c^4+c^2\mathbf p^2$ 检查归一化：
\begin{align*}
\bar u_su_{s'}
&=\left(\frac Ec+m_{\mathrm e} c\right)
\left[
1-\frac{\mathbf p^2}{(E/c+m_{\mathrm e} c)^2}
\right]\delta_{ss'}\\
&=2m_{\mathrm e} c\,\delta_{ss'}.
\end{align*}

对负频支令
\[
v_s(p)=\begin{pmatrix}\phi'_s\\ \chi'_s\end{pmatrix}.
\]
正文\eqnrefs{eq:v-dirac} 的两个分量方程为
\begin{align*}
(E+m_{\mathrm e} c^2)\phi'_s
-c\boldsymbol\sigma_{\rm P}\!\cdot\!\mathbf p\,\chi'_s&=0,\\
c\boldsymbol\sigma_{\rm P}\!\cdot\!\mathbf p\,\phi'_s
-(E-m_{\mathrm e} c^2)\chi'_s&=0.
\end{align*}
第一式给
\[
\phi'_s
=\frac{\boldsymbol\sigma_{\rm P}\cdot\mathbf p}{E/c+m_{\mathrm e} c}\chi'_s.
\]
取
$\chi'_s=\sqrt{E/c+m_{\mathrm e} c}\,\xi_s^{(-)}$，即得到正文\eqnrefs{eq:explicit-v}。同样利用质量壳，
\begin{align*}
\bar v_sv_{s'}
&=\left(\frac Ec+m_{\mathrm e} c\right)
\left[
\frac{\mathbf p^2}{(E/c+m_{\mathrm e} c)^2}-1
\right]\delta_{ss'}\\
&=-2m_{\mathrm e} c\,\delta_{ss'}.
\end{align*}
\end{derivation}





\begin{derivation}[从\eqnrefs{eq:u-dirac}到\eqnrefs{eq:v-complete}]

质量壳给 $\slashed p^{\,2}=m_{\mathrm e}^2c^2$，因此
\begin{align*}
\Lambda_+^2
&=\frac{(\slashed p+m_{\mathrm e} c)^2}{4m_{\mathrm e}^2c^2}
=\frac{\slashed p+m_{\mathrm e} c}{2m_{\mathrm e} c}=\Lambda_+,\\
\Lambda_-^2
&=\frac{(m_{\mathrm e} c-\slashed p)^2}{4m_{\mathrm e}^2c^2}
=\Lambda_-.
\end{align*}
并且
\begin{equation*}
(\slashed p-m_{\mathrm e} c)\Lambda_+=0,
\qquad
(\slashed p+m_{\mathrm e} c)\Lambda_-=0.
\end{equation*}
由于 $\operatorname{tr}\slashed p=0$，两投影的迹均为 $2$，所以其像空间分别与正能、负频 Dirac 解空间完全相同。由 $\bar u_su_{s'}=2m_{\mathrm e} c\,\delta_{ss'}$，正能解空间上的恒等投影为
\begin{equation*}
\Lambda_+=\frac1{2m_{\mathrm e} c}\sum_su_s\bar u_s,
\end{equation*}
故得到\eqnrefs{eq:u-complete}。负频基满足 $\bar v_sv_{s'}=-2m_{\mathrm e} c\,\delta_{ss'}$，于是
\begin{equation*}
\Lambda_-=-\frac1{2m_{\mathrm e} c}\sum_sv_s\bar v_s,
\end{equation*}
等价于\eqnrefs{eq:v-complete}。
\end{derivation}


\begin{derivation}[从\eqnrefs{eq:dispersion-taylor}到\eqnrefs{eq:third-order-control}]
为得到以动量宽度 $\sigma_p$ 表示的局部截断判据，取工作窗边界 $|\delta p|=\sigma_p$；正文\eqnrefs{eq:dispersion-taylor} 的最后保留项和第一个舍弃项分别为
\[
\Delta E_2=\frac12E''(p_0)\sigma_p^2,
\qquad
\Delta E_3=\frac16E^{(3)}(p_0)\sigma_p^3.
\]
先从精确色散 $E(p)=\sqrt{m_{\mathrm e}^2c^4+p^2c^2}$ 计算所需导数：
\begin{align*}
E'(p)&=\frac{pc^2}{E(p)},\\
E''(p)&=\frac{m_{\mathrm e}^2c^6}{E(p)^3},\\
E^{(3)}(p)
&=-\frac{3m_{\mathrm e}^2c^8p}{E(p)^5}.
\end{align*}
在 $p=p_0$ 处用 $E_0=\gamma_0m_{\mathrm e} c^2$ 与 $p_0c/E_0=\beta_0$，两项之比为
\begin{align*}
\epsilon_{\rm curv}
&=\left|\frac{\Delta E_3}{\Delta E_2}\right|\\
&=\frac{|E^{(3)}(p_0)|\sigma_p}{3|E''(p_0)|}\\
&=\frac{p_0c^2}{E_0^2}\sigma_p\\
&=\frac{p_0^2c^2}{E_0^2}\frac{\sigma_p}{p_0}\\
&=\beta_0^2\frac{\sigma_p}{p_0}.
\end{align*}
这就是正文\eqnrefs{eq:third-order-control}。
\end{derivation}

\begin{derivation}[从\eqnrefs{eq:four-transfer}到\eqnrefs{eq:recoil-phase-match}]
正文\eqnrefs{eq:four-transfer} 把正频模式与能量流向分开写成
\[
P_f=P_i+\sigma\hbar K,
\qquad \sigma^2=1.
\]
初态和末态都是同一静质量的真实电子，所以
\[
P_i^2=P_f^2=m_{\mathrm e}^2c^2.
\]
两式相减时静质量项完全消去，而 $\sigma^2=1$ 保留二阶反冲项：
\begin{align*}
0
&=P_f^2-P_i^2\\
&=(P_i+\sigma\hbar K)^2-P_i^2\\
&=2\sigma\hbar P_i\!\cdot K+\hbar^2K^2.
\end{align*}
除以 $\hbar$ 得到正文\eqnrefs{eq:exact-recoil-four}。因此吸收与发射的差别只出现在与 $P_i\!\cdot K$ 相乘的一阶流向符号中；模式本身仍始终取 $\omega>0$。

在一维纵向通道中，$P_f=P_i+\sigma\hbar K$ 给出
\begin{align*}
p_f&=p_0+\sigma\hbar\kappa_z,\\
E_f-E_i&=\sigma\hbar\omega,
\end{align*}
故正文\eqnrefs{eq:exact-longitudinal-match} 是四维母恒等式的能量表示。若 $|\hbar\kappa_z|$ 落在正文已经给出的 Taylor 控制窗内，对
$\delta p=\sigma\hbar\kappa_z$ 使用带积分余项的 Taylor 定理：
\begin{align*}
E(p_0+\delta p)-E(p_0)
&=E'(p_0)\delta p
+\frac{E''(p_0)}{2}\delta p^2
+\frac{\delta p^3}{2}
\int_0^1\dd s\,(1-s)^2E^{(3)}(p_0+s\delta p)\\
&=\sigma v_0\hbar\kappa_z
+\frac{\hbar^2\kappa_z^2}{2m_\parallel}
+\mathcal R_3^{(\sigma)}(\kappa_z).
\end{align*}
与 $\sigma\hbar\omega$ 比较，并利用 $1/\sigma=\sigma$，得到正文\eqnrefs{eq:recoil-phase-match}。积分余项还立刻给出
\begin{align*}
\left|\mathcal R_3^{(\sigma)}(\kappa_z)\right|
&\le
\frac{|\hbar\kappa_z|^3}{2}
\sup_{|p-p_0|\le\hbar|\kappa_z|}|E^{(3)}(p)|
\int_0^1(1-s)^2\dd s\\
&=\frac{\hbar^3|\kappa_z|^3}{6}
\sup_{|p-p_0|\le\hbar|\kappa_z|}|E^{(3)}(p)|.
\end{align*}
这里可以直接看出：线性同步项不区分吸收和发射，第一项真正区分两种流向的反冲修正出现在偶次色散项；舍弃高阶项时也有明确的误差上界，而不是只靠一个形式上的省略号。
\end{derivation}

\section{Maxwell Green 张量与散射}

Maxwell--Dirac 作用量给出的是局域微分方程，但真实实验通常指定的是电流源、材料和边界条件。与其为每一种几何重新解偏微分方程，更有用的做法是先问一个一般问题：在给定线性环境中，位于 $\mathbf r'$ 的单位源会在 $\mathbf r$ 产生怎样的场？这个“源到场”的线性响应核就是推迟 Green 张量。

把频域 Maxwell 方程写成线性算符以后，Green 张量就是该算符带因果边界条件的逆。由这个逆核可以统一建立三类对象：任意给定源的场响应；材料扰动引起的散射；以及由算符对称性和能量守恒导出的互易与光学恒等式。这样，几何、材料色散和边界条件都进入同一个响应核，而不必在主理论中提前列举具体实验名称。

这些结构首先属于经典线性响应。自由电子能量交换、辐射和量子化介质都使用同一个 Green 张量接口，只是源、材料本构和量子态不同；环境量子化属于另一个独立层级。

Maxwell--Dirac 场方程已经固定局域源项与真空中的因果传播。介质、边界和开放几何并不改变 Maxwell 方程的局域来源，但会改变“给定电流以后场如何传播”的逆问题。在线性时不变环境中，这个逆问题由频域 Maxwell 算符及其推迟 Green 张量统一描述；散射、耗散和后续宏观 QED 都从同一个逆核出发。

先考虑空间局域但允许各向异性和色散的线性介质。相对介电率与相对磁导率张量分别记为 $\boldsymbol\varepsilon_r(\mathbf r,\omega)$、$\boldsymbol\mu_r(\mathbf r,\omega)$，二者均无量纲；$\omega$ 是角频率，单位为 $\mathrm{s^{-1}}$。采用全文时间 Fourier 约定 $\ee^{-\ii\omega t}$。消去磁场以后，电场满足的二阶微分算符定义为
\begin{equation}
 \boxed{
 \mathsf L_{\rm EM}(\omega)
 \equiv
 \nabla\times\boldsymbol\mu_r^{-1}(\mathbf r,\omega)\nabla\times
 -\frac{\omega^2}{c^2}\boldsymbol\varepsilon_r(\mathbf r,\omega).}
 \label{eq:Maxwell-operator-SI-r12}
\end{equation}
$\mathsf L_{\rm EM}$ 的量纲为长度的负二次方。材料本构、几何和边界条件都进入这个算符；若材料响应具有空间非局域性，$\boldsymbol\varepsilon_r$ 与 $\boldsymbol\mu_r$ 应相应提升为积分核，而同一算符结构仍然成立。

给定辐射出射条件，推迟并矢 Green 张量 $\mathbf G^R(\mathbf r,\mathbf r';\omega)$ 定义为 $\mathsf L_{\rm EM}$ 的逆核：
\begin{equation}
 \boxed{
 \mathsf L_{\rm EM}(\omega)
 \mathbf G^R(\mathbf r,\mathbf r';\omega)
 =\mathbf I\,\delta^{(3)}(\mathbf r-\mathbf r').}
 \label{eq:GR-Maxwell-SI-r12}
\end{equation}
其中 $\mathbf I$ 是三维单位并矢，$[\mathbf G^R]=\mathrm{m^{-1}}$。上标 $R$ 表示推迟边界条件：源在给定频率产生的场必须满足向外辐射而不是从无穷远流入的条件。因而 Green 张量不仅包含材料参数，还包含开放系统的物理边界条件。

若时域外加电流密度为 $\mathbf J(\mathbf r,t)$，则按全文 Fourier convention 定义的频谱电流 $\widetilde{\mathbf J}(\mathbf r,\omega)$ 的单位为 $\mathrm{A\,s\,m^{-2}}$。频域 Maxwell 方程可压缩成
\begin{equation}
\boxed{
\mathsf L_{\rm EM}(\omega)\widetilde{\mathbf E}(\mathbf r,\omega)
=\ii\mu_0\omega\widetilde{\mathbf J}(\mathbf r,\omega).}
\label{eq:maxwell-source-operator-dp68}
\end{equation}
两条旋度方程合并后得到与具体几何无关的线性源方程。利用\eqnrefs{eq:GR-Maxwell-SI-r12} 的推迟逆核，得到电流到电场的线性响应母式
\begin{equation}
 \boxed{
 \widetilde{\mathbf E}(\mathbf r,\omega)
 =\ii\mu_0\omega
 \int\dd^3 r'\,
 \mathbf G^R(\mathbf r,\mathbf r';\omega)
 \cdot\widetilde{\mathbf J}(\mathbf r',\omega).}
 \label{eq:E-from-GJ-SI-r12}
\end{equation}
量纲检查给 $[\mu_0\omega G^R\widetilde J\,\dd^3 r']=\mathrm{V\,s\,m^{-1}}$，与时间 Fourier 电场 $\widetilde{\mathbf E}$ 的单位一致。这个公式是后续 EELS、辐射和环境诱导自能的经典响应起点；具体电子轨迹只会改变源 $\widetilde{\mathbf J}$，不会改变环境核 $\mathbf G^R$。

入射场与结构散射可由参考算符分解。选取参考算符 $\mathsf L_0$ 及其推迟 Green 张量 $\mathbf G_0^R=\mathsf L_0^{-1}$，并定义散射势算符
\begin{equation*}
 \mathsf V_{\rm EM}\equiv\mathsf L_0-\mathsf L_{\rm EM}.
\end{equation*}
同一源在真实环境中的场满足
\begin{equation}
\boxed{
\widetilde{\mathbf E}
=\widetilde{\mathbf E}_{\rm inc}+\mathbf G_0^R\mathsf V_{\rm EM}\widetilde{\mathbf E}.}
\label{eq:maxwell-lippmann-schwinger-dp68}
\end{equation}
右边第二项表示场先由结构散射一次，再由参考环境传播；把其中的总场再次用同一方程代回，就自动产生二次、三次直到任意次数的散射。只要相应预解式存在，这个级数可重求和为电磁 $T$ 算符
\begin{equation}
 \boxed{
 \mathsf T_{\rm EM}
 =\mathsf V_{\rm EM}
 (\mathbf I-\mathbf G_0^R\mathsf V_{\rm EM})^{-1}.}
 \label{eq:maxwell-t-operator-def-r16}
\end{equation}
于是频域总场 $\widetilde{\mathbf E}=\widetilde{\mathbf E}_{\rm inc}+\widetilde{\mathbf E}_{\rm sc}$ 的散射部分为
\begin{equation}
 \boxed{
 \widetilde{\mathbf E}_{\rm sc}
 =\mathbf G_0^R\mathsf T_{\rm EM}\widetilde{\mathbf E}_{\rm inc}.}
 \label{eq:maxwell_scattered_field_t_operator}
\end{equation}
这里 $\widetilde{\mathbf E}_{\rm inc}$ 是同一外部源在参考环境 $\mathsf L_0$ 中产生的 Fourier 电场。球、圆柱、平面界面和数值离散结构的差别，只在于如何表示并求解同一个 $\mathsf T_{\rm EM}$。

若局域介质满足互易本构条件
$\boldsymbol\varepsilon_r^T=\boldsymbol\varepsilon_r$ 与
$\boldsymbol\mu_r^T=\boldsymbol\mu_r$，并对两组场采用同一辐射边界条件，则分部积分后 Maxwell 算符满足
\begin{equation}
\boxed{
\int\dd^3 r\,\mathbf u\!\cdot\!\mathsf L_{\rm EM}\mathbf v
=\int\dd^3 r\,\mathbf v\!\cdot\!\mathsf L_{\rm EM}\mathbf u.}
\label{eq:maxwell-bilinear-reciprocity-dp68}
\end{equation}
这条式子只是说：在互易介质中，交换“源场”和“测试场”不会改变这一双线性响应。把两组测试场选成点源产生的 Green 场，就得到 Lorentz 互易关系
\begin{equation}
 \boxed{
 \mathbf G^R(\mathbf r,\mathbf r';\omega)
 =(\mathbf G^R)^T(\mathbf r',\mathbf r;\omega).}
 \label{eq:maxwell-green-reciprocity-r16}
\end{equation}
互易性约束交换源点与观测点时的张量转置；它与“无损”是不同条件，因此有吸收的互易介质仍可满足此式。

Green 张量还把材料耗散与场的谱密度联系起来。一般算符恒等式为
\begin{equation}
 \boxed{
 \mathbf G^R-(\mathbf G^R)^\dagger
 =(\mathbf G^R)^\dagger
 (\mathsf L_{\rm EM}^\dagger-\mathsf L_{\rm EM})
 \mathbf G^R.}
 \label{eq:resolvent-imag-identity-r12}
\end{equation}
左边是因果预解式的反厄米部分，右边则显示它只能来自 Maxwell 算符自身的反厄米部分，也就是吸收通道。若 $\boldsymbol\mu_r$ 为实厄米张量且耗散只来自电响应，则上式化为
\begin{equation}
 \boxed{
 \frac{\omega^2}{c^2}
 \int\dd^3 s\,
 \mathbf G^R(\mathbf r,\mathbf s;\omega)
 \cdot\operatorname{Im}\boldsymbol\varepsilon_r(\mathbf s,\omega)
 \cdot(\mathbf G^R)^\dagger(\mathbf r',\mathbf s;\omega)
 =\operatorname{Im}\mathbf G^R(\mathbf r,\mathbf r';\omega).}
 \label{eq:green-optical-identity-SI-r12}
\end{equation}
这条光学恒等式说明同一个 $\operatorname{Im}\mathbf G^R$ 同时记录传播与不可逆耗散。存在磁损耗时，右侧仍由\eqnrefs{eq:resolvent-imag-identity-r12} 控制，只需在耗散核中加入 $\operatorname{Im}\boldsymbol\mu_r^{-1}$ 的贡献。



\begin{derivation}[从\eqnrefs{eq:maxwell-lippmann-schwinger-dp68}到\eqnrefs{eq:maxwell_scattered_field_t_operator}]
把背景 Maxwell Green 算符记为 $\mathbf G_0^R$，介质相对于背景的散射势记为 $\mathsf V_{\rm EM}$。正文的 Lippmann--Schwinger 方程是
\[
\widetilde{\mathbf E}
=\widetilde{\mathbf E}_{\rm inc}
+\mathbf G_0^R\mathsf V_{\rm EM}\widetilde{\mathbf E}.
\]
先不要展开级数，而把未知总场的项移到左边：
\[
(\mathbf I-\mathbf G_0^R\mathsf V_{\rm EM})\widetilde{\mathbf E}
=\widetilde{\mathbf E}_{\rm inc}.
\]
只要所考察复频率不位于该算符的谱上，就存在逆算符，故
\[
\widetilde{\mathbf E}
=(\mathbf I-\mathbf G_0^R\mathsf V_{\rm EM})^{-1}
\widetilde{\mathbf E}_{\rm inc}.
\]
散射场因此为
\begin{align*}
\widetilde{\mathbf E}_{\rm sc}
&\equiv\widetilde{\mathbf E}-\widetilde{\mathbf E}_{\rm inc}\\
&=(\mathbf I-\mathbf G_0^R\mathsf V_{\rm EM})^{-1}
\mathbf G_0^R\mathsf V_{\rm EM}\widetilde{\mathbf E}_{\rm inc}.
\end{align*}
现在用可直接验证的算符恒等式
\[
(\mathbf I-AB)^{-1}A=A(\mathbf I-BA)^{-1}.
\]
令 $A=\mathbf G_0^R$、$B=\mathsf V_{\rm EM}$，便有
\[
\widetilde{\mathbf E}_{\rm sc}
=\mathbf G_0^R
\mathsf V_{\rm EM}
(\mathbf I-\mathbf G_0^R\mathsf V_{\rm EM})^{-1}
\widetilde{\mathbf E}_{\rm inc}.
\]
因此定义
\[
\boxed{
\mathsf T_{\rm EM}
\equiv\mathsf V_{\rm EM}
(\mathbf I-\mathbf G_0^R\mathsf V_{\rm EM})^{-1}}
\]
后立即得到
\[
\boxed{
\widetilde{\mathbf E}_{\rm sc}
=\mathbf G_0^R\mathsf T_{\rm EM}\widetilde{\mathbf E}_{\rm inc}},
\]
即正文\eqnrefs{eq:maxwell_scattered_field_t_operator}。

多次散射级数只是这个精确逆算符在收敛区域内的 Neumann 展开。若某个算符范数满足 $\lVert\mathbf G_0^R\mathsf V_{\rm EM}\rVert<1$，则
\begin{align*}
(\mathbf I-\mathbf G_0^R\mathsf V_{\rm EM})^{-1}
&=\sum_{n=0}^{\infty}(\mathbf G_0^R\mathsf V_{\rm EM})^n,\\
\mathsf T_{\rm EM}
&=\mathsf V_{\rm EM}
+\mathsf V_{\rm EM}\mathbf G_0^R\mathsf V_{\rm EM}
+\mathsf V_{\rm EM}\mathbf G_0^R\mathsf V_{\rm EM}\mathbf G_0^R\mathsf V_{\rm EM}+\cdots.
\end{align*}
第 $n$ 个额外因子 $\mathbf G_0^R\mathsf V_{\rm EM}$ 对应“再传播一次、再散射一次”。即使 Neumann 级数本身不收敛，$\mathsf T_{\rm EM}$ 仍由前面的算符逆定义；其极点对应齐次方程
\[
(\mathbf I-\mathbf G_0^R\mathsf V_{\rm EM})\mathbf E=0
\]
出现非零解，也就是
$\ker(\mathbf I-\mathbf G_0^R\mathsf V_{\rm EM})\neq\{0\}$；在有限维离散化后，这才等价于相应行列式为零。它对应无外源条件下仍可存在的本征模或共振。这样 $T$ 算符同时统一了弱散射 Born 级数与强多重散射的谱结构。
\end{derivation}

\begin{derivation}[从\eqnrefs{eq:maxwell-bilinear-reciprocity-dp68}到\eqnrefs{eq:maxwell-green-reciprocity-r16}]

正文\eqnrefs{eq:maxwell-bilinear-reciprocity-dp68} 已经把介质互易性写成 Maxwell 算符的双线性对称关系。要把它变成 Green 张量的源点交换关系，只需让两组测试场分别由两个点源产生。

取任意常矢量 \(\mathbf a,\mathbf b\)，令
\begin{equation*}
\mathbf u(\mathbf s)=\mathbf G^R(\mathbf s,\mathbf r;\omega)\cdot\mathbf a,
\qquad
\mathbf v(\mathbf s)=\mathbf G^R(\mathbf s,\mathbf r';\omega)\cdot\mathbf b.
\end{equation*}
由 Green 方程，
\begin{equation*}
\mathsf L_{\rm EM}\mathbf u
=\mathbf a\,\delta^{(3)}(\mathbf s-\mathbf r),
\qquad
\mathsf L_{\rm EM}\mathbf v
=\mathbf b\,\delta^{(3)}(\mathbf s-\mathbf r').
\end{equation*}
代入正文\eqnrefs{eq:maxwell-bilinear-reciprocity-dp68} 后，两个空间积分分别被 \(\delta\) 函数取样：
\begin{equation*}
\mathbf a\cdot\mathbf G^R(\mathbf r,\mathbf r';\omega)\cdot\mathbf b
=\mathbf b\cdot\mathbf G^R(\mathbf r',\mathbf r;\omega)\cdot\mathbf a.
\end{equation*}
因为 \(\mathbf a,\mathbf b\) 任意，所以
\begin{equation*}
\mathbf G^R(\mathbf r,\mathbf r';\omega)
=[\mathbf G^R(\mathbf r',\mathbf r;\omega)]^T,
\end{equation*}
即正文\eqnrefs{eq:maxwell-green-reciprocity-r16}。
\end{derivation}

\begin{derivation}[从\eqnrefs{eq:GR-Maxwell-SI-r12}到\eqnrefs{eq:green-optical-identity-SI-r12}]
令固定实频 $\omega>0$ 下的 Maxwell 算符为 $\mathsf L_{\rm EM}$，并定义
\[
\mathbf G^R=\mathsf L_{\rm EM}^{-1},
\qquad
\mathbf G^{R\dagger}=(\mathsf L_{\rm EM}^{\dagger})^{-1}.
\]
对任意两个可逆算符 $A,B$，有两种等价的预解式恒等式
\[
A^{-1}-B^{-1}
=A^{-1}(B-A)B^{-1}
=B^{-1}(B-A)A^{-1}.
\]
第一式例如由
$A^{-1}BB^{-1}-A^{-1}AB^{-1}=A^{-1}-B^{-1}$ 逐项得到。取
$A=\mathsf L_{\rm EM}$、$B=\mathsf L_{\rm EM}^{\dagger}$，先用第一种次序得到
\begin{align*}
\mathbf G^R-\mathbf G^{R\dagger}
&=\mathbf G^R
(\mathsf L_{\rm EM}^{\dagger}-\mathsf L_{\rm EM})
\mathbf G^{R\dagger}.
\end{align*}
材料的吸收与色散信息集中在算符差 $\mathsf L_{\rm EM}^{\dagger}-\mathsf L_{\rm EM}$ 中。对局域电介质且磁响应无损的情形，
\[
\mathsf L_{\rm EM}
=\nabla\times\nabla\times
-\frac{\omega^2}{c^2}\boldsymbol\varepsilon_r,
\]
故
\[
\mathsf L_{\rm EM}^{\dagger}-\mathsf L_{\rm EM}
=\frac{\omega^2}{c^2}
(\boldsymbol\varepsilon_r-\boldsymbol\varepsilon_r^{\dagger})
=2\ii\frac{\omega^2}{c^2}
\operatorname{Im}\boldsymbol\varepsilon_r.
\]
代回并除以 $2\ii$：
\[
\operatorname{Im}\mathbf G^R
=\frac{\omega^2}{c^2}
\mathbf G^R
(\operatorname{Im}\boldsymbol\varepsilon_r)
\mathbf G^{R\dagger}.
\]
把算符乘积写成坐标核。若介电张量局域，
$\operatorname{Im}\boldsymbol\varepsilon_r$ 的核为
$\operatorname{Im}\boldsymbol\varepsilon_r(\mathbf s,\omega)
\delta^{(3)}(\mathbf s-\mathbf s')$，于是
\[
\boxed{
\operatorname{Im}\mathbf G^R(\mathbf r,\mathbf r';\omega)
=\frac{\omega^2}{c^2}\int\dd^3s\,
\mathbf G^R(\mathbf r,\mathbf s;\omega)
\cdot\operatorname{Im}\boldsymbol\varepsilon_r(\mathbf s,\omega)
\cdot\mathbf G^{R\dagger}(\mathbf r',\mathbf s;\omega)}.
\]
这就是正文\eqnrefs{eq:green-optical-identity-SI-r12}。

该恒等式还解释了为什么 $\operatorname{Im}\mathbf G^R$ 能作为损失谱权重。为把正定性写成最直观的场能量形式，改用上面第二种等价次序，亦有
\[
\operatorname{Im}\mathbf G^R
=\frac{\omega^2}{c^2}
\mathbf G^{R\dagger}
(\operatorname{Im}\boldsymbol\varepsilon_r)
\mathbf G^R.
\]
对任意测试电流 $\mathbf J$，定义由它激发的场
\[
\mathbf E_J(\mathbf s)
\equiv\int\dd^3r'\,\mathbf G^R(\mathbf s,\mathbf r';\omega)\cdot\mathbf J(\mathbf r').
\]
于是左右与 $\mathbf J^*,\mathbf J$ 收缩得到
\[
\iint\dd^3r\dd^3r'\,
\mathbf J^*(\mathbf r)\cdot\operatorname{Im}\mathbf G^R(\mathbf r,\mathbf r')\cdot\mathbf J(\mathbf r')
=\frac{\omega^2}{c^2}\int\dd^3s\,
\mathbf E_J^*(\mathbf s)\cdot
\operatorname{Im}\boldsymbol\varepsilon_r(\mathbf s)\cdot\mathbf E_J(\mathbf s).
\]
对被动介质，$\operatorname{Im}\boldsymbol\varepsilon_r\succeq0$，所以右边非负。因而 Green 张量的反厄米部分不是一个形式上的“取虚部”规则，而是由材料局域耗散经过两次传播映射回源空间后得到的正半定谱核；这正是后文 EELS、LDOS 与涨落--耗散关系共同使用的数学结构。
\end{derivation}

